\chapter{Lie-theoretic inputs and the trust base}\label{ch:L02}
\begingroup\setlength{\emergencystretch}{3em}

\providecommand{\ltwog}{\mathfrak{g}}
\providecommand{\ltwoghat}{\widehat{\mathfrak{g}}}
\providecommand{\ltwop}{\mathfrak{p}}
\providecommand{\ltwon}{\mathfrak{n}}
\providecommand{\ltwovac}{|0\rangle}
\providecommand{\ltwoW}{\mathfrak{W}}
\providecommand{\ltwoVir}{\mathfrak{Vir}}
\providecommand{\ltwoUe}{\mathcal{U}}
\providecommand{\ltwoInd}{\operatorname{Ind}}
\providecommand{\ltwoJ}{\mathcal{J}}
\providecommand{\ltwoV}{\mathcal{V}}

This chapter documents the release of the LeanVOA library named in the front matter. Every Lean quotation below is checked by \texttt{scripts/textbook/sync.py}
to occur verbatim (up to whitespace and comments) in the source files of that release. The machine
checks reported below were run at an earlier state of the library, after \lean{pbw} became a theorem,
and were not re-run for this edition unless stated otherwise; a line \texttt{...} marks omitted structure fields, and a declaration head
ending in \lean{:= by} or \lean{:=} is printed without its proof. The chapter covers the files of
\leanpath{VOA/\allowbreak Foundations/\allowbreak } that are not about formal calculus:
\leanpath{AffineLie.lean}, \leanpath{RootDatum.lean}, \leanpath{PBWBasis.lean},
\leanpath{Virasoro.lean}, \leanpath{VirasoroModule.lean}, \leanpath{Lattices.lean},
\leanpath{LatticeCocycle.lean}, \leanpath{LatticeFinite.lean}, \leanpath{LatticeOrthoFrame.lean},
\leanpath{LatticeAlgebraTwisted.lean} and \leanpath{Parity.lean}, together with the four files of
\leanpath{VOA/\allowbreak Foundations/\allowbreak PBWTheorem/} (\leanpath{Construction.lean},
\leanpath{LieRep.lean}, \leanpath{Induced.lean}, \leanpath{VacuumBasis.lean}), which prove the
Poincar\'e--Birkhoff--Witt input of the library.

These files supply the \emph{Lie-theoretic raw material} from which the rest of the library builds
vertex operator algebras: the Witt and Virasoro algebras (the symmetry algebra of every conformal
field theory), the affine Kac--Moody algebras (the current algebras of Wess--Zumino--Witten
models), their vacuum modules, the root data of finite-dimensional simple Lie algebras, and the
lattice data (integral lattices, the Frenkel--Lepowsky--Meurman cocycle, the twisted group
algebra) from which lattice vertex algebras are made. They also contain the library's
Poincar\'e--Birkhoff--Witt input, the theorem
\begin{equation*}
 \lean{theorem pbw : PBWPrinciple}
\end{equation*}
in \leanpath{VOA/\allowbreak Foundations/\allowbreak AffineLie.lean}, proved from Mathlib alone through the
relative Poincar\'e--Birkhoff--Witt construction of \leanpath{VOA/\allowbreak Foundations/\allowbreak PBWTheorem/}.
\textbf{The library has no project axiom.} What the machine confirmed for this chapter is that
\texttt{\#print axioms VOA.pbw}, and the same command for every declaration named in a status box of
this chapter, reports only Lean's three standard axioms at that earlier commit (not re-run for this edition). The chapter did
not run the aggregate \texttt{lake build}, and the repository checker's overall verdict at that
commit is ``failed'', because two draft files are not imported (it found no axiom declaration). That
defect has since been fixed: the drafts were moved out of \leanpath{VOA/}, and a run of the checker on the
working tree of 28 October 2026 printed ``Repository check passed'' with 0 sanctioned axiom declarations.
Section~\ref{l02:sec:assumes} states exactly what the library rests on today and separates what was
machine-checked here from what is reported by the project. (Until October 2026 \lean{pbw} was an
axiom; that history is collected in Section~\ref{l02:sec:history}.)

Two jobs remain for this chapter, and they are the most important ones in it. The first is
\emph{fidelity}: \lean{VOA.pbw} is consumed by the affine vacuum modules and by the Zhu-theory
chain of Chapters~\ref{ch:L13} and~\ref{ch:L15}, and a referee must be able to agree that the Lean
statement \lean{PBWPrinciple} says what those consumers need it to say --- that it is the
Poincar\'e--Birkhoff--Witt consequence its name advertises, not something weaker or vacuous. A
proved statement can no longer make the library inconsistent, but a mis-stated one could still make
a downstream theorem mean less than its name. The second job is to explain the \emph{proof}: how
the injectivity of the bottom map of an induced module, and the nonvanishing of the affine vacuum,
follow from a relative version of the classical construction of Humphreys \S17.4.

Our verdict, argued in Section~\ref{l02:sec:axiom}, is that both fields of \lean{PBWPrinciple} are,
after every Lean definition they mention is unfolded, the standard consequences of the
Poincar\'e--Birkhoff--Witt theorem for Lie algebras over a field, valid without any finiteness
assumption, for every instance of their (very general) quantifiers. We check the quantifiers one by
one, check the degenerate instances by hand, list the slips in the encoding that would have made the
statement false or vacuous, and then present the proof that the library contains. For the results
that consume \lean{VOA.pbw} we say precisely which explicit hypotheses each still carries.

The chapter is organised as follows.
\begin{itemize}
\item Section~\ref{l02:sec:assumes} states what the library assumes today, and
  Section~\ref{l02:sec:history} records the history of the former PBW axiom.
\item Section~\ref{l02:sec:physics} gives the physical picture: currents, central extensions,
  stress tensors, and what PBW means for counting states.
\item Section~\ref{l02:sec:lean} recalls how Lie algebras, enveloping algebras and quotient
  modules look in Lean and Mathlib.
\item Sections~\ref{l02:sec:witt}--\ref{l02:sec:virmod} treat the Witt algebra, the Virasoro
  algebra and the universal Virasoro vacuum module.
\item Sections~\ref{l02:sec:affine}--\ref{l02:sec:grading} treat affine Kac--Moody algebras, their
  vacuum modules and the energy grading.
\item Section~\ref{l02:sec:axiom} is the heart: the classical PBW theorem, the exact Lean statement
  \lean{PBWPrinciple}, its fidelity check, and the proof of \lean{theorem pbw} by relative PBW.
\item Section~\ref{l02:sec:consequences} derives the consequences (nonvanishing of the vacuum, the
  simple quotient $L(0)$), lists the consumers of \lean{VOA.pbw}, and records the history of the
  axiom and of the axiom-free ``twin'' policy that accompanied it.
\item Section~\ref{l02:sec:pbwbasis} presents \lean{PBWBasis}, the full PBW basis of the vacuum
  module, which began as a \emph{named hypothesis} and is now proved for every $\ltwog$
  (\lean{pbwBasis_holds}).
\item Section~\ref{l02:sec:roots} presents root data, general affine Verma modules and the
  integrability target.
\item Section~\ref{l02:sec:lattices} treats lattices, cocycles, the twisted group algebra and parity
  briefly, pointing to Chapters~\ref{ch:02} and~\ref{ch:03} for the details.
\item Section~\ref{l02:sec:verification} is the verification record (the \texttt{\#print axioms} output
  behind every status box); the chapter ends with the referee's checklist and notes.
\end{itemize}
Vertex algebras themselves, modes, and the general trust model are those of
Chapter~\ref{ch:01}; the verification record of the whole library is Chapter~\ref{ch:15}.

%======================================================================
\section{What the library assumes today}\label{l02:sec:assumes}
%======================================================================

This section is in the present tense and states, in one place, what a reader must accept in order to
believe a theorem of the library as stated.

\paragraph{1. Lean's kernel and its three standard axioms.} Every proof is checked by Lean's kernel,
which accepts a proof only if it follows from the rules of Lean's type theory and from declared
axioms. Mathlib and this library use three axioms that belong to ordinary mathematics:
\begin{itemize}
\item \lean{propext} (propositional extensionality): two propositions that imply each other are equal;
\item \lean{Quot.sound}: in a quotient by a relation, related elements are equal; this is what makes
  quotient modules, quotient rings and the like behave as on paper (it also yields function
  extensionality);
\item \lean{Classical.choice}: from a proof that a type is nonempty one may pick an element; it gives the
  law of excluded middle, Zorn's lemma and the existence of bases.
\end{itemize}
Zermelo--Fraenkel set theory with choice (together with inaccessible cardinals, one for each universe
level used) provides a model of these, so they are the background logic of ordinary mathematics.
\texttt{\#print axioms foo} lists the axioms that the proof of \lean{foo} reaches; a report that is a
subset of these three is called \emph{\texttt{\#print axioms}-clean}.

\paragraph{2. Mathlib.} The pinned version is Mathlib v4.30.0; its definitions (Lie algebras,
universal enveloping algebras, tensor products, quotients, bases, \ldots) are taken to mean what their
names say. Chapter~\ref{ch:01} and Section~\ref{l02:sec:lean} explain the ones used here.

\paragraph{3. Fidelity of the library's own definitions.} A theorem is only as meaningful as the
definitions in its statement. Checking that each Lean definition means what the textbook means is the
job of this book's \emph{dictionary} boxes; it cannot be done by the machine.

\paragraph{4. No project axiom.} The library declares no axiom of its own; \lean{VOA.pbw} is a theorem
(Section~\ref{l02:sec:axiom}).

\paragraph{5. Explicit hypotheses of named theorems.} \texttt{\#print axioms} does \emph{not} report a
theorem's explicit hypotheses: a theorem \lean{bar (h : H) : C} proves only ``$H$ implies $C$'', however
clean its axiom list. Every status box of the book names these hypotheses. In this chapter and its
immediate consumers they are:
\begin{itemize}
\item \lean{highestVector H κ hκ b lam ≠ 0} for the general affine Verma module and its simple quotient
  (Section~\ref{l02:sec:roots}); true, but not proved in the library for general $\lambda$;
\item the hypothesis \lean{hKG} of the integrability wiring lemma, which is its whole content
  (Section~\ref{l02:sec:roots}; a tautology, not a result);
\item for Zhu's modularity theorem \lean{zhuModularity_}\allowbreak\lean{pinned_}\allowbreak\lean{of_}\allowbreak\lean{C2Cofinite_}\allowbreak\lean{voaPBW}:
  $C_2$-cofiniteness \lean{hC2} (Chapter~\ref{ch:L12}, Definition~\ref{l12:def:c2}), semisimplicity of Zhu's
  algebra \lean{hss} (Chapter~\ref{ch:L12}, Section~\ref{l12:sec:rationality}), the instance
  \lean{[Finite (Φ W)]} (finitely many irreducible modules, Definition~\ref{l12:def:phi}), and the family data
  \lean{hcard}, \lean{hcw}, \lean{hsep} (Chapter~\ref{ch:L15}, Theorem~\ref{l15:thm:headline}). These are
  the intended hypotheses of Zhu's theorem, not unproved inputs. The finiteness instance is in fact a
  consequence of \lean{hC2}, by the theorem \lean{VOA.Huang.ZhuFiniteIrr.}\allowbreak\lean{finite_Φ_of_C2Cofinite_pbw};
\item for \lean{zhuModularity_}\allowbreak\lean{strongly}\allowbreak\lean{Rational_}\allowbreak\lean{voaPBW}: strong rationality \lean{hSR}
  (Definition~\ref{l12:def:strong}), an intended hypothesis, and two \emph{literature theorems carried as
  hypotheses}, \lean{hFin} (Dong--Li--Mason finiteness) and \lean{hZhu} (Zhu's semisimplicity theorem)
  (Chapter~\ref{ch:L15}, Theorem~\ref{l15:thm:sr}). The wrapper's signature still carries them, but both are
  now discharged elsewhere in the library: \lean{dlmFinitelyManyIrredAdmissible_of_c2Cofinite} proves the
  first for every $C_2$-cofinite VOA and \lean{zhuAdmissibleSemisimple_holds} proves the second outright.
\end{itemize}
These inputs no longer stand between the reader and Zhu's theorem. The library proves Zhu modularity for the
canonical family of every $C_2$-cofinite \lean{AdmissiblyRational} VOA,
\lean{zhuModularity_}\allowbreak\lean{canonical_of_}\allowbreak\lean{C2Cofinite_}\allowbreak\lean{admissiblyRational} (Chapter~\ref{ch:L15}, Theorem~\ref{l15:thm:canonical}), from \lean{hC2} and
\lean{hrat : AdmissiblyRational W} alone: semisimplicity of Zhu's algebra is the theorem
\lean{isSemisimpleRing_}\allowbreak\lean{A_of_}\allowbreak\lean{admissiblyRational} (Theorem~\ref{l15:thm:zhuss}), and
\lean{[Finite (Φ W)]} is \lean{finite_Φ_of_C2Cofinite_pbw}. What remains open is only the relation of the
library's predicate \lean{AdmissiblyRational} to Dong--Li--Mason rationality as stated in the literature.
Results with no hypothesis beyond their data, such as \lean{VOA.pbw}, \lean{pbwBasis_holds} and
\lean{zhuBijective_}\allowbreak\lean{of_voaPBW}, are \emph{unconditional}. (The codomain of
\lean{zhuBijective_}\allowbreak\lean{of_voaPBW} is defined as an image, so only its injectivity has content; the
genuine Zhu correspondence is \lean{VOA.zhuCorrespondence_}\allowbreak\lean{of_C2Cofinite}, under \lean{hC2},
Chapter~\ref{ch:L13}.)

\paragraph{What the machine confirmed for this chapter, and what is reported.} \emph{Confirmed here}, at
an earlier commit (not re-run for this edition), with \texttt{lake env lean -DmaxSynthPendingDepth=3} on scratch files using the
existing compiled library: the \texttt{\#print axioms} reports of Section~\ref{l02:sec:verification}
(73 declarations in the displayed record, plus the five cocycle lemmas of Section~\ref{l02:sec:lattices}
in a separate probe), all exactly \lean{[propext, Classical.choice, Quot.sound]}; the verbatim occurrence
of every quoted Lean block in the source; and the output of \texttt{scripts/check\_repository.py}: 0 axiom
declarations, with overall verdict ``failed'' only because two draft files in \leanpath{VOA/Huang/} were not
imported by the umbrella file. (Re-run on the working tree of 28 October 2026, after the drafts were moved out
of \leanpath{VOA/}, the checker printed ``Repository check passed: 10353 library files scanned; \ldots;
0 sanctioned axiom declaration(s)''.) \emph{Not run here:} the aggregate \texttt{lake build} and the project's axiom
audit. \emph{Reported by the project} (its build record): the aggregate
\texttt{lake build VOA} was last run green on 26 October 2026 (13868 jobs) with the axiom audit passing.
A reader without access to the repository should treat that last item as the project's claim.

%======================================================================
\section{History of the former PBW axiom}\label{l02:sec:history}
%======================================================================

Nothing in this section is needed to read the rest of the chapter; it is kept because earlier editions
of this book, many library docstrings and other chapters refer to it.

\paragraph{The axiom and its removal.} Until 19 October 2026 the library had
exactly one project axiom, \lean{axiom pbw : PBWPrinciple}, a statement Lean accepted without proof, with
the same name and type as today's theorem. A revision of that date replaced it by the theorem of
Section~\ref{l02:sec:pbwproof}; no use site had to change. In the first edition of this chapter,
\texttt{\#print axioms} of every declaration depending on it listed \lean{pbw} in addition to the three
standard axioms (the record is kept at the end of Section~\ref{l02:sec:verification}). The commit message
of that revision states that the aggregate build was not yet certified at that time.

\paragraph{When the fidelity check was the only safeguard.}\label{l02:rem:axiomhistory}
Until 19 October 2026 the same statement was an axiom, and the fidelity check of
Section~\ref{l02:sec:axiom} was then the \emph{only} reason to believe the library consistent. The reasoning was: Lean's logic (with
\lean{propext}, \lean{Classical.choice}, \lean{Quot.sound}) has a set-theoretic model --- ZFC with a
sufficient supply of inaccessible cardinals, one for each universe --- in which types are sets, Lie
algebras are Lie algebras, and Mathlib's universal enveloping algebra, tensor product and quotients are
the usual ones; a statement that is true in that model, as Theorems~\ref{l02:thm:A} and~\ref{l02:thm:B}
show \lean{PBWPrinciple} to be, cannot introduce an inconsistency beyond whatever the base theory has.
The Lean proof of Section~\ref{l02:sec:pbwproof} has superseded this argument; it is kept here because
the earlier editions of this book, and many docstrings of the library, still refer to it.

\paragraph{The axiom-free twin policy.}
While \lean{pbw} was an axiom, the project
kept every result that needed it in two forms, so that \texttt{\#print axioms} would separate the
unconditional from the PBW-dependent. Each \lean{_voaPBW} wrapper had an \emph{axiom-free twin} that took
the PBW consequence as an explicit hypothesis --- \lean{Pr : BottomProjection …} or a
\lean{PBWFreeness} input --- and reported the three standard axioms only; for instance
\lean{zhuBijective} (twin of \lean{zhuBijective_}\allowbreak\lean{of_voaPBW}) and
\lean{zhuModularity_}\allowbreak\lean{pinned_}\allowbreak\lean{of_}\allowbreak\lean{C2Cofinite} (twin of the \lean{_voaPBW} version). The
same policy introduced the named hypothesis \lean{PBWBasis} (Section~\ref{l02:sec:pbwbasis}) rather than a
second axiom. Since \lean{pbw} became a theorem the distinction is moot: the hypothesis
\lean{Pr : BottomProjection (parabolicRepOfAMod W (regularAMod W))} of a twin is supplied, for every
\lean{W}, by the definition \lean{voaPBWBottom}\allowbreak\lean{Projection} built from the now-proved
\lean{voaPBWBottom}\allowbreak\lean{Injective}, and each \lean{_voaPBW} wrapper is exactly its twin with that
hypothesis discharged. The twins remain in the library and remain correct; their \lean{Pr} (or
\lean{PBWFreeness}, or \lean{h : PBWBasis …}) hypotheses are now \emph{redundant} --- always satisfiable by
a theorem of the library --- and a referee may read each twin together with its \lean{_voaPBW} wrapper as
one unconditional-in-PBW result. The project lists dropping such redundant hypotheses as an optional
clean-up; it had not been done when this was checked at an earlier commit (not re-checked for this edition).

%======================================================================
\section{The physical picture}\label{l02:sec:physics}
%======================================================================

\subsection{Currents and the affine algebra}

Consider a two-dimensional quantum field theory with a continuous internal symmetry group $G$, for
instance the Wess--Zumino--Witten (WZW) model, a sigma model whose target is the group manifold
$G$ itself. Noether's theorem gives conserved currents $J^a_\mu$, one for each element $x_a$ of a
basis of the Lie algebra $\ltwog$ of $G$. In a conformal field theory with holomorphic
factorisation the holomorphic component $J^a(z)$ is a holomorphic field of conformal weight one,
and it has a Laurent expansion
\begin{equation*}
 J^a(z) = \sum_{n\in\Z} J^a_n\, z^{-n-1}.
\end{equation*}
The modes $J^a_n$ are operators on the space of states. The operator product expansion of two
currents in the WZW model at level $k$ is
\begin{equation}\label{l02:eq:JJope}
 J^a(z)\,J^b(w) \sim \frac{k\,\kappa(x_a,x_b)}{(z-w)^2} + \frac{\sum_c f_{ab}{}^c\,J^c(w)}{z-w},
\end{equation}
where $[x_a,x_b] = \sum_c f_{ab}{}^c x_c$ and $\kappa$ is an invariant symmetric bilinear form on
$\ltwog$ (for simple $\ltwog$ a multiple of the Killing form); the symbol $\sim$ means ``equal up to terms
regular at $z = w$''. Extracting modes by contour integrals, $J^a_m = \oint\frac{dz}{2\pi i}z^mJ^a(z)$, turns
\eqref{l02:eq:JJope} into the commutation relations
\begin{equation}\label{l02:eq:affinemodes}
 [J^a_m, J^b_n] = \sum_c f_{ab}{}^c\, J^c_{m+n} + k\, m\, \kappa(x_a,x_b)\,\delta_{m+n,0}.
\end{equation}
The first term is just the Lie bracket of the loop algebra $\ltwog\otimes\C[t,t^{-1}]$ (where
$J^a_n$ corresponds to $x_a\otimes t^n$). The second term is a \emph{Schwinger term}: a c-number
that appears because the product of two quantum fields at the same point must be regularised.
Abstractly it is a central extension: one adjoins a central element $C$, writes the bracket as
$[x\otimes t^m, y\otimes t^n] = [x,y]\otimes t^{m+n} + m\,\kappa(x,y)\,\delta_{m+n,0}\,C$, and in
a representation $C$ acts as the scalar $k$, the \emph{level}. The resulting Lie algebra is the
affine Kac--Moody algebra $\ltwoghat$. One usually also adjoins the derivation $\ell_0$ with
$[\ell_0, x\otimes t^n] = n\,(x\otimes t^n)$; on states it measures (minus) the energy contributed
by the currents.

\subsection{The stress tensor and the Virasoro algebra}

Every conformal field theory has a stress tensor $T(z) = \sum_n L_n z^{-n-2}$, the conserved current
of conformal symmetry, with operator product
\begin{equation*}
 T(z)T(w) \sim \frac{c/2}{(z-w)^4} + \frac{2T(w)}{(z-w)^2} + \frac{\partial T(w)}{z-w}.
\end{equation*}
In modes this is the Virasoro algebra
\begin{equation}\label{l02:eq:vir}
 [L_m,L_n] = (m-n)L_{m+n} + \frac{c}{12}(m^3-m)\,\delta_{m+n,0}.
\end{equation}
Without the last term, \eqref{l02:eq:vir} is the Lie algebra of infinitesimal conformal
transformations of the punctured plane, the Witt algebra, realised by the vector fields
$\ell_n = -z^{n+1}\partial_z$. The central term is a quantum anomaly: the classical conformal
symmetry is realised projectively on the quantum state space, and the obstruction is measured by
the central charge $c$. That $(m^3-m)/12$ is the only possible shape of such a term (up to
redefinitions $L_0\mapsto L_0 + \text{const}$) is the statement that the second cohomology of the
Witt algebra with trivial coefficients is one-dimensional; the library proves that the given term is
a cocycle but not this uniqueness.

\subsection{Vacuum modules and what PBW says}

In both cases the space of states of the ``vacuum sector'' is built in the same way. There is a
vacuum state $\ltwovac$ with the following properties: $J^a_n\ltwovac = 0$ for $n\ge0$ (regularity of $J^a(z)\ltwovac$ at $z = 0$), and
$L_n\ltwovac = 0$ for $n\ge-1$ (regularity of $T(z)\ltwovac$ at $z=0$, including translation
invariance $L_{-1}\ltwovac = 0$). The other states are obtained by applying creation operators,
\begin{equation*}
 J^{a_1}_{-n_1}\cdots J^{a_r}_{-n_r}\ltwovac \quad (n_i\ge1), \qquad
 L_{-n_1}\cdots L_{-n_r}\ltwovac\quad(n_i\ge2).
\end{equation*}
A physicist counting the states of the vacuum sector assumes, without thinking about it, that
these states \emph{with the creation operators in a fixed order} are linearly independent, so that
the partition function of the universal vacuum module of $\ltwoghat$ is
$\prod_{n\ge1}(1-q^n)^{-\dim\ltwog}$ (up to the conventional prefactor $q^{-c/24}$). The
mathematical justification of this counting is the Poincar\'e--Birkhoff--Witt (PBW) theorem: the
universal enveloping algebra of a Lie algebra has a basis of ordered monomials. The commutation
relations allow any monomial to be reordered (that is the easy half), and PBW says that reordering
is the \emph{only} source of relations --- there are no ``unexpected'' linear dependencies, and in
particular the vacuum module is not zero. This last statement, that the construction does not
collapse, is exactly what the first field of the library's theorem \lean{VOA.pbw} asserts for affine
algebras (it was an axiom until October 2026; it is now proved).

There is an alternative physical route to the same conclusion: if one can write down \emph{some}
concrete representation in which the vacuum conditions hold and the vacuum is nonzero (a free
field realisation: a free boson for Virasoro, Wakimoto's $\beta\gamma$-plus-boson realisation for
$\widehat{\mathfrak{sl}}_2$), then the universal vacuum module maps onto something nonzero and
cannot be zero. The library uses this route for the Virasoro algebra (where it is the only proof)
and for $\widehat{\mathfrak{sl}}_2$ and $\widehat{\mathfrak{sl}}_3$ (where it was developed, while PBW
was still an axiom, to keep those parts of the library independent of it; see
Sections~\ref{l02:sec:virmod} and~\ref{l02:sec:consequences}).

\subsection{Lattices}

The third ingredient, lattices, enters through compactified bosons. A free boson on a circle of
radius $R$, or more generally $r$ bosons on a torus $\mathbb{R}^r/2\pi L$, has momentum and winding
sectors labelled by points of a lattice; for the chiral (holomorphic) theory one needs an even
lattice, and the vertex operators $e^{i\alpha\cdot X}$ of the different sectors are mutually local
only after they are multiplied by \emph{cocycle factors} (``Klein factors'') $\varepsilon(\alpha,\beta) =
\pm1$. The FLM cocycle of Section~\ref{l02:sec:lattices} is the mathematical form of these Klein
factors.

%======================================================================
\section{Lie algebras in Lean: what the reader needs}\label{l02:sec:lean}
%======================================================================

We collect the Mathlib notions that occur in every statement of this chapter. Chapter~\ref{ch:01},
Section~\ref{ch01:sec:lean}, explains the general vocabulary (types, structures, typeclasses,
\lean{Submodule}, linear maps \lean{→ₗ[ℂ]}); here we add the Lie-theoretic part.

\begin{itemize}
\item \lean{[LieRing 𝔤] [LieAlgebra ℂ 𝔤]}: the type \lean{𝔤} carries an additive group with a
  bracket \lean{⁅x, y⁆} which is biadditive, alternating (\lean{⁅x, x⁆ = 0}) and satisfies the
  Jacobi identity in the Leibniz form \lean{⁅x, ⁅y, z⁆⁆ = ⁅⁅x, y⁆, z⁆ + ⁅y, ⁅x, z⁆⁆}; and a
  $\C$-vector-space structure for which the bracket is $\C$-bilinear. So \lean{𝔤} is a complex Lie
  algebra, of \emph{arbitrary} dimension: no finiteness is implied.
\item \lean{L₁ →ₗ⁅ℂ⁆ L₂}: a homomorphism of Lie algebras ($\C$-linear and bracket-preserving).
\item \lean{Module.End ℂ V}: the associative algebra $\End_\C(V)$ of $\C$-linear maps $V\to V$.
  Mathlib makes every associative ring a Lie ring with the commutator bracket
  \lean{⁅A, B⁆ = A * B - B * A} (the lemma \lean{LieRing.}\allowbreak\lean{of_}\allowbreak\lean{associative_}\allowbreak\lean{ring_}\allowbreak\lean{bracket}). So a Lie
  homomorphism \lean{ρ : 𝔤 →ₗ⁅ℂ⁆ Module.End ℂ V} is precisely a representation of $\ltwog$ on $V$:
  $\rho([x,y]) = \rho(x)\rho(y)-\rho(y)\rho(x)$.
\item \lean{LieSubalgebra ℂ 𝔤}: a Lie subalgebra $\ltwop\subseteq\ltwog$; used as a type, \lean{p}
  denotes the elements of $\ltwop$ with the restricted bracket, and \lean{(x : 𝔤)} is the inclusion.
  \lean{LieSubalgebra.lieSpan ℂ _ S} is the Lie subalgebra generated by a set $S$.
\item \lean{UniversalEnvelopingAlgebra ℂ 𝔤}: the universal enveloping algebra $U(\ltwog)$,
  constructed by Mathlib as the quotient of the tensor algebra $T(\ltwog)$ by the two-sided ideal
  generated by $x\otimes y - y\otimes x - [x,y]$ (more precisely by the relation
  $\iota[x,y] + \iota y\,\iota x \sim \iota x\,\iota y$, which is the same thing). The map
  \lean{ι ℂ : 𝔤 →ₗ⁅ℂ⁆ UniversalEnvelopingAlgebra ℂ 𝔤} is the canonical Lie homomorphism, and
  \lean{lift} is the universal property: every Lie homomorphism $\ltwog\to A$ into an associative
  algebra extends uniquely to an algebra homomorphism $U(\ltwog)\to A$. \emph{Mathlib proves
  nothing more about $U(\ltwog)$}: not the PBW theorem, not even that \lean{ι} is injective. (We
  checked \leanpath{Mathlib/\allowbreak Algebra/\allowbreak Lie/\allowbreak UniversalEnveloping.lean} at the pinned
  version; the files \leanpath{Free.lean} and \leanpath{SerreConstruction.lean} of the same directory
  mention PBW only as a missing ingredient.)
\item \lean{Submodule R R} for a ring $R$: a \emph{left ideal} of $R$ ($R$ acting on itself by left
  multiplication). \lean{Submodule.span R S} is the left ideal $\sum_{s\in S}Rs$ generated by $S$.
  The quotient type \lean{R ⧸ I} is the left $R$-module $R/I$, and
  \lean{Submodule.Quotient.mk a} is the class $[a]$.
\item \lean{ℤ →₀ 𝔤}: finitely supported functions $\Z\to\ltwog$ (``\lean{Finsupp}''). The element
  \lean{single n x} is the function with value $x$ at $n$ and $0$ elsewhere. We read
  $f\in\lean{ℤ →₀ 𝔤}$ as the Laurent polynomial $\sum_n f(n)\otimes t^n\in\ltwog\otimes\C[t,t^{-1}]$;
  this is an isomorphism of vector spaces, because a Laurent polynomial has finitely many nonzero
  coefficients.
\item \lean{Type u}: Lean's types live in a hierarchy of universes \lean{Type 0}, \lean{Type 1},
  \ldots; a statement quantifying over \lean{𝔤 : Type u} speaks about all Lie algebras whose
  underlying type lives in universe $u$. Set-theoretically one can think of universes as Grothendieck
  universes; every classical theorem of algebra holds in each of them.
\end{itemize}

%======================================================================
\section{The Witt algebra and the Virasoro algebra}\label{l02:sec:witt}
%======================================================================

All of this section is in \leanpath{VOA/\allowbreak Foundations/\allowbreak Virasoro.lean}.

\subsection{The Witt algebra}

\begin{definition}[Witt algebra]\label{l02:def:witt}
The \emph{Witt algebra} $\ltwoW$ is the complex vector space with basis $\{\ell_n\}_{n\in\Z}$ and
bracket $[\ell_m,\ell_n] = (m-n)\,\ell_{m+n}$, extended bilinearly. It is realised faithfully on the
Laurent polynomials $\C[z,z^{-1}]$ by $\ell_n = -z^{n+1}\frac{d}{dz}$, i.e.\
$\ell_n\cdot z^k = -k\,z^{k+n}$.
\end{definition}

File \leanpath{VOA/Foundations/Virasoro.lean}.
\begin{leancode}
abbrev Witt : Type := ℤ →₀ ℂ

namespace Witt

noncomputable def gen (n : ℤ) : Witt := single n 1

noncomputable def op (n : ℤ) : Module.End ℂ (ℤ →₀ ℂ) :=
  Finsupp.lsum ℂ fun k => LinearMap.toSpanSingleton ℂ (ℤ →₀ ℂ) ((-(k : ℂ)) • single (k + n) 1)

noncomputable def bracketWith (m : ℤ) : Witt →ₗ[ℂ] Witt :=
  Finsupp.lsum ℂ fun n => LinearMap.toSpanSingleton ℂ Witt (single (m + n) ((m : ℂ) - n))

noncomputable def bracketₗ : Witt →ₗ[ℂ] Witt →ₗ[ℂ] Witt :=
  Finsupp.lsum ℂ fun m => LinearMap.toSpanSingleton ℂ _ (bracketWith m)

noncomputable def rep : Witt →ₗ[ℂ] Module.End ℂ (ℤ →₀ ℂ) :=
  Finsupp.lsum ℂ fun n => LinearMap.toSpanSingleton ℂ _ (op n)
\end{leancode}
together with the Lie structure and the defining relation:
\begin{leancode}
theorem rep_injective : Function.Injective rep := by

noncomputable instance : LieRing Witt where
  add_lie x y z := by simp only [bracket_def, map_add, LinearMap.add_apply]
  lie_add x y z := by simp only [bracket_def, map_add]
  lie_self x := rep_injective (by rw [rep_bracket, lie_self, map_zero])
  leibniz_lie x y z := rep_injective (by simp only [rep_bracket, map_add]; exact leibniz_lie _ _ _)

theorem gen_bracket (m n : ℤ) : ⁅gen m, gen n⁆ = ((m : ℂ) - n) • gen (m + n) := by
\end{leancode}

\begin{dictionary}
\textbf{Decoding.} The carrier \lean{Witt} is \emph{literally} \lean{ℤ →₀ ℂ}, finitely supported
functions $\Z\to\C$; \lean{gen n = single n 1} is the basis vector $\ell_n$, and an element
\lean{w} is $\sum_n w(n)\ell_n$. \lean{op n} is the linear map on \lean{ℤ →₀ ℂ} (now read as
$\C[z,z^{-1}]$, \lean{single k 1} $= z^k$) sending $z^k\mapsto -k\,z^{k+n}$; this is
$-z^{n+1}\,d/dz$. \lean{bracketWith m} sends $\ell_n\mapsto (m-n)\ell_{m+n}$, and \lean{bracketₗ}
assembles these into the bilinear bracket; \lean{rep} sends $\ell_n\mapsto$ \lean{op n}.

\textbf{Why the Lie axioms hold.} Instead of proving the Jacobi identity by a direct computation,
the file proves (\lean{rep_bracket}) that \lean{rep} turns the bracket into the commutator of
operators, and (\lean{rep_injective}) that \lean{rep} is injective; then alternation and Jacobi are
transported from the associative algebra $\End(\C[z,z^{-1}])$. The faithfulness proof evaluates
\lean{rep x} on $z = $ \lean{single 1 1}: $\ell_n z = -z^{n+1}$, so the coefficients of $x$ are read
off (shifted by one and negated).

\textbf{Conventions.} The sign $\ell_n = -z^{n+1}\partial_z$ is the standard one for which
$[\ell_m,\ell_n] = (m-n)\ell_{m+n}$; with $+z^{n+1}\partial_z$ the bracket would have the opposite
sign. The Lean bracket \lean{gen_bracket} is exactly Definition~\ref{l02:def:witt}.

\textbf{Could it be vacuous?} No: \lean{Witt} is a concrete type, its bracket is computed on basis
vectors by \lean{bracketₗ_single}, and \lean{gen_bracket} pins it to the textbook formula. The
structure is the Witt algebra, not something weaker.
\end{dictionary}

\begin{physics}
The Witt algebra is the Lie algebra of infinitesimal holomorphic coordinate changes
$z\mapsto z + \epsilon z^{n+1}$ of the punctured plane: $\ell_{-1}$ generates translations, $\ell_0$
dilatations (and rotations), $\ell_1$ special conformal transformations; $\ell_{-1},\ell_0,\ell_1$
span $\mathfrak{sl}_2$, the global conformal algebra of the Riemann sphere. It is the classical
conformal symmetry algebra of a two-dimensional field theory, before quantisation.
\end{physics}

\begin{proof}[Proof on paper]
Write $D_n = -z^{n+1}\partial_z$. Then $[D_m,D_n]z^k = D_m(-kz^{k+n}) - D_n(-kz^{k+m}) =
k(k+n)z^{k+m+n} - k(k+m)z^{k+m+n} = (n-m)k\,z^{k+m+n} = (m-n)D_{m+n}z^k$. Hence the linear map
$\ell_n\mapsto D_n$ intertwines the bracket with the commutator. It is injective because
$D_n z = -z^{n+1}$ and the $z^{n+1}$ are linearly independent. Therefore the bracket on
$\ltwoW$ is the pull-back of the commutator bracket along an injective linear map, and inherits
bilinearity, alternation and the Jacobi identity.
\end{proof}

\begin{status}
\lean{VOA.Witt.}\allowbreak\lean{gen_bracket}, \lean{VOA.Witt.}\allowbreak\lean{rep_injective}: \texttt{\#print axioms} gives
\lean{[propext, Classical.choice, Quot.sound]} (run for this chapter, see
Section~\ref{l02:sec:verification}). Unconditional.
\end{status}

\subsection{The Gel'fand--Fuks cocycle and the Virasoro algebra}

\begin{definition}[Virasoro algebra]\label{l02:def:vir}
Let $\psi : \ltwoW\times\ltwoW\to\C$ be the bilinear map with
$\psi(\ell_m,\ell_n) = \frac{1}{12}(m^3-m)\,\delta_{m+n,0}$. It is alternating and satisfies the
2-cocycle identity $\psi(x,[y,z]) = \psi([x,y],z) + \psi(y,[x,z])$. The \emph{Virasoro algebra}
$\ltwoVir$ is the central extension $\ltwoW\oplus\C C$ with bracket
$[(x,a),(y,b)] = ([x,y],\psi(x,y))$; writing $L_n = (\ell_n,0)$ and $C = (0,1)$,
\begin{equation*}
 [L_m,L_n] = (m-n)L_{m+n} + \frac{m^3-m}{12}\,\delta_{m+n,0}\,C,\qquad [C,\cdot\,] = 0.
\end{equation*}
A representation of $\ltwoVir$ on $V$ has \emph{central charge} $c$ if $C$ acts as $c\cdot\mathrm{id}$.
\end{definition}

File \leanpath{VOA/Foundations/Virasoro.lean}.
\begin{leancode}
noncomputable def gfCoeff (m n : ℤ) : ℂ := if m + n = 0 then ((m : ℂ) ^ 3 - m) / 12 else 0

theorem gfCoeff_cocycle (a b c : ℤ) :
    ((b : ℂ) - c) * gfCoeff a (b + c)
      = ((a : ℂ) - b) * gfCoeff (a + b) c + ((a : ℂ) - c) * gfCoeff b (a + c) := by

theorem cocycle_cond (x y z : Witt) :
    cocycleₗ x ⁅y, z⁆ = cocycleₗ ⁅x, y⁆ z + cocycleₗ y ⁅x, z⁆ := by

noncomputable def gfCochain : twoCochain ℂ Witt ℂ := ⟨Witt.cocycleₗ, Witt.cocycleₗ_self⟩

noncomputable def gfCocycle : twoCocycle ℂ Witt ℂ :=
  ⟨gfCochain, (mem_twoCocycle_iff_of_trivial (a := gfCochain)).mpr
    fun x y z => Witt.cocycle_cond x y z⟩

abbrev Virasoro : Type := LieAlgebra.ofTwoCocycle gfCocycle

noncomputable def L (n : ℤ) : Virasoro := ofProd gfCocycle (Witt.gen n, 0)

noncomputable def C : Virasoro := ofProd gfCocycle (0, 1)

theorem virasoro_bracket (m n : ℤ) :
    ⁅L m, L n⁆ = ((m : ℂ) - n) • L (m + n)
      + (if m + n = 0 then (((m : ℂ) ^ 3 - m) / 12) • C else 0) := by

structure VirRep (V : Type*) [AddCommGroup V] [Module ℂ V] where
  ρ : Virasoro →ₗ⁅ℂ⁆ Module.End ℂ V
  centralCharge : ℂ
  acts_central : ρ C = centralCharge • LinearMap.id
\end{leancode}

\begin{dictionary}
\textbf{Decoding.} \lean{gfCoeff m n} is $\psi(\ell_m,\ell_n)$. \lean{cocycleₗ} is its bilinear
extension (defined exactly like \lean{bracketₗ}). \lean{twoCochain ℂ Witt ℂ} and
\lean{twoCocycle ℂ Witt ℂ} are Mathlib's alternating bilinear maps $\ltwoW\times\ltwoW\to\C$ and,
among them, the 2-cocycles for the Chevalley--Eilenberg differential with coefficients in the
trivial module $\C$ (the trivial action is installed by the \lean{LieRingModule Witt ℂ} instance with
zero bracket). \lean{mem_}\allowbreak\lean{twoCocycle_}\allowbreak\lean{iff_}\allowbreak\lean{of_}\allowbreak\lean{trivial} says that for trivial coefficients the cocycle
condition is $\psi(x,[y,z]) = \psi([x,y],z)+\psi(y,[x,z])$, which is \lean{cocycle_cond}.
\lean{LieAlgebra.}\allowbreak\lean{ofTwoCocycle} is Mathlib's central extension of a Lie algebra by a 2-cocycle;
\lean{ofProd} identifies it with the product $\ltwoW\times\C$, so \lean{L n} is $(\ell_n,0)$ and
\lean{C} is $(0,1)$. \lean{virasoro_bracket} is the displayed relation of
Definition~\ref{l02:def:vir}, including the normalisation $1/12$, and \lean{C_central} (in
\leanpath{VirasoroModule.lean}) says $C$ is central. \lean{VirRep V} bundles a representation $\rho$
with a number \lean{centralCharge} and the proof that $\rho(C) = c\cdot\mathrm{id}$.

\textbf{Same as the textbook?} Yes, with the standard normalisation: the anomaly coefficient
$(m^3-m)/12$ makes $L_{-1},L_0,L_1$ span an $\mathfrak{sl}_2$ with no central term
(\lean{[L₁, L₋₁] = 2L₀}, one of the file's smoke tests) and gives $[L_2,L_{-2}] = 4L_0 + C/2$
(the other smoke test).

\textbf{Could it be vacuous or wrong?} The scalar identity \lean{gfCoeff_cocycle} is the whole
content of the cocycle condition on basis vectors. If $a+b+c\ne0$ every term vanishes; if
$c=-a-b$ it is a polynomial identity in $a,b$, which the proof checks by \lean{ring}. A wrong
coefficient would have been caught: with $m^5$ in place of $m^3-m$ the identity at $(a,b)=(1,2)$
reads $5 = -115$, and the definition \lean{gfCocycle} would not have compiled, because it requires a
proof of the cocycle condition. (The cocycle condition does \emph{not} pin the normalisation:
$\alpha(m^3-m)+\beta m$ is a cocycle for all $\alpha,\beta$, the $\beta m$ part being a coboundary.
The choice $\alpha = 1/12$, $\beta=0$ is the convention of the physics literature, fixed by
\lean{virasoro_bracket}.)
\end{dictionary}

\begin{physics}
$\ltwoVir$ is the algebra of modes of the stress tensor $T(z) = \sum L_nz^{-n-2}$ of any
two-dimensional conformal field theory; the central element is the anomaly, and its value $c$ is
the central charge (Casimir energy on the cylinder $-c/24$, trace anomaly
$\langle T^\mu{}_\mu\rangle\propto c\,R$ on curved surfaces). The Schwinger term
$\frac{c}{12}(m^3-m)$ is what the $(z-w)^{-4}$ term of the $TT$ operator product expansion becomes in
modes. \lean{VirRep} is the statement ``this space of states carries the conformal symmetry with
central charge $c$''.
\end{physics}

\begin{proof}[Proof on paper of the cocycle identity]
On basis vectors $x = \ell_a$, $y=\ell_b$, $z=\ell_c$ the identity reads
$(b-c)\psi(\ell_a,\ell_{b+c}) = (a-b)\psi(\ell_{a+b},\ell_c) + (a-c)\psi(\ell_b,\ell_{a+c})$,
using $[\ell_b,\ell_c] = (b-c)\ell_{b+c}$ etc. All three $\psi$-values vanish unless $a+b+c=0$. If
$c = -a-b$, multiply by 12; the left side is $(a+2b)(a^3-a)$, the right side is
$(a-b)((a+b)^3-(a+b)) + (2a+b)(b^3-b)$, and expanding both gives
$a^4+2a^3b-a^2-2ab$ on each side. Alternation is $\psi(\ell_m,\ell_{-m}) = -\psi(\ell_{-m},\ell_m)$,
i.e.\ $m^3-m$ is odd in $m$.
\end{proof}

\begin{status}
\lean{VOA.Witt.}\allowbreak\lean{cocycle_}\allowbreak\lean{cond}, \lean{VOA.virasoro_}\allowbreak\lean{bracket}, \lean{VOA.C_central}:
\lean{[propext, Classical.choice, Quot.sound]} (run for this chapter). Unconditional.
Not formalised: uniqueness of the central extension, $H^2(\ltwoW;\C)\cong\C$ (the file's docstring
records it as a stretch goal).
\end{status}

%======================================================================
\section{The universal Virasoro vacuum module}\label{l02:sec:virmod}
%======================================================================

All declarations are in \leanpath{VOA/\allowbreak Foundations/\allowbreak VirasoroModule.lean} except the
nonvanishing theorem, which is in \leanpath{VOA/\allowbreak Examples/\allowbreak VirasoroNonvanishing.lean} and is
documented in detail in Chapter~\ref{ch:L06}; we include it here because it shows the free-field
pattern that let parts of the library avoid the former PBW axiom, and because the Virasoro algebra
is not covered by \lean{VOA.pbw} at all (its vacuum module is not an affine vacuum module).

\begin{definition}[Universal Virasoro vacuum module]\label{l02:def:viruniv}
Let $J(c)\subseteq U(\ltwoVir)$ be the left ideal generated by $\{L_n : n\ge -1\}\cup\{C-c\}$. The
\emph{universal vacuum module} at central charge $c$ is $V(c,0) = U(\ltwoVir)/J(c)$ with
$\ltwoVir$ acting by left multiplication, and its vacuum vector is $\ltwovac = [1]$.
\end{definition}

File \leanpath{VOA/Foundations/VirasoroModule.lean}.
\begin{leancode}
abbrev UVir : Type := UniversalEnvelopingAlgebra ℂ Virasoro

def vacGens (c : ℂ) : Set UVir :=
  {a | (∃ n : ℤ, -1 ≤ n ∧ a = ι ℂ (L n)) ∨ a = ι ℂ C - algebraMap ℂ UVir c}

noncomputable def vacIdeal (c : ℂ) : Submodule UVir UVir :=
  Submodule.span _ (vacGens c)

abbrev VirUniv (c : ℂ) : Type := UVir ⧸ vacIdeal c

noncomputable def virVacRho (c : ℂ) : Virasoro →ₗ⁅ℂ⁆ Module.End ℂ (VirUniv c) :=
  (Algebra.lsmul ℂ ℂ (VirUniv c)).toLieHom.comp (ι ℂ)

noncomputable def vacVec (c : ℂ) : VirUniv c := Submodule.Quotient.mk 1

theorem vac_acts_central (c : ℂ) : virVacRho c C = c • LinearMap.id := by

theorem vac_pos_annihilate (c : ℂ) (n : ℤ) (hn : -1 ≤ n) :
    virVacRho c (L n) (vacVec c) = 0 := by

theorem vac_cyclic (c : ℂ) :
    Submodule.span UVir {vacVec c} = ⊤ := by
\end{leancode}
and, from \leanpath{VOA/\allowbreak Examples/\allowbreak VirasoroNonvanishing.lean},
File \leanpath{VOA/Examples/VirasoroNonvanishing.lean}.
\begin{leancode}
theorem vacVec_ne_zero (c : ℂ) : VirModule.vacVec c ≠ 0 := by
\end{leancode}

\begin{dictionary}
\textbf{Decoding.} \lean{vacIdeal c} is \lean{Submodule.span} over \lean{UVir} $=U(\ltwoVir)$ of
the generators, i.e.\ the \emph{left} ideal $\sum_{n\ge-1}U(\ltwoVir)L_n + U(\ltwoVir)(C-c)$
(\lean{algebraMap ℂ UVir c} is the scalar $c\cdot 1$). \lean{VirUniv c} is the quotient module.
\lean{Algebra.lsmul ℂ ℂ V} is the action of $U(\ltwoVir)$ on $V$ by left multiplication, as an
algebra homomorphism into $\End_\C V$; composing with \lean{ι} gives the Lie action, so
\lean{virVacRho c x [a] = [ι x * a]} (\lean{virVacRho_mk}). \lean{vac_cyclic} says the
$U(\ltwoVir)$-submodule generated by $\ltwovac$ is everything (\lean{⊤} is the top submodule).

\textbf{The three properties are exactly the defining ones}: $C$ acts by $c$ on all of $V(c,0)$
(not only on $\ltwovac$ --- this uses that $C$ is central, so $[(C-c)a] = [a(C-c)] = 0$),
$L_n\ltwovac = 0$ for $n\ge -1$, and $\ltwovac$ generates. The name \lean{VirUniv} (``universal'')
was chosen after an earlier docstring wrongly called this module $L(c,0)$, which denotes the
\emph{simple} quotient; the file records the renaming.

\textbf{Vacuity.} Every statement in the code block before \lean{vacVec_ne_zero} would be true
if $V(c,0)$ were the zero module. The content that excludes this is \lean{vacVec_ne_zero}: the
vacuum vector is not zero, so $V(c,0)\ne0$. It is proved for \emph{every} $c\in\C$, by a free-field
argument that does not use \lean{VOA.pbw}. (The relative PBW theorem of
Section~\ref{l02:sec:axiom} would also give it, by inducing from the subalgebra spanned by the
$L_n$, $n\ge-1$, and $C$; the library does not take that route.)
\end{dictionary}

\begin{physics}
$V(c,0)$ is the space of all states one can make from the vacuum with the stress tensor alone:
$L_{-n_1}\cdots L_{-n_r}\ltwovac$, $n_i\ge2$ --- the conformal family of the identity operator,
before any null states are removed. The condition $L_{-1}\ltwovac = 0$ is translation invariance
of the vacuum; $L_0\ltwovac = 0$ says the vacuum has energy zero; $L_n\ltwovac=0$ for $n\ge1$
says $T(z)\ltwovac$ is regular at $z=0$.
\end{physics}

\begin{proof}[Proof on paper of \lean{vacVec_ne_zero}]
It suffices to find a representation $\rho$ of $\ltwoVir$ on a space $F$ with central charge $c$
and a vector $v\ne0$ in $F$ with $\rho(L_n)v = 0$ for all $n\ge-1$. Indeed, by the universal property
$\rho$ extends to an algebra map $\tilde\rho : U(\ltwoVir)\to\End F$, and $a\mapsto\tilde\rho(a)v$
kills every generator of $J(c)$, hence (being left $U$-linear) all of $J(c)$; if $1\in J(c)$ we
would get $v = \tilde\rho(1)v = 0$. The library takes $F$ to be the Fock space of one free boson
with the ``linear dilaton'' (background charge) stress tensor
$L^\lambda_n = \frac12\sum_m{:}a_{n-m}a_m{:} + \lambda(n+1)a_n$, which has central charge
$1-12\lambda^2$ (the term $\lambda(n+1)a_n$ vanishes on the Fock vacuum for $n=-1$ because of the factor
$n+1$, and for $n\ge0$ because $a_n$ annihilates it); since $\C$ is algebraically closed every $c$ is of this
form, and $v$ is the Fock vacuum. (Chapter~\ref{ch:L06} gives the details.)
\end{proof}

\begin{status}
\lean{VOA.}\allowbreak\lean{VirModule.}\allowbreak\lean{vac_}\allowbreak\lean{acts_}\allowbreak\lean{central}, \lean{VOA.}\allowbreak\lean{VirModule.}\allowbreak\lean{vac_}\allowbreak\lean{pos_}\allowbreak\lean{annihilate},
\lean{VOA.}\allowbreak\lean{VirModule.}\allowbreak\lean{vac_}\allowbreak\lean{cyclic}, \lean{VOA.vacVec_}\allowbreak\lean{ne_zero}, \lean{VOA.}\allowbreak\lean{vacuumSpace_}\allowbreak\lean{nontrivial} (the Virasoro version, $V(c,0)$ is nontrivial,
in \leanpath{Examples/\allowbreak VirasoroNonvanishing.lean}; not the affine \lean{AffineLie.vacuumSpace_nontrivial}):
\lean{[propext, Classical.choice, Quot.sound]} (run for this chapter). Unconditional; the
proof does not use \lean{VOA.pbw}.
\end{status}

%======================================================================
\section{Affine Kac--Moody algebras}\label{l02:sec:affine}
%======================================================================

All declarations of this section and the next two are in
\leanpath{VOA/\allowbreak Foundations/\allowbreak AffineLie.lean}.

\subsection{Invariant forms}

\begin{definition}[Invariant form]\label{l02:def:invariant}
A bilinear form $\kappa$ on a Lie algebra $\ltwog$ is \emph{invariant} (in the sense used here) if it
is symmetric, $\kappa(x,y) = \kappa(y,x)$, and ad-invariant, $\kappa([x,y],z) = \kappa(x,[y,z])$.
\end{definition}

File \leanpath{VOA/Foundations/AffineLie.lean}.
\begin{leancode}
structure _root_.LinearMap.BilinForm.IsInvariant {𝔤 : Type*} [LieRing 𝔤] [LieAlgebra ℂ 𝔤]
    (κ : LinearMap.BilinForm ℂ 𝔤) : Prop where
  symm : ∀ x y : 𝔤, κ x y = κ y x
  invariant : ∀ x y z : 𝔤, κ ⁅x, y⁆ z = κ x ⁅y, z⁆
\end{leancode}

\begin{dictionary}
\lean{LinearMap.BilinForm ℂ 𝔤} is the type of $\C$-bilinear forms $\ltwog\times\ltwog\to\C$
(curried: \lean{κ x y}). The predicate has two fields, symmetry and ad-invariance, exactly as in
Definition~\ref{l02:def:invariant}. \emph{Nothing more is required}: $\kappa$ may be degenerate,
even zero, and $\ltwog$ need not be simple, semisimple or finite-dimensional. This generality
matters in Section~\ref{l02:sec:axiom}, because \lean{PBWPrinciple} quantifies over all such $\kappa$; we
shall check that it is the intended (and true) statement even for degenerate $\kappa$. The file's docstring explains why
symmetry is spelled out rather than taken from Mathlib (a sesquilinear notion of symmetry is
avoided). Examples: the Killing form $\kappa(x,y) = \tr(\mathrm{ad}x\,\mathrm{ad}y)$ of any
finite-dimensional Lie algebra; the trace form $\tr(xy)$ on $\mathfrak{gl}_n$; any symmetric form on an abelian
Lie algebra (where ad-invariance is vacuous); the zero form.
\end{dictionary}

\begin{physics}
$\kappa$ is the normalisation of the current--current two-point function,
$\langle J^a(z)J^b(w)\rangle = k\kappa(x_a,x_b)/(z-w)^2$; symmetry is Bose symmetry of the two
currents, and invariance is the statement that the two-point function is $G$-invariant. A
degenerate $\kappa$ corresponds to currents with vanishing two-point functions (null currents),
which are allowed algebraically though not in a unitary theory.
\end{physics}

\begin{status}
A definition; no axioms involved.
\end{status}

\subsection{The loop algebra and its cocycle}

\begin{definition}[Loop bracket, degree derivation, Kac--Moody cocycle]\label{l02:def:loop}
On $L\ltwog = \ltwog\otimes\C[t,t^{-1}]$ define
\begin{itemize}
\item the loop bracket $[x\otimes t^m, y\otimes t^n] = [x,y]\otimes t^{m+n}$;
\item the degree derivation $D(x\otimes t^n) = n\,x\otimes t^n$, i.e.\ $D = t\frac{d}{dt}$;
\item the cocycle $\omega(x\otimes t^m, y\otimes t^n) = m\,\kappa(x,y)\,\delta_{m+n,0}$.
\end{itemize}
\end{definition}

File \leanpath{VOA/Foundations/AffineLie.lean}.
\begin{leancode}
noncomputable def loopBracket : (ℤ →₀ 𝔤) →ₗ[ℂ] (ℤ →₀ 𝔤) →ₗ[ℂ] (ℤ →₀ 𝔤) :=
  Finsupp.lsum ℂ loopBracketLeft

@[simp] theorem loopBracket_single (m n : ℤ) (x y : 𝔤) :
    loopBracket (single m x) (single n y) = single (m + n) ⁅x, y⁆ := by

noncomputable def loopDeriv : (ℤ →₀ 𝔤) →ₗ[ℂ] (ℤ →₀ 𝔤) :=
  Finsupp.lsum ℂ fun n => (n : ℂ) • Finsupp.lsingle n

noncomputable def cocycle (κ : LinearMap.BilinForm ℂ 𝔤) :
    (ℤ →₀ 𝔤) →ₗ[ℂ] (ℤ →₀ 𝔤) →ₗ[ℂ] ℂ :=
  Finsupp.lsum ℂ (cocycleLeft κ)

@[simp] theorem cocycle_single (κ : LinearMap.BilinForm ℂ 𝔤) (m n : ℤ) (x y : 𝔤) :
    cocycle κ (single m x) (single n y) = if m + n = 0 then (m : ℂ) * κ x y else 0 := by
\end{leancode}
and the four identities from which the Jacobi identity of $\ltwoghat$ is derived:
\begin{leancode}
theorem loop_leibniz (f g h : ℤ →₀ 𝔤) :
    loopBracket f (loopBracket g h)
      = loopBracket (loopBracket f g) h + loopBracket g (loopBracket f h) := by

theorem deriv_loopBracket (f g : ℤ →₀ 𝔤) :
    loopDeriv (loopBracket f g)
      = loopBracket (loopDeriv f) g + loopBracket f (loopDeriv g) := by

theorem cocycle_cocycle (κ : LinearMap.BilinForm ℂ 𝔤) (hκ : κ.IsInvariant) (f g h : ℤ →₀ 𝔤) :
    cocycle κ f (loopBracket g h)
      = cocycle κ (loopBracket f g) h + cocycle κ g (loopBracket f h) := by

theorem cocycle_deriv_skew (κ : LinearMap.BilinForm ℂ 𝔤) (f g : ℤ →₀ 𝔤) :
    cocycle κ (loopDeriv f) g + cocycle κ f (loopDeriv g) = 0 := by
\end{leancode}

\begin{dictionary}
\textbf{Decoding.} An element $f$ of \lean{ℤ →₀ 𝔤} is the Laurent polynomial
$\sum_n f(n)\otimes t^n$; \lean{single m x} is $x\otimes t^m$. \lean{Finsupp.lsum} extends a family
of linear maps defined on the summands \lean{single n} to all of \lean{ℤ →₀ 𝔤}; the auxiliary
\lean{loopBracketLeft m x} sends $g\mapsto\sum_n[x,g(n)]\otimes t^{m+n}$, and \lean{cocycleLeft κ m x}
sends $g\mapsto m\,\kappa(x,g(-m))$. The simp lemmas \lean{loopBracket_}\allowbreak\lean{single} and
\lean{cocycle_single} confirm that on basis elements the definitions are those of
Definition~\ref{l02:def:loop}; since all three maps are (bi)linear, this determines them.
\lean{loopDeriv} sends \lean{single n x} to $n\cdot$\lean{single n x}, i.e.\ $D = t\,d/dt$.

\textbf{Hypotheses.} Only \lean{cocycle_cocycle} needs \lean{hκ}; \lean{cocycle_skew} (alternation
of $\omega$, not displayed) also needs it. The derivation compatibility
\lean{cocycle_}\allowbreak\lean{deriv_skew} holds for every bilinear $\kappa$: on basis vectors it reads
$(m+n)\,m\,\kappa(x,y)\,\delta_{m+n,0} = 0$.
\end{dictionary}

\begin{physics}
$L\ltwog$ is the Lie algebra of the loop group $LG$ of maps $S^1\to G$, the classical symmetry of
a chiral current algebra. $\omega$ is the Schwinger term of \eqref{l02:eq:affinemodes}: the
two-point function of the currents, read in modes: in $[J^a_m,J^b_n] = \oint_w\oint_{z}z^mw^n\,J^a(z)J^b(w)$
the double pole $k\kappa(x_a,x_b)/(z-w)^2$ contributes
$k\kappa(x_a,x_b)\oint_w\partial_w(w^m)\,w^n = k\,m\,\kappa(x_a,x_b)\,\delta_{m+n,0}$, which is where the factor $m$
comes from.
\end{physics}

\begin{proof}[Proof on paper of the 2-cocycle identity]
Take $f = w\otimes t^l$, $g = x\otimes t^m$, $h = y\otimes t^n$. All three terms vanish unless
$l+m+n=0$. Then the left side is $l\,\kappa(w,[x,y])$, and the right side is
$(l+m)\kappa([w,x],y) + m\,\kappa(x,[w,y])$. By invariance $\kappa([w,x],y) = \kappa(w,[x,y])$, and
by symmetry and invariance $\kappa(x,[w,y]) = \kappa([w,y],x) = \kappa(w,[y,x]) = -\kappa(w,[x,y])$. So
the right side is $(l+m-m)\kappa(w,[x,y])$, equal to the left side.
\end{proof}

\begin{status}
All four identities: \lean{[propext, Classical.choice, Quot.sound]} (they are ingredients of
\lean{affine_bracket}, whose report is quoted below).
\end{status}

\subsection{The affine algebra}

\begin{definition}[Affine Kac--Moody algebra]\label{l02:def:affine}
Let $\ltwog$ be a complex Lie algebra with invariant form $\kappa$. The \emph{affine algebra} is
$\ltwoghat = L\ltwog\oplus\C C\oplus\C\ell_0$ with bracket
\begin{equation}\label{l02:eq:bracketP}
 [(f,a,p),(g,b,q)] = \bigl([f,g] + p\,Dg - q\,Df,\ \omega(f,g),\ 0\bigr).
\end{equation}
In particular
\begin{equation*}
 [x\otimes t^m, y\otimes t^n] = [x,y]\otimes t^{m+n} + m\,\kappa(x,y)\,\delta_{m+n,0}\,C,\qquad
 [C,\cdot\,] = 0,\qquad [\ell_0,x\otimes t^n] = n\,x\otimes t^n.
\end{equation*}
A representation of $\ltwoghat$ has \emph{level} $k$ if $C$ acts as $k\cdot\mathrm{id}$.
\end{definition}

File \leanpath{VOA/Foundations/AffineLie.lean}.
\begin{leancode}
structure AffineLie (𝔤 : Type*) [LieRing 𝔤] [LieAlgebra ℂ 𝔤]
    (κ : LinearMap.BilinForm ℂ 𝔤) (hκ : κ.IsInvariant) where
  carrier : (ℤ →₀ 𝔤) × ℂ × ℂ

noncomputable def bracketP (a b : (ℤ →₀ 𝔤) × ℂ × ℂ) : (ℤ →₀ 𝔤) × ℂ × ℂ :=
  (loopBracket a.1 b.1 + a.2.2 • loopDeriv b.1 - b.2.2 • loopDeriv a.1, cocycle κ a.1 b.1, 0)

noncomputable instance : LieRing (AffineLie 𝔤 κ hκ) where
  add_lie X Y Z := by
    simp only [bracket_def, of_symm_add, bracketP_add_left, of_add]
  lie_add X Y Z := by
    simp only [bracket_def, of_symm_add, bracketP_add_right, of_add]
  lie_self X := by
    rw [bracket_def, bracketP_self κ hκ, of_zero]

noncomputable def ofLoop (x : 𝔤) (m : ℤ) : AffineLie 𝔤 κ hκ := ofProd κ hκ (single m x, 0, 0)

noncomputable def C : AffineLie 𝔤 κ hκ := ofProd κ hκ (0, 1, 0)

noncomputable def ell0 : AffineLie 𝔤 κ hκ := ofProd κ hκ (0, 0, 1)

theorem affine_bracket (x y : 𝔤) (m n : ℤ) :
    ⁅ofLoop κ hκ x m, ofLoop κ hκ y n⁆
      = ofLoop κ hκ ⁅x, y⁆ (m + n)
        + (if m + n = 0 then ((m : ℂ) * κ x y) • C κ hκ else 0) := by

@[simp] theorem C_lie (X : AffineLie 𝔤 κ hκ) : ⁅C κ hκ, X⁆ = 0 := by

theorem ell0_lie_ofLoop (x : 𝔤) (n : ℤ) :
    ⁅ell0 κ hκ, ofLoop κ hκ x n⁆ = (n : ℂ) • ofLoop κ hκ x n := by

structure AffineRep (V : Type*) [AddCommGroup V] [Module ℂ V] where
  ρ : AffineLie 𝔤 κ hκ →ₗ⁅ℂ⁆ Module.End ℂ V
  level : ℂ
  acts_central : ρ (C κ hκ) = level • LinearMap.id
\end{leancode}

\begin{dictionary}
\textbf{Decoding the carrier.} \lean{AffineLie 𝔤 κ hκ} is a one-field wrapper around the product
\lean{(ℤ →₀ 𝔤) × ℂ × ℂ}: a triple $(f,a,p)$ stands for $f + aC + p\ell_0$. The wrapper exists so
that the Lie structure, which depends on $\kappa$, is attached to a type that remembers $\kappa$;
\lean{ofProd} is the identification, and the vector space structure is transported along it. The
argument \lean{hκ} is a proof, and in Lean any two proofs of the same proposition are equal, so
the type does not depend on \emph{which} proof of invariance is supplied.

\textbf{Decoding the bracket.} In \lean{bracketP}, \lean{a.1} is the loop part, \lean{a.2.1} the
$C$-coefficient (it never enters: $C$ is central) and \lean{a.2.2} the $\ell_0$-coefficient. So
\lean{bracketP} is formula \eqref{l02:eq:bracketP} letter for letter. The result has $\ell_0$-component
$0$: $\ell_0$ is never produced by a bracket, as in the textbook. The \lean{LieRing} instance
(of which the first three fields are displayed; the fourth, \lean{leibniz_lie}, is the Jacobi
identity, proved from \lean{bracketP_leibniz}) shows that $\ltwoghat$ \emph{is} a Lie algebra: this is
not assumed. The proof of \lean{bracketP_leibniz} uses the four identities of the previous
subsection and \lean{hκ}; without symmetry and invariance of $\kappa$ the Jacobi identity would fail
and the instance could not have been written.

\textbf{The defining relations as theorems.} \lean{ofLoop κ hκ x m} is $x\otimes t^m$, \lean{C κ hκ}
is $C$, \lean{ell0 κ hκ} is $\ell_0$. \lean{affine_bracket} is the displayed relation of
Definition~\ref{l02:def:affine} (the Lean \lean{if m + n = 0 then … else 0} is
$\delta_{m+n,0}$), \lean{C_lie} says $C$ is central, \lean{ell0_lie_ofLoop} gives
$[\ell_0,x\otimes t^n] = n\,x\otimes t^n$.

\textbf{The sign of $\ell_0$.} The code implements $\ell_0 = +t\frac{d}{dt}$, as \lean{ell0_lie_ofLoop}
certifies: $[t\frac{d}{dt}, x\otimes t^n] = n\,x\otimes t^n$. The physically natural energy operator is $-\ell_0$: on the vacuum module
the creation operator $x\otimes t^{-n}$ raises $-\ell_0$ by $n$ (Section~\ref{l02:sec:grading}).

\textbf{\lean{AffineRep}.} A representation of level $k$: a Lie homomorphism $\rho$ to
$\End_\C V$, a number \lean{level}, and the proof that $\rho(C) = k\cdot\mathrm{id}$.

\textbf{Same as the textbook?} Yes. Kac's book defines $\ltwoghat$ for simple $\ltwog$ with the
normalised invariant form; the library's definition is the general one (any $\ltwog$, any invariant
$\kappa$), which specialises to Kac's. For $\ltwog$ abelian of dimension $r$ and $\kappa$ an inner
product it is the rank-$r$ Heisenberg algebra (extended by $\ell_0$).
\end{dictionary}

\begin{physics}
$\ltwoghat$ is the algebra of the modes $J^a_n$ of the conserved currents of a WZW model or of any
CFT with a continuous symmetry (current algebra); $C$ is the level $k$, which in the WZW model is
the integer coefficient of the Wess--Zumino term; $\ell_0$ is (minus) the part of the Hamiltonian
that counts the mode number. The cocycle is the Schwinger term, i.e.\ the chiral anomaly of the
currents.
\end{physics}

\begin{proof}[Proof on paper of the Jacobi identity]
Write $X = (f,a,p)$ etc. The $\ell_0$-component of any bracket is $0$, so in a double bracket
$[X,[Y,Z]]$ the inner bracket has no $\ell_0$-part. The loop component of
$[X,[Y,Z]] - [[X,Y],Z] - [Y,[X,Z]]$ is a combination of: the Jacobi identity of the loop bracket
(which follows from that of $\ltwog$, termwise), terms $p(D[g,h] - [Dg,h]-[g,Dh])$ and its analogues,
which vanish because $D$ is a derivation (on basis vectors: $(m+n) = m + n$), and terms with two
$D$'s which cancel by skew-symmetry. The $C$-component is
$\omega(f,[g,h]+qDh-rDg) - \omega([f,g]+pDg-qDf,h) - \omega(g,[f,h]+pDh-rDf)$; the terms without $D$
cancel by the 2-cocycle identity, and the terms with $D$ cancel in pairs by
$\omega(Du,v) + \omega(u,Dv) = 0$ and skew-symmetry of $\omega$.
\end{proof}

\begin{status}
\lean{VOA.}\allowbreak\lean{AffineLie.}\allowbreak\lean{affine_}\allowbreak\lean{bracket}, \lean{VOA.Affine}\allowbreak\lean{Lie.C_lie},
\lean{VOA.}\allowbreak\lean{AffineLie.}\allowbreak\lean{ell0_}\allowbreak\lean{lie_}\allowbreak\lean{ofLoop}: \lean{[propext, Classical.choice, Quot.sound]} (run for this
chapter). Unconditional; the Lie algebra structure is proved, not assumed.
\end{status}

\subsection{The triangular pieces}

File \leanpath{VOA/Foundations/AffineLie.lean}.
\begin{leancode}
noncomputable def posPart : LieSubalgebra ℂ (AffineLie 𝔤 κ hκ) :=
  LieSubalgebra.lieSpan ℂ _ {X | ∃ (x : 𝔤) (n : ℤ), 0 < n ∧ X = ofLoop κ hκ x n}

noncomputable def negPart : LieSubalgebra ℂ (AffineLie 𝔤 κ hκ) :=
  LieSubalgebra.lieSpan ℂ _ {X | ∃ (x : 𝔤) (n : ℤ), n < 0 ∧ X = ofLoop κ hκ x n}

noncomputable def parabolic : LieSubalgebra ℂ (AffineLie 𝔤 κ hκ) :=
  LieSubalgebra.lieSpan ℂ _
    ({X | ∃ (x : 𝔤) (n : ℤ), 0 ≤ n ∧ X = ofLoop κ hκ x n} ∪ {C κ hκ, ell0 κ hκ})
\end{leancode}

\begin{dictionary}
These are the Lie subalgebras \emph{generated} by the indicated elements: $\ltwon_+$ by the
$x\otimes t^n$ with $n>0$, $\ltwon_-$ by those with $n<0$, and the parabolic $\hat{\ltwop}$ by the
$x\otimes t^n$ with $n\ge0$ together with $C$ and $\ell_0$. Since the loop bracket adds degrees and
the cocycle is supported on $m+n = 0$, the generated subalgebras are the expected subspaces:
$\ltwon_\pm = \ltwog\otimes t^{\pm1}\C[t^{\pm1}]$ (for $m,n$ of the same strict sign, $m+n\ne0$, so no $C$
appears) and $\hat{\ltwop} = \ltwog[t]\oplus\C C\oplus\C\ell_0$ (for $m,n\ge0$, $m+n = 0$ forces $m = 0$,
and the cocycle carries the factor $m$). The library does not prove these equalities of
subspaces; they are not needed, because the vacuum module below is defined directly by
generators.
\end{dictionary}

\begin{physics}
$\ltwon_-$ consists of the creation operators $J^a_{-n}$, $n\ge1$; $\hat{\ltwop}$ of the operators that
annihilate the vacuum or act on it by scalars (zero modes $J^a_0$, the level, the energy).
\end{physics}

\begin{status}
Definitions; no axioms involved.
\end{status}

%======================================================================
\section{The affine vacuum module}\label{l02:sec:vacuum}
%======================================================================

\begin{definition}[Affine vacuum module]\label{l02:def:vacuum}
Let $U(\ltwoghat)$ be the universal enveloping algebra and $J(k)\subseteq U(\ltwoghat)$ the left ideal
generated by
\begin{equation*}
 \{x\otimes t^n : x\in\ltwog,\ n\ge0\}\ \cup\ \{C - k\}\ \cup\ \{\ell_0\}.
\end{equation*}
The \emph{level-$k$ vacuum module} is $\ltwoV(\ltwog,k) = U(\ltwoghat)/J(k)$ (often written
$V^k(\ltwog)$), with $\ltwoghat$ acting by left multiplication; its vacuum vector is $\ltwovac = [1]$.
\end{definition}

File \leanpath{VOA/Foundations/AffineLie.lean}.
\begin{leancode}
abbrev UEnv (𝔤 : Type*) [LieRing 𝔤] [LieAlgebra ℂ 𝔤]
    (κ : LinearMap.BilinForm ℂ 𝔤) (hκ : κ.IsInvariant) : Type _ :=
  UniversalEnvelopingAlgebra ℂ (AffineLie 𝔤 κ hκ)

def vacuumGens (k : ℂ) : Set (UEnv 𝔤 κ hκ) :=
  {a | (∃ (x : 𝔤) (n : ℤ), 0 ≤ n ∧ a = ι ℂ (ofLoop κ hκ x n))
        ∨ a = ι ℂ (C κ hκ) - algebraMap ℂ (UEnv 𝔤 κ hκ) k
        ∨ a = ι ℂ (ell0 κ hκ)}

noncomputable def vacuumIdeal (k : ℂ) : Submodule (UEnv 𝔤 κ hκ) (UEnv 𝔤 κ hκ) :=
  Submodule.span _ (vacuumGens κ hκ k)

abbrev vacuumSpace (k : ℂ) : Type _ := UEnv 𝔤 κ hκ ⧸ vacuumIdeal κ hκ k

noncomputable def vacuumRho (k : ℂ) :
    AffineLie 𝔤 κ hκ →ₗ⁅ℂ⁆ Module.End ℂ (vacuumSpace κ hκ k) :=
  (Algebra.lsmul ℂ ℂ (vacuumSpace κ hκ k)).toLieHom.comp (ι ℂ)

noncomputable def vacuumVector (k : ℂ) : vacuumSpace κ hκ k := Submodule.Quotient.mk 1
\end{leancode}
and its basic properties:
\begin{leancode}
theorem vacuum_acts_central (k : ℂ) :
    vacuumRho κ hκ k (C κ hκ) = k • LinearMap.id := by

theorem vacuum_pos_annihilate (k : ℂ) (x : 𝔤) (n : ℤ) (hn : 0 ≤ n) :
    vacuumRho κ hκ k (ofLoop κ hκ x n) (vacuumVector κ hκ k) = 0 := by

theorem vacuum_ell0_annihilate (k : ℂ) :
    vacuumRho κ hκ k (ell0 κ hκ) (vacuumVector κ hκ k) = 0 := by

theorem vacuum_cyclic (k : ℂ) :
    Submodule.span (UEnv 𝔤 κ hκ) {vacuumVector κ hκ k} = ⊤ := by
\end{leancode}

\begin{dictionary}
\textbf{Decoding.} \lean{UEnv 𝔤 κ hκ} is $U(\ltwoghat)$. \lean{vacuumGens κ hκ k} is the set of three
kinds of elements of Definition~\ref{l02:def:vacuum}: \lean{ι ℂ (ofLoop κ hκ x n)} for $n\ge0$ (note
\lean{0 ≤ n}: the zero modes $x\otimes t^0$ are included), $\iota(C) - k\cdot1$, and $\iota(\ell_0)$.
\lean{vacuumIdeal} is the \emph{left} ideal they generate (\lean{Submodule.span} over $U(\ltwoghat)$
acting on itself by left multiplication), \lean{vacuumSpace} the quotient module, \lean{vacuumRho}
the action by left multiplication, \lean{vacuumVector} the class of $1$. The theorems say: $C$ acts
by $k$ on the whole module; the modes $x\otimes t^n$, $n\ge0$, and $\ell_0$ kill $\ltwovac$; and
$\ltwovac$ generates.

\textbf{This is the standard induced module.} In textbooks $V^k(\ltwog)$ is defined as
$U(\ltwoghat)\otimes_{U(\hat{\ltwop})}\C_k$, where $\C_k$ is the one-dimensional $\hat{\ltwop}$-module on
which $\ltwog[t]$ and $\ell_0$ act by $0$ and $C$ by $k$. The quotient $U(\ltwoghat)/J(k)$ is the same
module: in general, for a Lie subalgebra $\ltwop$ and a character $\chi:\ltwop\to\C$, one has
$U(\ltwog)\otimes_{U(\ltwop)}\C_\chi\cong U(\ltwog)/U(\ltwog)\{a-\chi(a): a\in\ltwop\}$, and the
generators of $J(k)$ are exactly the elements $a-\chi(a)$ for $a$ running over the spanning set
$\{x\otimes t^n\}_{n\ge0}\cup\{C,\ell_0\}$ of $\hat{\ltwop}$ (with $\chi(x\otimes t^n) = 0$,
$\chi(C) = k$, $\chi(\ell_0) = 0$). Section~\ref{l02:sec:axiom} uses this identification.

\textbf{Some texts omit $\ell_0$} and induce from $\ltwog[t]\oplus\C C$ in the algebra without
derivation; the resulting space is the same, and including $\ell_0$ with eigenvalue $0$ on $\ltwovac$
is harmless (and it is what makes the energy operator available inside the module).
(Chapters~\ref{ch:01} and~\ref{ch:15} describe the ideal as ``generated by $\ltwog[t]$ and $C-k$'';
the precise Lean generating set also contains $\ell_0$.)

\textbf{Vacuity.} As for Virasoro, every theorem in this block would hold for the zero module.
Whether $\ltwovac\ne0$ is \emph{the} question; it is the first field of \lean{PBWPrinciple}
(Section~\ref{l02:sec:axiom}).
\end{dictionary}

\begin{physics}
$V^k(\ltwog)$ is the space of states $J^{a_1}_{-n_1}\cdots J^{a_r}_{-n_r}\ltwovac$ of the vacuum
sector of the current algebra at level $k$ --- the WZW vacuum module before null vectors are removed
(for $k\in\N$ and $\ltwog$ simple, the WZW model's actual vacuum module is the simple quotient
$L_k(\ltwog)$). The vacuum conditions $J^a_n\ltwovac = 0$, $n\ge0$, say that $J^a(z)\ltwovac$ is regular
at $z=0$ and that the vacuum is $G$-invariant ($J^a_0\ltwovac = 0$).
\end{physics}

\begin{proof}[Proof on paper of \lean{vacuum_acts_}\allowbreak\lean{central}]
For $a\in U(\ltwoghat)$, $C[a] - k[a] = [(C-k)a] = [a(C-k)]$ because $C$ is central in $\ltwoghat$,
hence $\iota(C)$ is central in $U(\ltwoghat)$ (\lean{commute_ι_C}); and $a(C-k)\in J(k)$ since $J(k)$
is a left ideal containing $C-k$. The other three statements are immediate from the definitions:
$x\otimes t^n\cdot[1] = [x\otimes t^n]$ and $x\otimes t^n\in J(k)$; and $[a] = a\cdot[1]$.
\end{proof}

\begin{status}
\lean{VOA.}\allowbreak\lean{AffineLie.}\allowbreak\lean{vacuum_}\allowbreak\lean{acts_}\allowbreak\lean{central}, \lean{VOA.}\allowbreak\lean{AffineLie.}\allowbreak\lean{vacuum_}\allowbreak\lean{pos_}\allowbreak\lean{annihilate},
\lean{VOA.}\allowbreak\lean{AffineLie.}\allowbreak\lean{vacuum_}\allowbreak\lean{cyclic}: \lean{[propext, Classical.choice, Quot.sound]} (run for this
chapter). Unconditional, and \lean{VOA.pbw}-free.
\end{status}

%======================================================================
\section{Spanning and the energy grading}\label{l02:sec:grading}
%======================================================================

\begin{theorem}[Spanning; the easy half of PBW]\label{l02:thm:span}
$V^k(\ltwog)$ is spanned by the creation monomials
\begin{equation*}
 (x_1\otimes t^{-n_1-1})(x_2\otimes t^{-n_2-1})\cdots(x_r\otimes t^{-n_r-1})\ltwovac,\qquad r\ge0,\ x_i\in\ltwog,\ n_i\in\N.
\end{equation*}
Each such monomial is an $\ell_0$-eigenvector with eigenvalue $-\sum_i(n_i+1)$. Consequently
$V^k(\ltwog)$ is the direct sum of its $\ell_0$-eigenspaces; the only eigenvalues that occur are
$-d$, $d\in\N$; the eigenvalue-$0$ space is $\C\ltwovac$; and if $\ltwog$ is finite-dimensional, every
eigenspace is finite-dimensional.
\end{theorem}

File \leanpath{VOA/Foundations/AffineLie.lean}.
\begin{leancode}
noncomputable def negWord (k : ℂ) : List (𝔤 × ℕ) → vacuumSpace κ hκ k
  | [] => vacuumVector κ hκ k
  | (x, n) :: L => vacuumRho κ hκ k (ofLoop κ hκ x (-(n : ℤ) - 1)) (negWord k L)

noncomputable def negSpan (k : ℂ) : Submodule ℂ (vacuumSpace κ hκ k) :=
  Submodule.span ℂ (Set.range (negWord κ hκ k))

theorem negSpan_eq_top (k : ℂ) : negSpan κ hκ k = ⊤ := by

def wordDeg : List (𝔤 × ℕ) → ℕ
  | [] => 0
  | (_, n) :: L => (n + 1) + wordDeg L

theorem vacuumRho_ell0_negWord_eq (k : ℂ) (L : List (𝔤 × ℕ)) :
    vacuumRho κ hκ k (ell0 κ hκ) (negWord κ hκ k L)
      = (-(wordDeg L : ℂ)) • negWord κ hκ k L := by

noncomputable def ell0Op (k : ℂ) : Module.End ℂ (vacuumSpace κ hκ k) :=
  vacuumRho κ hκ k (ell0 κ hκ)

theorem ell0_isInternal (k : ℂ) :
    DirectSum.IsInternal (Module.End.eigenspace (ell0Op κ hκ k)) :=

theorem eigenspace_eq_bot_of_not_energy (k : ℂ) (μ : ℂ) (hμ : ∀ d : ℕ, μ ≠ -(d : ℂ)) :
    Module.End.eigenspace (ell0Op κ hκ k) μ = ⊥ := by

theorem eigenspace_zero_eq (k : ℂ) :
    Module.End.eigenspace (ell0Op κ hκ k) 0 = Submodule.span ℂ {vacuumVector κ hκ k} := by

theorem eigenspace_finiteDimensional [FiniteDimensional ℂ 𝔤] (k : ℂ) (μ : ℂ) :
    FiniteDimensional ℂ (Module.End.eigenspace (ell0Op κ hκ k) μ) := by
\end{leancode}

\begin{dictionary}
\textbf{Decoding the index convention.} A word is a list of pairs $(x,n)$ with $n\in\N$, and the
pair $(x,n)$ stands for the mode $x\otimes t^{-n-1}$: \emph{depth} $n$ means mode $-n-1$. The first
element of the list is the \emph{leftmost} operator. \lean{wordDeg} is the energy
$\sum(n_i+1)$. \lean{negSpan} is the $\C$-span of all words --- \emph{all} words, in any order and with
repetitions, not only ordered ones.

\textbf{Decoding the conclusions.} \lean{Module.End.eigenspace T μ} is $\ker(T-\mu)$.
\lean{DirectSum.IsInternal (fun μ => E μ)} says the canonical map $\bigoplus_\mu E_\mu\to V$ is
bijective, i.e.\ $V = \bigoplus_\mu E_\mu$. \lean{⊥} is the zero subspace. The four grading
theorems are the four sentences after ``Consequently'' in Theorem~\ref{l02:thm:span}; only the last
needs \lean{[FiniteDimensional ℂ 𝔤]}.

\textbf{Sign.} $\ell_0$ has eigenvalue $-(\text{energy})$, because $\ell_0 = t\,d/dt$ and creation
operators have negative mode. The library's affine VOA (Chapter~\ref{ch:L08}) uses the grading by
$-\ell_0$. Its carrier for $\mathfrak{sl}_2$ is exactly this universal vacuum module, under the Lean name
\lean{AffineVOA.WakModule} (an abbreviation for \lean{vacuumSpace} in \leanpath{VOA/\allowbreak Examples/\allowbreak AffineSl2.lean}); despite the name it is the vacuum module
$V^k(\mathfrak{sl}_2)$, \emph{not} a Wakimoto free-field module (release audit 2026-10-27, finding m-6). The VOA built on it is
of CFT type (Chapter~\ref{ch:L03}): $\deg 0$ is $\C\ltwovac$ and there are no negative degrees, as this section shows for
the energy grading.

\textbf{What is \emph{not} claimed.} Linear independence of anything. In particular
\lean{eigenspace_}\allowbreak\lean{zero_eq} says the degree-$0$ piece is \emph{spanned} by $\ltwovac$; it is
one-dimensional only if $\ltwovac\ne0$, which these theorems do not show (the module docstring says so
explicitly). This is why nothing in this section needs \lean{VOA.pbw}; when the PBW principle was
an axiom, this was what kept the energy grading free of it.
\end{dictionary}

\begin{physics}
The energy grading is the $L_0$-grading of the current algebra's vacuum sector (for the Sugawara
stress tensor $L_0$ agrees with $-\ell_0$ on $V^k(\ltwog)$ at noncritical level). Finite-dimensionality
of the graded pieces is what makes the partition function $\tr\,q^{L_0}$ a formal power series with
finite coefficients; boundedness below says there is a ground state.
\end{physics}

\begin{proof}[Proof on paper]
\emph{Spanning.} Let $S$ be the span of all creation monomials. It contains $\ltwovac$. It is stable
under every $x\otimes t^n$: by induction on the length of the word, write
$(x\otimes t^n)(y\otimes t^{-m-1})w = (y\otimes t^{-m-1})(x\otimes t^n)w + [x\otimes t^n, y\otimes t^{-m-1}]w$;
the first term is in $S$ by induction and because $S$ is closed under prepending creation operators,
the second is $([x,y]\otimes t^{n-m-1})w + \delta\cdot n\,\kappa(x,y)\,k\,w$, in $S$ by induction. At length
$0$: $x\otimes t^n\ltwovac$ is $0$ if $n\ge0$ and a monomial if $n<0$. It is stable under $C$ (scalar
$k$) and $\ell_0$ (next paragraph). Since these generate $\ltwoghat$ as a vector space, $S$ is a
$U(\ltwoghat)$-submodule containing the cyclic vector, so $S = V^k(\ltwog)$.

\emph{Eigenvalues.} $\ell_0(x\otimes t^{-n-1})w = (x\otimes t^{-n-1})\ell_0w - (n+1)(x\otimes t^{-n-1})w$;
induction on length gives eigenvalue $-\sum(n_i+1)$.

\emph{Grading.} Eigenvectors with distinct eigenvalues are independent (always), and they span by the
previous two points; hence the direct sum. An eigenvector with eigenvalue $\mu$ is, by independence,
a combination of monomials of energy $-\mu$ only; so the $\mu$-eigenspace is the span of the monomials
of energy $-\mu$ (\lean{eigenspace_}\allowbreak\lean{eq_degSpan}). This span is $0$ if $-\mu\notin\N$, and is $\C\ltwovac$ if
$\mu=0$ (only the empty word has energy $0$). For finite-dimensional $\ltwog$, the energy-$d$ span is
contained in $\C\ltwovac + \sum_{n<d}(\ltwog\otimes t^{-n-1})\cdot(\text{energy }d-n-1\text{ span})$, a finite
sum of images of bilinear maps on finite-dimensional spaces; strong induction on $d$ gives finite
dimension.
\end{proof}

\begin{status}
\lean{VOA.}\allowbreak\lean{AffineLie.}\allowbreak\lean{negSpan_}\allowbreak\lean{eq_}\allowbreak\lean{top}, \lean{VOA.}\allowbreak\lean{AffineLie.}\allowbreak\lean{ell0_}\allowbreak\lean{isInternal},
\lean{VOA.}\allowbreak\lean{AffineLie.}\allowbreak\lean{eigenspace_}\allowbreak\lean{eq_}\allowbreak\lean{bot_}\allowbreak\lean{of_}\allowbreak\lean{not_}\allowbreak\lean{energy}, \lean{VOA.}\allowbreak\lean{AffineLie.}\allowbreak\lean{eigenspace_}\allowbreak\lean{zero_}\allowbreak\lean{eq},
\lean{VOA.}\allowbreak\lean{AffineLie.}\allowbreak\lean{eigenspace_}\allowbreak\lean{finiteDimensional}: \lean{[propext, Classical.choice, Quot.sound]}
(run for this chapter). Unconditional; does not use \lean{VOA.pbw}.
\end{status}

%======================================================================
\section{The PBW principle \texorpdfstring{\texttt{VOA.pbw}}{VOA.pbw}: statement and proof}\label{l02:sec:axiom}
%======================================================================

\subsection{The classical theorem}

\begin{theorem}[Poincar\'e--Birkhoff--Witt]\label{l02:thm:pbw}
Let $\ltwog$ be a Lie algebra over a field $F$ (of any dimension), and let $(e_i)_{i\in I}$ be a
basis of $\ltwog$ indexed by a totally ordered set $I$. Then the ordered monomials
\begin{equation*}
 \iota(e_{i_1})\iota(e_{i_2})\cdots\iota(e_{i_r}),\qquad r\ge0,\quad i_1\le i_2\le\dots\le i_r,
\end{equation*}
form a basis of $U(\ltwog)$ (the empty monomial being $1$). Equivalently, the associated graded
algebra of $U(\ltwog)$ for its natural filtration is the symmetric algebra $S(\ltwog)$.
\end{theorem}

\begin{physics}
In a quantum theory whose operators are generated by modes $X_i$ with commutation relations
$[X_i,X_j] = \sum f_{ij}{}^kX_k$, PBW says that every operator product can be normal-ordered in a
fixed order, uniquely: the normal-ordered monomials are a basis. The passage $U(\ltwog)\leadsto S(\ltwog)$
(associated graded) is the classical limit, in which operators commute and become functions on $\ltwog^*$.
\end{physics}

Every Lie algebra over a field has a basis (Zorn's lemma), and every index set can be totally
ordered (by the well-ordering theorem), so the theorem applies to every Lie algebra. Two
consequences are what the library needs.

\begin{corollary}\label{l02:cor:free}
Let $\ltwop\subseteq\ltwog$ be a Lie subalgebra.
\begin{enumerate}
\item The algebra map $U(\ltwop)\to U(\ltwog)$ induced by the inclusion is injective; identify
  $U(\ltwop)$ with its image.
\item Choose a basis $(e_i)_{i\in I_1}$ of a vector-space complement of $\ltwop$ and a basis
  $(e_i)_{i\in I_2}$ of $\ltwop$, and order $I = I_1\sqcup I_2$ so that $I_1 < I_2$. Then $U(\ltwog)$ is a
  free right $U(\ltwop)$-module with basis the ordered monomials in the $e_i$, $i\in I_1$:
  \begin{equation*}
   U(\ltwog) = \bigoplus_{m} m\cdot U(\ltwop),\qquad m = \iota(e_{i_1})\cdots\iota(e_{i_r}),\ i_1\le\dots\le i_r\in I_1.
  \end{equation*}
\item If moreover the complement $\ltwon$ is itself a Lie subalgebra, multiplication
  $U(\ltwon)\otimes U(\ltwop)\to U(\ltwog)$ is a linear isomorphism.
\end{enumerate}
\end{corollary}

\begin{proof}[Proof on paper]
Apply Theorem~\ref{l02:thm:pbw} to $\ltwog$ with the ordered basis of (2): each ordered monomial
factors uniquely as (ordered monomial in $I_1$)$\cdot$(ordered monomial in $I_2$), and the
ordered monomials in $I_2$ are, by PBW for $\ltwop$, a basis of $U(\ltwop)$ which maps to linearly
independent elements of $U(\ltwog)$; this gives (1) and (2). For (3), the ordered monomials in $I_1$
are a basis of $U(\ltwon)$ by PBW for $\ltwon$, so (2) says exactly that multiplication is
bijective.
\end{proof}

\begin{proof}[Sketch of the proof of Theorem~\ref{l02:thm:pbw}]
The spanning half is elementary: in $U(\ltwog)$, $\iota(e_j)\iota(e_i) = \iota(e_i)\iota(e_j) +
\iota([e_j,e_i])$, so any monomial can be rewritten, by induction on length and on the number of
inversions, as a combination of ordered monomials of the same or smaller length. The hard half is
independence. The classical proof (Humphreys \S17.4, Dixmier 2.1.11) constructs a representation of
$\ltwog$ on the polynomial algebra $F[z_i : i\in I]$ such that $e_i\cdot z_{i_1}\cdots z_{i_r} =
z_iz_{i_1}\cdots z_{i_r}$ whenever $i\le i_1\le\dots\le i_r$, by induction on the degree; the
compatibility with the bracket is the only computation, and uses the Jacobi identity. By the
universal property this representation extends to $U(\ltwog)$, and applying an ordered monomial to
$1$ gives the corresponding monomial $z_{i_1}\cdots z_{i_r}$; distinct ordered monomials thus have
linearly independent images, so they are independent. An alternative proof uses Bergman's diamond
lemma: the rewriting rules $e_je_i\to e_ie_j + [e_j,e_i]$ ($j>i$) have only resolvable ambiguities,
precisely by the Jacobi identity. No finiteness of $\dim\ltwog$ is used in either proof.
\end{proof}

\begin{physics}
PBW is the statement that the Fock-type module obtained by acting with creation operators on a
vacuum has the ``obvious'' basis of normal-ordered monomials, i.e.\ that the canonical commutation
relations generate \emph{all} the relations among products of generators, and there are no hidden
ones. For a free boson this is the familiar fact that $a_{-n_1}\cdots a_{-n_r}\ltwovac$ with
$n_1\ge\dots\ge n_r\ge1$ are independent; for a current algebra it is the corresponding statement for
$J^{a}_{-n}$, which is what makes the character of the universal vacuum module
$\prod_n(1-q^n)^{-\dim\ltwog}$.
\end{physics}

\subsection{Induced modules and the bottom map}\label{l02:sec:induced}

\paragraph{Notation.} In this section $U(\cdot)$, always written with its argument in parentheses, is the
universal enveloping algebra; a plain $U$ is the vector space of a representation $\rho$ of a Lie
subalgebra. (The Lean files use the same letters, which is why we keep them.) Later in the chapter $M$
denotes the space of Theorem~\ref{l02:thm:relrep} and $N$ a maximal submodule; $C$ is the central element
of $\ltwoghat$, $\mathfrak c$ a complement of $\ltwop$; Mathlib's canonical map $\ltwog\to U(\ltwog)$ is $\iota$,
while \lean{ι} in Section~\ref{l02:sec:pbwbasis} is an index type; and the Witt generator $\ell_0 = -z\,d/dz$
of Section~\ref{l02:sec:witt} is a different object from the affine derivation $\ell_0 = +t\,d/dt$.

\paragraph{The induced module.} Let $\ltwop\subseteq\ltwog$ be a Lie subalgebra and $\rho$ a representation of
$\ltwop$ on $U$. The \emph{induced module} is
\begin{equation*}
 \ltwoInd_{\ltwop}^{\ltwog}U = U(\ltwog)\otimes_{U(\ltwop)}U,
\end{equation*}
with $\ltwog$ acting by left multiplication on the first factor. It is the ``freest'' $\ltwog$-module
generated by a copy of $U$ on which $\ltwop$ acts through $\rho$: every $\ltwop$-equivariant map from $U$
to a $\ltwog$-module extends uniquely to a $\ltwog$-map out of it. The \emph{bottom map} is
$u\mapsto1\otimes u$; its image is the ``bottom'' layer, the copy of $U$ from which the rest of the module is
generated (in the vertex-algebra application, the lowest-weight space of a generalised Verma module;
in the affine vacuum module, the line $\C\ltwovac$).

\paragraph{Why injectivity is not automatic.} The tensor product over the noncommutative ring
$U(\ltwop)$ imposes the relations $y\,\iota(x)\otimes u = y\otimes\rho(x)u$, and these could in principle
identify distinct vectors of $U$ or kill them. Item~1 of the warning in
Section~\ref{l02:sec:axiom} shows that this really happens if $\rho$ is only linear and not a
representation. That the bottom map is injective for every representation is a consequence of PBW
(Theorem~\ref{l02:thm:B}); the nonvanishing of the affine vacuum is the special case $U = \C$
(Remark~\ref{l02:rem:AfromB}).

\paragraph{Which proof the library contains.} This section gives two proofs of the two fields of
\lean{PBWPrinciple}. The classical one (Theorems~\ref{l02:thm:A} and~\ref{l02:thm:B}, via the full PBW basis of
Corollary~\ref{l02:cor:free}) is a paper argument, used here only to confirm that the Lean statement is
true and means what it should. The proof \emph{formalised} in the library is the relative one of
Section~\ref{l02:sec:pbwproof}.

\subsection{The Lean statement}

The library states the two consequences it needs. First, induction from a Lie subalgebra, in a
namespace \lean{PBWInduction}:

File \leanpath{VOA/Foundations/AffineLie.lean}.
\begin{leancode}
universe u

namespace PBWInduction

variable {g U : Type u} [LieRing g] [LieAlgebra ℂ g]
variable [AddCommGroup U] [Module ℂ U]

abbrev Cover (g U : Type u) [LieRing g] [LieAlgebra ℂ g]
    [AddCommGroup U] [Module ℂ U] :=
  UniversalEnvelopingAlgebra ℂ g ⊗[ℂ] U

def balancingSet (p : LieSubalgebra ℂ g) (rho : p →ₗ⁅ℂ⁆ Module.End ℂ U) :
    Set (Cover g U) :=
  {t | ∃ (y : UniversalEnvelopingAlgebra ℂ g) (x : p) (u : U),
    t = (y * UniversalEnvelopingAlgebra.ι ℂ (x : g)) ⊗ₜ[ℂ] u - y ⊗ₜ[ℂ] (rho x u)}

noncomputable def balancing (p : LieSubalgebra ℂ g)
    (rho : p →ₗ⁅ℂ⁆ Module.End ℂ U) : Submodule ℂ (Cover g U) :=
  Submodule.span ℂ (balancingSet p rho)

abbrev Induced (p : LieSubalgebra ℂ g) (rho : p →ₗ⁅ℂ⁆ Module.End ℂ U) : Type u :=
  Cover g U ⧸ balancing p rho

noncomputable def bottom (p : LieSubalgebra ℂ g)
    (rho : p →ₗ⁅ℂ⁆ Module.End ℂ U) : U →ₗ[ℂ] Induced p rho :=
  (balancing p rho).mkQ.comp
    (TensorProduct.mk ℂ (UniversalEnvelopingAlgebra ℂ g) U 1)

end PBWInduction
\end{leancode}
Then the vacuum statement, the bundle, and the theorem (whose one-line proof term is shown because it
is informative: it names the two results, proved in Section~\ref{l02:sec:pbwproof}, from which the
two fields come):
\begin{leancode}
def PBW {𝔤 : Type*} [LieRing 𝔤] [LieAlgebra ℂ 𝔤]
    (κ : LinearMap.BilinForm ℂ 𝔤) (hκ : κ.IsInvariant) (k : ℂ) : Prop :=
  AffineLie.vacuumVector κ hκ k ≠ 0

structure PBWPrinciple : Prop where
  affineVacuum : ∀ {𝔤 : Type u} [LieRing 𝔤] [LieAlgebra ℂ 𝔤]
      (κ : LinearMap.BilinForm ℂ 𝔤) (hκ : κ.IsInvariant) (k : ℂ),
    PBW κ hκ k
  inducedBottom : ∀ {g U : Type u} [LieRing g] [LieAlgebra ℂ g]
      [AddCommGroup U] [Module ℂ U] (p : LieSubalgebra ℂ g)
      (rho : p →ₗ⁅ℂ⁆ Module.End ℂ U),
    Function.Injective (PBWInduction.bottom p rho)

theorem pbw : PBWPrinciple :=
  ⟨fun κ hκ k => AffineLie.vacuumVector_ne_zero_of_induction κ hκ k,
    fun p rho => PBWTheorem.bottom_injective p rho⟩
\end{leancode}

\begin{dictionary}
\textbf{Line-by-line decoding of \lean{PBWInduction}.}
\begin{itemize}
\item \lean{Cover g U} is the $\C$-vector space $U(\ltwog)\otimes_\C U$ (\lean{⊗[ℂ]} is Mathlib's
  tensor product of $\C$-modules; $U(\ltwog)$ is a $\C$-module through its $\C$-algebra structure).
\item \lean{balancingSet p rho} is the set of all elements $y\,\iota(x)\otimes u - y\otimes\rho(x)u$
  with $y\in U(\ltwog)$, $x\in\ltwop$, $u\in U$. Here \lean{(x : g)} is $x$ viewed in $\ltwog$,
  \lean{UniversalEnvelopingAlgebra.ι ℂ (x : g)} is its image in $U(\ltwog)$, \lean{y * ι x} is the
  product in $U(\ltwog)$ with $\iota(x)$ on the \emph{right}, and \lean{a ⊗ₜ[ℂ] u} is the pure
  tensor $a\otimes u$.
\item \lean{balancing p rho} is the $\C$-linear span of that set (not a $U(\ltwog)$-submodule; but it is
  automatically stable under left multiplication by $U(\ltwog)$, since the set is: multiply $y$ on the
  left).
\item \lean{Induced p rho} is the quotient vector space $(U(\ltwog)\otimes U)/\text{balancing}$; by
  Proposition~\ref{l02:prop:induced} below it is $U(\ltwog)\otimes_{U(\ltwop)}U$, the induced module
  $\ltwoInd_{\ltwop}^{\ltwog} U$.
\item \lean{bottom p rho} is the composite $U\to U(\ltwog)\otimes U\to\text{Induced}$,
  $u\mapsto 1\otimes u\mapsto[1\otimes u]$ (\lean{TensorProduct.mk ℂ A U 1} is $u\mapsto 1\otimes u$,
  \lean{mkQ} the quotient map).
\end{itemize}
\textbf{Decoding \lean{PBW} and \lean{PBWPrinciple}.} \lean{PBW κ hκ k} is the proposition
$\ltwovac\ne0$ in $V^k(\ltwog) = U(\ltwoghat)/J(k)$ of Definition~\ref{l02:def:vacuum}.
\lean{PBWPrinciple} is a structure in \lean{Prop} --- a conjunction of two propositions with named
components:
\begin{itemize}
\item[(A)] \lean{affineVacuum}: \emph{for every complex Lie algebra $\ltwog$ (in universe $u$), every
  symmetric invariant bilinear form $\kappa$ on it, and every $k\in\C$, the vacuum vector of
  $V^k(\ltwog)$ is nonzero.}
\item[(B)] \lean{inducedBottom}: \emph{for every complex Lie algebra $\ltwog$, every complex vector space
  $U$ (both in universe $u$), every Lie subalgebra $\ltwop\subseteq\ltwog$ and every representation
  $\rho$ of $\ltwop$ on $U$, the map $U\to U(\ltwog)\otimes_{U(\ltwop)}U$, $u\mapsto[1\otimes u]$, is
  injective.}
\end{itemize}
The braces \lean{∀ {𝔤 : Type u}} make $\ltwog$ an implicit argument (inferred from $\kappa$); the
square brackets are typeclass arguments (the Lie structure). \lean{theorem pbw : PBWPrinciple} proves
that this conjunction holds. Because the declaration mentions the universe variable \lean{u}, it is
\emph{universe-polymorphic}: it is in effect one theorem \texttt{pbw.\{u\}} for each universe level, and a
use site instantiates the level it needs (\lean{voaPBWBottom}\allowbreak\lean{Injective} writes
\texttt{(VOA.pbw : VOA.PBWPrinciple.\{u\})}).

\textbf{Decoding the proof term.} The anonymous constructor \lean{⟨…, …⟩} builds the structure from its
two fields, in order. The first field is supplied by \lean{AffineLie.}\allowbreak\lean{vacuumVector_}\allowbreak\lean{ne_}\allowbreak\lean{zero_}\allowbreak\lean{of_}\allowbreak\lean{induction}
(Theorem~\ref{l02:thm:vacind} below) and the second by \lean{PBWTheorem.}\allowbreak\lean{bottom_}\allowbreak\lean{injective}
(Theorem~\ref{l02:thm:bottominj} below). The second is stated in the library for the copies
\lean{PBWTheorem.balancing} and \lean{PBWTheorem.bottom} of the definitions above, over an arbitrary
field \lean{K}; at \lean{K = ℂ} and with $\ltwog$, $U$ in the same universe they unfold to
\lean{PBWInduction.balancing} and \lean{PBWInduction.bottom} symbol for symbol, which is why the
term \lean{fun p rho => PBWTheorem.bottom_injective p rho} is accepted at the type of the field
\lean{inducedBottom} without any conversion lemma.

\textbf{Why two fields in one structure (history).} While \lean{pbw} was an axiom (until
19 October 2026), bundling both consequences into one structure gave the library a single named axiom
with a single universe parameter; the structure was kept unchanged when the axiom became a theorem, so
that none of its use sites had to be edited. The type \lean{PBWPrinciple} is exactly the type the axiom
had; only the keyword and the proof changed.
\end{dictionary}

\begin{physics}
\lean{PBWPrinciple} packages the two places where a conformal field theorist silently uses the
Poincar\'e--Birkhoff--Witt theorem: that the vacuum sector of a current algebra, built by acting with the
creation currents $J^a_{-n}$ on a vacuum annihilated by the $J^a_{n\ge0}$, is not the zero space; and that
when one builds a module of the mode algebra of a vertex operator algebra by acting freely on a space of
lowest-weight states (a generalised Verma module), those lowest-weight states survive undamaged.
\end{physics}

\subsection{The fidelity check: is \texorpdfstring{\texttt{PBWPrinciple}}{PBWPrinciple} the intended statement?}

Now that \lean{pbw} is a theorem, Lean itself guarantees that \lean{PBWPrinciple} is \emph{true} (relative
to Lean's logic and its three standard axioms). What Lean cannot guarantee is that the Lean statement
\emph{means} what its name says. The danger is entirely in the \emph{encoding}: a Lean definition that
does not mean what it looks like could make \lean{PBWPrinciple} a true but weaker, or even vacuous,
statement, and then the consumers of \lean{pbw} (Section~\ref{l02:sec:consequences}) would be proving
less than they claim. We therefore identify, on paper, each field with the classical PBW consequence it
is meant to be: the following three results do this, and the list of traps after them explains which
slips in the encoding would have made the statement false (and the Lean proof impossible) or would have
emptied it of content.

(Before October 2026 this check was the only safeguard against an inconsistent library;
Section~\ref{l02:sec:history} explains why.)

\begin{proposition}[The Lean induced module is the induced module]\label{l02:prop:induced}
Let $\ltwop\subseteq\ltwog$ be a Lie subalgebra, $\rho$ a representation of $\ltwop$ on $U$, and
$\varphi : U(\ltwop)\to U(\ltwog)$ the algebra map induced by the inclusion. Make $U(\ltwog)$ a right
$U(\ltwop)$-module through $\varphi$ and $U$ a left $U(\ltwop)$-module through $\rho$. Let
$N\subseteq U(\ltwog)\otimes_\C U$ be the span of $\{y\iota(x)\otimes u - y\otimes\rho(x)u : y\in U(\ltwog),
x\in\ltwop,u\in U\}$. Then $N$ equals the span $N'$ of
$\{y\varphi(a)\otimes u - y\otimes a\cdot u : y\in U(\ltwog), a\in U(\ltwop), u\in U\}$, so
\begin{equation*}
 \lean{Induced p rho} = (U(\ltwog)\otimes_\C U)/N = U(\ltwog)\otimes_{U(\ltwop)}U.
\end{equation*}
\end{proposition}

\begin{proof}
$N\subseteq N'$: take $a = \iota_{\ltwop}(x)$, noting $\varphi(\iota_{\ltwop}(x)) = \iota(x)$ and
$\iota_{\ltwop}(x)\cdot u = \rho(x)u$. For $N'\subseteq N$, let $A$ be the set of $a\in U(\ltwop)$ such that
$y\varphi(a)\otimes u - y\otimes a u\in N$ for all $y, u$. $A$ is a linear subspace, contains $1$
(the element is $0$) and contains $\iota_{\ltwop}(\ltwop)$. It is closed under products: for $a,b\in A$,
\begin{equation*}
 y\varphi(ab)\otimes u - y\otimes ab\,u = \bigl[(y\varphi(a))\varphi(b)\otimes u - y\varphi(a)\otimes b u\bigr]
 + \bigl[y\varphi(a)\otimes (bu) - y\otimes a(bu)\bigr],
\end{equation*}
and both brackets are in $N$ (the first because $b\in A$, with $y\varphi(a)$ in place of $y$; the
second because $a\in A$, with $bu$ in place of $u$). Since $U(\ltwop)$ is generated as an algebra by
$\iota_{\ltwop}(\ltwop)$, $A = U(\ltwop)$. Finally $(U(\ltwog)\otimes_\C U)/N'$ is by definition (the
standard presentation of the tensor product over a ring) $U(\ltwog)\otimes_{U(\ltwop)}U$.
\end{proof}

\begin{physics}
$U(\ltwog)\otimes_{U(\ltwop)}U$ is the module obtained by declaring the states of $U$ to be ``lowest''
states on which the subalgebra $\ltwop$ acts as prescribed by $\rho$, and then acting freely with the
remaining generators. The balancing relations say only that moving an element of $\ltwop$ through to the
lowest states is the same as letting it act there; Proposition~\ref{l02:prop:induced} says that imposing
this for single generators of $\ltwop$ is enough.
\end{physics}

\begin{theorem}[Field (B) is true]\label{l02:thm:B}
For every complex Lie algebra $\ltwog$, Lie subalgebra $\ltwop$, and representation $\rho$ of $\ltwop$ on
any complex vector space $U$, the map $U\to U(\ltwog)\otimes_{U(\ltwop)}U$, $u\mapsto 1\otimes u$, is
injective.
\end{theorem}

\begin{proof}
By Corollary~\ref{l02:cor:free}(2), $U(\ltwog) = \bigoplus_m m\,U(\ltwop)$ as right $U(\ltwop)$-modules,
with $m$ running over ordered monomials in a basis of a complement of $\ltwop$, and the empty
monomial $m=1$ among them. Tensor product over $U(\ltwop)$ commutes with direct sums, and
$m\,U(\ltwop)\cong U(\ltwop)$ as right modules (freeness), so
\begin{equation*}
 U(\ltwog)\otimes_{U(\ltwop)}U \cong \bigoplus_m \bigl(m\,U(\ltwop)\bigr)\otimes_{U(\ltwop)}U\cong\bigoplus_m U,
\end{equation*}
and under this isomorphism $1\otimes u$ goes to $u$ in the summand $m=1$. Hence $u\mapsto 1\otimes u$
is injective. By Proposition~\ref{l02:prop:induced} this is the Lean map \lean{bottom p rho}.
\end{proof}

\begin{physics}
In the vertex-algebra application (Zhu's theory) $U$ is a space of lowest-weight states and the induced
module is the generalised Verma module built on it. Theorem~\ref{l02:thm:B} says that the lowest-weight
states one started with survive in the generalised Verma module, undeformed: nothing in the commutation
relations forces two of them to be identified.
\end{physics}

\begin{theorem}[Field (A) is true]\label{l02:thm:A}
For every complex Lie algebra $\ltwog$, every symmetric invariant bilinear form $\kappa$ (possibly
degenerate or zero) and every $k\in\C$ (including the critical level when $\ltwog$ is simple), the
vacuum vector $[1]\in U(\ltwoghat)/J(k)$ is nonzero; equivalently $1\notin J(k)$.
\end{theorem}

\begin{proof}
Let $\ltwon_- = \bigoplus_{n<0}\ltwog\otimes t^n$ and $\hat{\ltwop} = \bigoplus_{n\ge0}\ltwog\otimes t^n
\oplus\C C\oplus\C\ell_0$. As vector spaces $\ltwoghat = \ltwon_-\oplus\hat{\ltwop}$, and both are Lie
subalgebras: brackets of negative modes have negative degree and no central term (the cocycle needs
$m+n=0$); brackets of non-negative modes have non-negative degree and central term
$m\,\kappa(x,y)\delta_{m+n,0}C$ with $m=n=0$, hence zero; $[\ell_0, x\otimes t^n] = n\,x\otimes t^n$ stays
in the same degree; $C$ is central. Define $\chi:\hat{\ltwop}\to\C$ by $\chi(x\otimes t^n) = 0$,
$\chi(C) = k$, $\chi(\ell_0) = 0$. Then $\chi$ vanishes on $[\hat{\ltwop},\hat{\ltwop}]$, because by the
computation just made every bracket of two elements of $\hat{\ltwop}$ lies in $\bigoplus_{n\ge0}\ltwog\otimes
t^n$ (no $C$- and no $\ell_0$-component). So $\chi$ is a Lie algebra homomorphism to the abelian Lie
algebra $\C$, and extends to an algebra homomorphism $\tilde\chi : U(\hat{\ltwop})\to\C$.

By Corollary~\ref{l02:cor:free}(3), multiplication $\mu: U(\ltwon_-)\otimes U(\hat{\ltwop})\to U(\ltwoghat)$ is
a linear isomorphism. Define $\Phi : U(\ltwoghat)\to U(\ltwon_-)$ by $\Phi(\mu(u\otimes v)) =
\tilde\chi(v)\,u$. We claim $\Phi(J(k)) = 0$. The left ideal $J(k)$ is spanned by the elements
$y\,(a-\chi(a))$ with $y\in U(\ltwoghat)$ and $a$ one of the generators $x\otimes t^n$ ($n\ge0$), $C$,
$\ell_0$ (for which $a - \chi(a)$ is $x\otimes t^n$, $C-k$, $\ell_0$ respectively). Write
$y = \sum_i u_iv_i$ with $u_i\in U(\ltwon_-)$, $v_i\in U(\hat{\ltwop})$. Then
$y(a-\chi(a)) = \sum_i u_i\,\bigl(v_i(a-\chi(a))\bigr)$ with $v_i(a-\chi(a))\in U(\hat{\ltwop})$, so
\begin{equation*}
 \Phi\bigl(y(a-\chi(a))\bigr) = \sum_i\tilde\chi\bigl(v_i(a-\chi(a))\bigr)\,u_i =
 \sum_i\tilde\chi(v_i)\bigl(\chi(a)-\chi(a)\bigr)u_i = 0.
\end{equation*}
But $\Phi(1) = \Phi(\mu(1\otimes1)) = 1\ne0$ in $U(\ltwon_-)$ (the augmentation $U(\ltwon_-)\to\C$ sends
$1\mapsto1$, so $1\ne0$). Hence $1\notin J(k)$.
\end{proof}

\begin{remark}[Field (A) is a special case of field (B)]\label{l02:rem:AfromB}
Mathematically, (A) follows from (B): with $\ltwop = $ \lean{parabolic} and $U = \C$ with the character
$\chi$ above, the maps $[y]\mapsto[y\otimes1]$ and $[y\otimes c]\mapsto c[y]$ are mutually inverse
isomorphisms between $U(\ltwoghat)/J(k)$ and $\ltwoInd_{\hat{\ltwop}}^{\ltwoghat}\C_\chi$, sending $\ltwovac$ to
$[1\otimes1] = $ \lean{bottom} $1$, which is nonzero by (B). (One must check that the Lean
\lean{parabolic}, defined as a generated subalgebra, is the subspace $\hat{\ltwop}$, and that $\chi$ is a
Lie homomorphism on it; both are the computation in the proof of Theorem~\ref{l02:thm:A}.) Since
\lean{pbw} became a theorem the library \emph{does} formalise a version of this derivation: it is how the
first field of \lean{pbw} is now proved (Theorem~\ref{l02:thm:vacind}). The formal version avoids
\lean{parabolic} and instead uses a subalgebra \lean{nonnegSub} defined directly as the subspace
$\hat{\ltwop}$, and it needs only one direction of the isomorphism (that $1\in J(k)$ would force
$[1\otimes1] = 0$). The structure \lean{PBWPrinciple} still lists (A) as a separate field, and its
docstring's remark ``no implication from the first field to the second is claimed'' remains accurate:
(B) yields (A) by the short argument above, whereas no comparably short argument gives (B) from (A).
\end{remark}

\subsection{How the statement could have been false or empty, and why it is not}

Before \lean{pbw} became a theorem, a one-character slip in any of the definitions above could have made the
Lean statement false, and with it the library inconsistent. Today such a slip would simply have made the
Lean proof impossible. The list below is kept because it is also the list of the places where the
statement could have drifted away from its intended meaning, and a referee should still check each one:
the evidence is that each slip is absent.

\begin{warning}
\textbf{1. The representation must be a Lie homomorphism.} \lean{rho} has type
\lean{p →ₗ⁅ℂ⁆ Module.End ℂ U}, a Lie algebra homomorphism. Had it been only a linear map
\lean{p →ₗ[ℂ] Module.End ℂ U}, field (B) would be \emph{false}. Counterexample: $\ltwog = \ltwop$ the
two-dimensional abelian Lie algebra with basis $x,y$; $U = \C^2$; $\rho(x)$, $\rho(y)$ two
non-commuting matrices. In $U(\ltwog)$, $\iota x\,\iota y = \iota y\,\iota x$, and the balancing relations
give $[\iota x\iota y\otimes u] = [\iota x\otimes\rho(y)u] = [1\otimes\rho(x)\rho(y)u]$ and likewise
$[\iota y\iota x\otimes u] = [1\otimes\rho(y)\rho(x)u]$; so the bottom map identifies $\rho(x)\rho(y)u$ with
$\rho(y)\rho(x)u$ and is not injective. The Lean type rules this out.

\textbf{2. The side of the balancing relation.} The relation is $y\,\iota(x)\otimes u\sim y\otimes\rho(x)u$:
$\iota(x)$ multiplies $y$ on the \emph{right}, which is what defines the tensor product over $U(\ltwop)$
of the \emph{right} module $U(\ltwog)$ with the left module $U$. The Lean term is
\lean{(y * UniversalEnvelopingAlgebra.ι ℂ (x : g)) ⊗ₜ[ℂ] u - y ⊗ₜ[ℂ] (rho x u)}: correct side.

\textbf{3. The vacuum ideal must be generated by $C-k$, not $C$.} If $J(k)$ contained $C$ and $C-k$
with $k\ne0$, it would contain $k$, hence $1$. The Lean generator is
\lean{ι ℂ (C κ hκ) - algebraMap ℂ (UEnv 𝔤 κ hκ) k}, i.e.\ $C-k$, and $C$ itself is not a generator.

\textbf{4. The vacuum ideal must not contain modes of both signs.} If $J(k)$ contained $x\otimes t$ and
$y\otimes t^{-1}$ with $\kappa(x,y) = 1$, then it would contain
$[x\otimes t, y\otimes t^{-1}] = [x,y]\otimes1 + C$ (a commutator of two elements of a left ideal is in the
ideal: $ab - ba$ with $a,b\in J$), hence, since $[x,y]\otimes 1\in J$, also $C$, and with $C-k$ also $k$;
for $k\ne0$, $1\in J$. The Lean condition is \lean{0 ≤ n} on all loop generators: no negative modes.

\textbf{5. The zero modes and the cocycle.} The generators include the zero modes $x\otimes t^0$. This
is consistent only because the Kac--Moody cocycle is $m\,\kappa(x,y)\delta_{m+n,0}$ with the factor $m$,
which vanishes at $m=n=0$. With a term such as $\kappa(x,y)\delta_{m+n,0}$ (no factor $m$; this is
not even alternating, so not a valid cocycle) one would get $[x\otimes1,y\otimes1]\ni\kappa(x,y)C$ and the same collapse as in item 4.
The Lean \lean{cocycle_single} has the factor \lean{(m : ℂ)}.

\textbf{6. The derivation.} $\ell_0$ is a generator, i.e.\ $\ell_0\ltwovac=0$. This is consistent
because $[\ell_0, x\otimes t^n] = n\,x\otimes t^n$ maps generators to multiples of generators, and $\ell_0$
never appears in a bracket (the third component of \lean{bracketP} is \lean{0}). Had brackets of
non-negative modes produced an $\ell_0$- or $C$-component, the character $\chi$ of the proof of
Theorem~\ref{l02:thm:A} would not exist and the argument would break.

\textbf{7. Hidden typeclass mismatch.} The tensor product \lean{UniversalEnvelopingAlgebra ℂ g ⊗[ℂ] U}
uses the $\C$-module structure coming from the $\C$-algebra structure of $U(\ltwog)$, and \lean{U}'s given
\lean{[Module ℂ U]}; \lean{Module.End ℂ U} uses the same. There is no second, incompatible module
structure in play, so the Lean objects are the mathematical ones.

\textbf{8. Is the statement secretly about something trivial?} No: it quantifies over all Lie
algebras and all representations, and its instances include genuinely nontrivial ones (for example
$\ltwog = \mathfrak{sl}_2$, $\ltwop$ a Borel subalgebra, $U$ a one-dimensional weight: (B) is then the
nonvanishing of the Verma module highest-weight vector).
\end{warning}

\begin{dictionary}
\textbf{Edge cases, checked by hand.} A false universal statement is often false at a degenerate
instance; we check them.
\begin{itemize}
\item $U = 0$: any map out of the zero space is injective.
\item $\ltwop = \ltwog$: $U(\ltwog)\otimes_{U(\ltwog)}U\cong U$ via $y\otimes u\mapsto\tilde\rho(y)u$, and this
  sends $[1\otimes u]\mapsto u$; injectivity needs no PBW.
\item $\ltwop = 0$: the balancing relations are $y\iota(0)\otimes u - y\otimes\rho(0)u = 0$, so
  \lean{Induced} $=U(\ltwog)\otimes U$ and \lean{bottom} is $u\mapsto1\otimes u$; it has the left inverse
  $\varepsilon\otimes\mathrm{id}$, where $\varepsilon : U(\ltwog)\to\C$ is the augmentation (the algebra map
  extending the zero Lie map), which satisfies $\varepsilon(1) = 1$. Injective, again without PBW.
\item $\ltwog = 0$ in (A): $\ltwoghat = \C C\oplus\C\ell_0$ is abelian, $J(k)$ is generated by $C-k$ and
  $\ell_0$, and $U(\ltwoghat) = \C[C,\ell_0]$; $1\notin(C-k,\ell_0)$ because evaluation at $(k,0)$ sends
  $1\mapsto1$ and the ideal to $0$.
\item $\kappa = 0$ in (A): $\ltwoghat = L\ltwog\oplus\C C\oplus\C\ell_0$ with $C$ central and never a
  bracket; the character $\chi$ extends to all of $\ltwoghat$ ($x\otimes t^n\mapsto0$ for all $n$,
  $C\mapsto k$, $\ell_0\mapsto0$) because $C$ and $\ell_0$ are not in $[\ltwoghat,\ltwoghat]$; so the
  one-dimensional module $\C_\chi$ is a nonzero quotient of $V^k(\ltwog)$ and $\ltwovac\ne0$, without PBW.
\item $k = 0$, or $k = -h^\vee$ (critical level) for simple $\ltwog$: nothing changes in the proof of
  Theorem~\ref{l02:thm:A}, which never divides by $k+h^\vee$. (At the critical level the Sugawara
  construction fails, but that is irrelevant to the nonvanishing of the vacuum.)
\item Infinite-dimensional $\ltwog$ or $U$: PBW holds without finiteness; the bases exist by Zorn's lemma,
  which is available through \lean{Classical.choice}.
\end{itemize}
\textbf{Verdict.} After every definition is unfolded, (A) is Theorem~\ref{l02:thm:A} and (B) is
Theorem~\ref{l02:thm:B}, both true for all instances of the quantifiers. \lean{PBWPrinciple} is therefore
exactly the pair of classical PBW consequences its name advertises --- neither stronger nor weaker --- and
it is in fact a fairly weak consequence of PBW: it asserts injectivity of one map, not the existence of a
basis. Field (A) follows from (B) by a short argument (Remark~\ref{l02:rem:AfromB}), so (B) carries the substance. Its
instances include genuinely nontrivial ones (item 8 of the warning above), so it is not vacuous.
\end{dictionary}

\begin{physics}
Field (A) says that the WZW-type vacuum sector built from the vacuum by the creation currents
$J^a_{-n}$ is not the zero space --- the vacuum conditions $J^a_{n\ge0}\ltwovac = 0$, $C = k$,
$\ell_0\ltwovac = 0$ are mutually consistent at every level. Field (B) is the general version used in Zhu's
theory: when one builds a module of a vertex operator algebra by ``acting with all modes'' on a space of
lowest-weight states (a generalised Verma module), the lowest-weight states one started with survive
intact --- the construction does not identify two of them or kill any. Both are statements that
creation operators act freely, with no unexpected relations.
\end{physics}

%----------------------------------------------------------------------
\subsection{The proof in the library: relative PBW}\label{l02:sec:pbwproof}
%----------------------------------------------------------------------

The classical route to field (B) goes through the full PBW basis of $U(\ltwog)$
(Corollary~\ref{l02:cor:free}). The library takes a shorter route, which proves exactly what (B) needs
and nothing more. It is the \emph{relative} form of the key lemma in Humphreys' proof of PBW
(\emph{Introduction to Lie Algebras and Representation Theory}, \S17.4, Lemma A): instead of building a
representation of $\ltwog$ on a polynomial algebra $S(\ltwog)$ in which $1$ generates freely, one builds a
representation of $\ltwog$ on $S(\mathfrak c)\otimes U$, where $\mathfrak c$ is a complement of $\ltwop$, in
which the given representation $U$ of $\ltwop$ sits as the degree-zero part. Then (B) follows from the
universal property of induction, and (A) follows from (B). Everything is proved from Mathlib alone
(\leanpath{VOA/\allowbreak Foundations/\allowbreak PBWTheorem/\allowbreak Construction.lean},
\leanpath{LieRep.lean}, \leanpath{Induced.lean}); no PBW statement is assumed anywhere. The three files
work over an arbitrary field $K$, not only over $\C$.

\subsubsection{The relative PBW representation}

\begin{theorem}[Relative PBW representation; Humphreys \S17.4, Lemma A, relative form]\label{l02:thm:relrep}
Let $\ltwog$ be a Lie algebra over a field $K$ (of any dimension), $\ltwop\subseteq\ltwog$ a Lie subalgebra,
$\mathfrak c$ a vector-space complement of $\ltwop$ in $\ltwog$, $(y_a)_{a\in A}$ a basis of $\mathfrak c$ indexed by
a totally ordered set $A$, and $\rho$ a representation of $\ltwop$ on a $K$-vector space $U$. Let $M$ be the
vector space of finitely supported functions from finite multisets of $A$ to $U$; write $m\otimes u$ for
the function with value $u$ at the multiset $m$ and $0$ elsewhere, $|m|$ for the number of elements of
$m$ counted with multiplicity, and $a\le m$ for ``$a\le b$ for every letter $b$ of $m$''. Then there is a
representation $\sigma$ of $\ltwog$ on $M$ such that, for all $u\in U$, $x\in\ltwop$ and all multisets $m$:
\begin{align*}
 \text{(E1)}\quad& \sigma(y_a)(m\otimes u) = (a\,m)\otimes u &&\text{if } a\le m,\\
 \text{(E2)}\quad& \sigma(y_a)\bigl((b\,m)\otimes u\bigr) = \sigma(y_b)\sigma(y_a)(m\otimes u) + \sigma([y_a,y_b])(m\otimes u)
   &&\text{if } b<a,\ b\le m,\\
 \text{(E3)}\quad& \sigma(x)(\varnothing\otimes u) = \varnothing\otimes\rho(x)u,\\
 \text{(E4)}\quad& \sigma(x)\bigl((b\,m)\otimes u\bigr) = \sigma(y_b)\sigma(x)(m\otimes u) + \sigma([x,y_b])(m\otimes u)
   &&\text{if } b\le m,
\end{align*}
where $a\,m$ denotes the multiset $m$ with one more copy of $a$. In particular $u\mapsto\varnothing\otimes u$ is
an injective linear map $U\to M$ which intertwines $\rho(x)$ with $\sigma(x)$ for every $x\in\ltwop$.
\end{theorem}

The carrier and the standing variables.

File \leanpath{VOA/Foundations/PBWTheorem/Construction.lean}.
\begin{leancode}
variable {K : Type*} [Field K]
variable {g : Type*} [LieRing g] [LieAlgebra K g]
variable {U : Type*} [AddCommGroup U] [Module K U]
variable {A : Type*} [LinearOrder A]

abbrev Car (A U : Type*) [AddCommGroup U] [Module K U] := Multiset A →₀ U
\end{leancode}
and, with the further variables \lean{p}, \lean{hc}, \lean{bc}, \lean{ρ} (declared in the same way as in the
\leanpath{LieRep.lean} block below), the four defining equations of the limiting action \lean{act : g →ₗ[K] Module.End K (Car A U)} (a linear map, not yet known to be a Lie
homomorphism):
\begin{leancode}
theorem act_y_low {a : A} {m : Multiset A} (h : ∀ b ∈ m, a ≤ b) (u : U) :
    act p hc bc ρ (bc a : g) (single m u) = single (a ::ₘ m) u := by

theorem act_p_zero (x : p) (u : U) :
    act p hc bc ρ (x : g) (single 0 u) = single 0 (ρ x u) := by

theorem act_y_cons {a b : A} {m : Multiset A} (hba : b < a) (hb : ∀ c ∈ m, b ≤ c) (u : U) :
    act p hc bc ρ (bc a : g) (single (b ::ₘ m) u) =
      act p hc bc ρ (bc b : g) (act p hc bc ρ (bc a : g) (single m u)) +
        act p hc bc ρ ⁅(bc a : g), (bc b : g)⁆ (single m u) := by

theorem act_p_cons (x : p) {b : A} {m : Multiset A} (hb : ∀ c ∈ m, b ≤ c) (u : U) :
    act p hc bc ρ (x : g) (single (b ::ₘ m) u) =
      act p hc bc ρ (bc b : g) (act p hc bc ρ (x : g) (single m u)) +
        act p hc bc ρ ⁅(x : g), (bc b : g)⁆ (single m u) := by
\end{leancode}
Then the representation and its restriction to degree zero.

File \leanpath{VOA/Foundations/PBWTheorem/LieRep.lean}.
\begin{leancode}
variable (p : LieSubalgebra K g) {c : Submodule K g} (hc : IsCompl p.toSubmodule c)
variable (bc : Basis A K c) (ρ : p →ₗ⁅K⁆ Module.End K U)

theorem defect_eq_zero (z z' : g) (v : Car (K := K) A U) : defect p hc bc ρ z z' v = 0 := by

noncomputable def rep : g →ₗ⁅K⁆ Module.End K (Car (K := K) A U) :=

theorem rep_apply (z : g) (v : Car (K := K) A U) : rep p hc bc ρ z v = act p hc bc ρ z v := rfl

theorem rep_p_zero (x : p) (u : U) :
    rep p hc bc ρ (x : g) (single 0 u) = single 0 (ρ x u) :=
  act_p_zero p hc bc ρ x u
\end{leancode}
where \lean{defect p hc bc ρ z z'} is, by definition, the linear map
$\sigma(z)\sigma(z') - \sigma(z')\sigma(z) - \sigma([z,z'])$ (with \lean{act} in place of $\sigma$).

\begin{dictionary}
\textbf{Decoding.} \lean{Multiset A} is the type of finite multisets of elements of \lean{A}, i.e.\
unordered monomials $y_{a_1}\cdots y_{a_r}$ in the complement basis; \lean{Multiset A →₀ U} is the space
$M$ of the theorem, which is $S(\mathfrak c)\otimes_K U$ written in the monomial basis of $S(\mathfrak c)$.
\lean{single m u} is $m\otimes u$, \lean{single 0 u} is $\varnothing\otimes u$ (the empty multiset is
\lean{0}), and \lean{a ::ₘ m} is the multiset $a\,m$. \lean{hc : IsCompl p.toSubmodule c} says that
$\mathfrak c$ is a complement of $\ltwop$ ($\ltwop\cap\mathfrak c = 0$, $\ltwop+\mathfrak c = \ltwog$); \lean{bc} is a
basis of $\mathfrak c$, so \lean{(bc a : g)} is $y_a$ viewed in $\ltwog$; \lean{[LinearOrder A]} is the total
order. The hypothesis \lean{∀ b ∈ m, a ≤ b} is ``$a\le m$''; \lean{b < a} together with
\lean{∀ c ∈ m, b ≤ c} says that $b$ is the least letter of $b\,m$ and is strictly below $a$. The four
theorems are (E1), (E3), (E2), (E4) in this order. \lean{g →ₗ⁅K⁆ Module.End K M} is a Lie algebra
homomorphism into $\End_K M$ with the commutator bracket (Section~\ref{l02:sec:lean}), i.e.\ a
representation: so \lean{rep} is $\sigma$, and \lean{rep_apply} says it is the operator family
\lean{act} of the four equations. \lean{defect_eq_zero} is the statement that \lean{act} respects brackets,
which is exactly what is needed to upgrade \lean{act} to the Lie homomorphism \lean{rep}.
\lean{rep_p_zero} is the final sentence of Theorem~\ref{l02:thm:relrep}: on degree zero, $\ltwop$ acts
through $\rho$.

\textbf{Hypotheses and generality.} No finiteness of $\ltwog$, $\ltwop$, $U$ or $A$ is assumed, and $K$ is any
field. The equations do not mention $\ltwop$-components explicitly: \lean{act} is linear in its
$\ltwog$-argument, so \lean{act ⁅(bc a : g), (bc b : g)⁆} means ``decompose $[y_a,y_b]$ into its $\ltwop$- and
$\mathfrak c$-parts and act by each''. The construction depends on the choice of $\mathfrak c$, of the basis and of
the order; Theorem~\ref{l02:thm:relrep} asserts existence for every such choice.

\textbf{Could it be vacuous?} No: \lean{rep} is a genuine Lie homomorphism (its type requires the bracket
relation, proved as \lean{defect_eq_zero}), and \lean{rep_p_zero} together with the injectivity of
$u\mapsto\varnothing\otimes u$ (\lean{Finsupp.single_injective 0}) says that $U$ sits inside it as an honest
$\ltwop$-subrepresentation. Those two facts are exactly what is used downstream.
\end{dictionary}

\begin{proof}[Proof on paper]
\emph{Construction.} We define $\sigma(z)$ on $M_{\le d} = \operatorname{span}\{m\otimes u : |m|\le d\}$ by
induction on $d$, for $z$ running over the generators $y_a$ ($a\in A$) and $x\in\ltwop$, and extend linearly
in $z$ using $\ltwog = \ltwop\oplus\mathfrak c$. We carry along the \emph{degree bounds}
\begin{equation*}
 \sigma(y_a)(m\otimes u) - (a\,m)\otimes u\in M_{\le|m|},\qquad \sigma(x)(m\otimes u)\in M_{\le|m|}.
\end{equation*}
In degree $0$, (E1) and (E3) define everything. Suppose $\sigma$ is defined on $M_{\le d-1}$ and satisfies the
bounds, and let $|m| = d - 1$, $b\,m$ a monomial of degree $d$ with least letter $b$. If $a\le b\,m$, (E1)
defines $\sigma(y_a)((b\,m)\otimes u)$. If $b<a$, write $\sigma(y_a)(m\otimes u) = (a\,m)\otimes u + w$ with
$w\in M_{\le d-1}$ (degree bound at degree $d-1$). Since $b\le a\,m$, (E1) gives
$\sigma(y_b)((a\,m)\otimes u) = (b\,a\,m)\otimes u$, and $\sigma(y_b)w$ is already defined; so the right-hand
side of (E2) makes sense and \emph{defines} the left-hand side. The same argument, using the bound
$\sigma(x)(m\otimes u)\in M_{\le d-1}$, makes (E4) a definition. The new values satisfy the degree bounds, which
closes the induction.

\emph{The bracket relation.} We show $\sigma(z)\sigma(z')v - \sigma(z')\sigma(z)v = \sigma([z,z'])v$ for
generators $z, z'$ and basis vectors $v = m\otimes u$, by induction on $|m|$; bilinearity then gives it for
all $z, z'$, $v$. Call a generator $y_a$ \emph{low} for $m$ if $a\le m$.

\emph{Low case.} If one of the two generators is low for $m$ (if both are, take the one with the smaller
letter), put it in the second slot using skew-symmetry of both sides. With $z' = y_a$ low, $\sigma(y_a)v =
(a\,m)\otimes u$ by (E1), and the relation to be proved, $\sigma(z)((a\,m)\otimes u) = \sigma(y_a)\sigma(z)v +
\sigma([z,y_a])v$, is literally (E2) (if $z = y_c$ with $c>a$), (E4) (if $z\in\ltwop$), or trivial (if $z = y_a$).
This case needs no induction.

\emph{Degree zero.} If $m = \varnothing$ and $z, z'\in\ltwop$, the relation is $\rho([x,x']) =
[\rho(x),\rho(x')]$, i.e.\ that $\rho$ is a representation; every other pair at $m=\varnothing$ has a low
generator.

\emph{Inductive case.} Otherwise $m = b\,m'$ with $b$ the least letter, and both generators are
\emph{above} $b$ (each is some $y_a$ with $a>b$, or an element of $\ltwop$). Put $v' = m'\otimes u$, so
$v = \sigma(y_b)v'$, and assume the relation for all vectors of degree $<|m|$. By (E2)/(E4),
$\sigma(z)v = \sigma(y_b)\sigma(z)v' + \sigma([z,y_b])v'$, and similarly for $z'$. Apply $\sigma(z')$ to the
first expression and expand $\sigma(z')\sigma(y_b)(\sigma(z)v')$. Here $\sigma(z)v'$ is not a single monomial,
so we split it. By the degree bound, if $z = y_a$ then $\sigma(z)v' = (a\,m')\otimes u + w$ with $w$ of degree
$<|m|$, and $y_b$ is low for $a\,m'$ because $b<a$ and $b\le m'$; if $z\in\ltwop$ there is no leading monomial
and $\sigma(z)v' = w$ has degree $<|m|$. On the leading monomial the relation for $(z',y_b)$ holds by the
low case; on $w$ it holds by the inductive hypothesis (applied monomial by monomial and summed, since both
sides are linear). Hence $\sigma(z')\sigma(y_b)\sigma(z)v' = \sigma(y_b)\sigma(z')\sigma(z)v' +
\sigma([z',y_b])\sigma(z)v'$. Doing the same
with $z,z'$ exchanged and subtracting, the terms collect into
\begin{equation*}
 \sigma(y_b)\bigl([\sigma(z),\sigma(z')]v'\bigr) + \sigma([[z,y_b],z'])v' + \sigma([z,[z',y_b]])v',
\end{equation*}
where the inductive hypothesis (degree $|m'| < |m|$) was used for the pairs $(z,z')$, $([z,y_b],z')$ and
$(z,[z',y_b])$ on $v'$. The first term is $\sigma(y_b)\sigma([z,z'])v'$, again by induction, and by the Jacobi
identity $[[z,y_b],z'] + [z,[z',y_b]] = [[z,z'],y_b]$. Finally
$\sigma(y_b)\sigma([z,z'])v' + \sigma([[z,z'],y_b])v' = \sigma([z,z'])\sigma(y_b)v' = \sigma([z,z'])v$ by the low case
applied to $v'$ (the generator $y_b$ is low for $m'$), for each generator component of $[z,z']$. This is
Humphreys' computation; the only property of $\ltwog$ used is the Jacobi identity, and the only property of
$\rho$ used is that it is a representation.
\end{proof}

\begin{example}[The Verma module of $\mathfrak{sl}_2$]\label{l02:ex:relrep}
Let $\ltwog = \mathfrak{sl}_2$ with $[h,e] = 2e$, $[h,f] = -2f$, $[e,f] = h$, let $\ltwop = \operatorname{span}\{h,e\}$
(the Borel subalgebra), $\mathfrak c = \C f$ (one letter, so a monomial is $f^n$), and $U = \C u$ with
$\rho(h) = \lambda$, $\rho(e) = 0$. Then $M$ has basis $f^n\otimes u$, $n\ge0$. (E1) gives
$\sigma(f)(f^n\otimes u) = f^{n+1}\otimes u$. (E3) gives $\sigma(h)(\varnothing\otimes u) = \lambda\,\varnothing\otimes u$ and
$\sigma(e)(\varnothing\otimes u) = 0$. (E4) with $x = h$ gives
$\sigma(h)(f^n\otimes u) = \sigma(f)\sigma(h)(f^{n-1}\otimes u) - 2f^n\otimes u$, hence $\sigma(h)(f^n\otimes u) =
(\lambda-2n)f^n\otimes u$; with $x = e$ it gives $\sigma(e)(f^n\otimes u) = \sigma(f)\sigma(e)(f^{n-1}\otimes u) +
\sigma(h)(f^{n-1}\otimes u)$, hence $\sigma(e)(f^n\otimes u) = n(\lambda-n+1)f^{n-1}\otimes u$. These are the familiar
formulas of the Verma module $M(\lambda)$, and one checks directly that
$[\sigma(e),\sigma(f)] = \sigma(h)$ on $f^n\otimes u$: $(n+1)(\lambda-n) - n(\lambda-n+1) = \lambda-2n$.
\end{example}

\begin{physics}
The example is the highest-weight representation of the $SU(2)$ current zero modes built on a state of
$J^3_0$-eigenvalue $\lambda/2$: the lowering operator acts freely, and the raising operator acts by the
standard matrix elements.
\end{physics}

\begin{physics}
$M = S(\mathfrak c)\otimes U$ is the ``Fock space'' one expects for a generalised Verma module: states are
created from the lowest-weight space $U$ by the creation operators $y_a$ (a basis of the modes not in
$\ltwop$), \emph{normal-ordered} (the letters of a monomial are listed in increasing order). (E1) says that
applying a creation operator that is already in normal order just adds it to the monomial; (E2)--(E4) say
that any other operator is moved to its normal-ordered place by commuting it through, picking up the
commutator. Theorem~\ref{l02:thm:relrep} says that this normal-ordering prescription is consistent, i.e.\
defines an honest representation; physically, it is the statement that the canonical commutation
relations are the only relations among the modes.
\end{physics}

\begin{status}
\lean{VOA.}\allowbreak\lean{PBWTheorem.}\allowbreak\lean{rep}, \lean{VOA.}\allowbreak\lean{PBWTheorem.}\allowbreak\lean{rep_}\allowbreak\lean{p_}\allowbreak\lean{zero}, and the four defining equations are
proved from Mathlib alone; they enter the report of \lean{PBWTheorem.}\allowbreak\lean{bottom_}\allowbreak\lean{injective}
below, which is \lean{[propext, Classical.choice, Quot.sound]} (run for this chapter). Unconditional.
\end{status}

\subsubsection{Injectivity of the bottom map}

\begin{theorem}[Bottom injectivity, relative PBW]\label{l02:thm:bottominj}
Let $\ltwog$ be a Lie algebra over a field $K$, $\ltwop$ a Lie subalgebra and $\rho$ a representation of $\ltwop$
on $U$. Then the map $U\to(U(\ltwog)\otimes_K U)/N$, $u\mapsto[1\otimes u]$, is injective, where $N$ is the span
of the balancing relations $y\,\iota(x)\otimes u - y\otimes\rho(x)u$. By Proposition~\ref{l02:prop:induced} this is
the map $U\to U(\ltwog)\otimes_{U(\ltwop)}U$, and over $K = \C$ it is field (B) of \lean{PBWPrinciple}.
More generally: if there are a representation $\sigma$ of $\ltwog$ on some $M$ and an injective linear
map $j : U\to M$ with $\sigma(x)j(u) = j(\rho(x)u)$ for $x\in\ltwop$, then the bottom map is injective.
\end{theorem}

File \leanpath{VOA/Foundations/PBWTheorem/Induced.lean}.
\begin{leancode}
def balancingSet (p : LieSubalgebra K g) (rho : p →ₗ⁅K⁆ Module.End K U) :
    Set (UniversalEnvelopingAlgebra K g ⊗[K] U) :=
  {t | ∃ (y : UniversalEnvelopingAlgebra K g) (x : p) (u : U),
    t = (y * UniversalEnvelopingAlgebra.ι K (x : g)) ⊗ₜ[K] u - y ⊗ₜ[K] (rho x u)}

noncomputable def balancing (p : LieSubalgebra K g) (rho : p →ₗ⁅K⁆ Module.End K U) :
    Submodule K (UniversalEnvelopingAlgebra K g ⊗[K] U) :=
  Submodule.span K (balancingSet p rho)

noncomputable def bottom (p : LieSubalgebra K g) (rho : p →ₗ⁅K⁆ Module.End K U) :
    U →ₗ[K] (UniversalEnvelopingAlgebra K g ⊗[K] U) ⧸ balancing p rho :=
  (balancing p rho).mkQ.comp (TensorProduct.mk K (UniversalEnvelopingAlgebra K g) U 1)

theorem bottom_injective_of_embedding (p : LieSubalgebra K g) (rho : p →ₗ⁅K⁆ Module.End K U)
    {M : Type*} [AddCommGroup M] [Module K M] (σ : g →ₗ⁅K⁆ Module.End K M) (j : U →ₗ[K] M)
    (hj : Function.Injective j) (hσ : ∀ (x : p) (u : U), σ (x : g) (j u) = j (rho x u)) :
    Function.Injective (bottom p rho) := by

theorem bottom_injective (p : LieSubalgebra K g) (rho : p →ₗ⁅K⁆ Module.End K U) :
    Function.Injective (bottom p rho) := by
\end{leancode}

\begin{dictionary}
\textbf{Same definitions as \lean{PBWInduction}.} Compare with the block of
Section~\ref{l02:sec:axiom} that defines \lean{PBWInduction.balancingSet}, \lean{balancing} and
\lean{bottom}: the three definitions here are the same text with \lean{ℂ} replaced by a field
\lean{K}, with the abbreviations \lean{Cover} and \lean{Induced} unfolded, and with $\ltwog$ and $U$
allowed to live in different universes. At \lean{K = ℂ} they are \emph{definitionally} equal to the
\lean{PBWInduction} ones, which is why \lean{pbw} can use \lean{bottom_injective} for the field
\lean{inducedBottom} directly. So Theorem~\ref{l02:thm:bottominj} is field (B), for every field.

\textbf{The embedding criterion.} In \lean{bottom_injective_of_embedding}, \lean{σ} is a representation
of $\ltwog$ on $M$, \lean{j} a linear map with \lean{hj} (injective) and \lean{hσ} ($\ltwop$-equivariant: the
restriction of $\sigma$ to $\ltwop$, applied to $j(u)$, is $j$ of $\rho$). The conclusion is injectivity of
the bottom map. This is the second paragraph of Theorem~\ref{l02:thm:bottominj}.

\textbf{No hidden hypothesis.} \lean{bottom_injective} has no hypothesis beyond the data: every Lie
subalgebra, every representation. The choices needed in the proof (a complement, a basis, a total order)
are made inside the proof, using Mathlib's \lean{Submodule.exists_isCompl}, \lean{Basis.ofVectorSpace}
(which uses Zorn's lemma, via \lean{Classical.choice}) and a well-ordering of the index set
(\lean{WellOrderingRel}).
\end{dictionary}

\begin{proof}[Proof on paper]
\emph{The criterion.} By the universal property of $U(\ltwog)$, $\sigma$ extends to an algebra map
$\tilde\sigma : U(\ltwog)\to\End_K M$. The bilinear map $(y,u)\mapsto\tilde\sigma(y)j(u)$ gives a linear map
$L_0 : U(\ltwog)\otimes U\to M$, and $L_0$ kills every balancing relation:
$\tilde\sigma(y\iota(x))j(u) - \tilde\sigma(y)j(\rho(x)u) = \tilde\sigma(y)\bigl(\sigma(x)j(u) - j(\rho(x)u)\bigr) = 0$. So $L_0$
descends to $L$ on the quotient, with $L([1\otimes u]) = j(u)$. If $[1\otimes u_1] = [1\otimes u_2]$, then
$j(u_1) = j(u_2)$ and $u_1 = u_2$.

\emph{The theorem.} Choose a complement $\mathfrak c$ of $\ltwop$ (every subspace of a vector space has one), a
basis of $\mathfrak c$ (every vector space has one), and a total order on its index set (every set can be
well-ordered). Theorem~\ref{l02:thm:relrep} gives $\sigma$ on $M$, and $j(u) = \varnothing\otimes u$ is
injective and $\ltwop$-equivariant by (E3). Apply the criterion.
\end{proof}

\begin{physics}
The criterion is the familiar physicist's argument ``to show that a state is nonzero, exhibit a
representation in which it is nonzero'': the generalised Verma module is the \emph{freest} module
generated by $U$ with the prescribed $\ltwop$-action, so it maps onto any other module containing such a
copy of $U$, and if $U$ survives there it survived in the Verma module. Theorem~\ref{l02:thm:relrep}
supplies the other module: the normal-ordered Fock space $S(\mathfrak c)\otimes U$.
\end{physics}

\begin{status}
\lean{VOA.}\allowbreak\lean{PBWTheorem.}\allowbreak\lean{bottom_}\allowbreak\lean{injective}: \lean{[propext, Classical.choice, Quot.sound]}
(run for this chapter, Section~\ref{l02:sec:verification}). Unconditional.
\end{status}

\subsubsection{The affine vacuum from relative PBW}

\begin{theorem}[Nonvanishing of the affine vacuum, from bottom injectivity]\label{l02:thm:vacind}
For every complex Lie algebra $\ltwog$ with symmetric invariant form $\kappa$ and every $k\in\C$, the vacuum
vector of $V^k(\ltwog)$ is nonzero. The proof uses the subalgebra
$\hat{\ltwop} = \ltwog[t]\oplus\C C\oplus\C\ell_0$ of $\ltwoghat$ and its character $\chi_k(f + aC + p\ell_0) = ka$.
\end{theorem}

File \leanpath{VOA/Foundations/AffineLie.lean}.
\begin{leancode}
noncomputable def nonnegSub : LieSubalgebra ℂ (AffineLie 𝔤 κ hκ) where
  carrier := {X | X.carrier.1 ∈ supported 𝔤 ℂ (Set.Ici (0 : ℤ))}
  ...

noncomputable def vacuumChar (k : ℂ) : nonnegSub κ hκ →ₗ⁅ℂ⁆ Module.End ℂ ℂ where
  toFun X := (k * (X : AffineLie 𝔤 κ hκ).carrier.2.1) • LinearMap.id
  ...

theorem tmul_one_mem_balancing_of_mem_vacuumIdeal (k : ℂ) {a : UEnv 𝔤 κ hκ}
    (ha : a ∈ vacuumIdeal κ hκ k) :
    a ⊗ₜ[ℂ] (1 : ℂ) ∈ PBWTheorem.balancing (nonnegSub κ hκ) (vacuumChar κ hκ k) := by

theorem vacuumVector_ne_zero_of_induction (k : ℂ) : vacuumVector κ hκ k ≠ 0 := by
\end{leancode}

\begin{dictionary}
\textbf{Decoding.} \lean{supported 𝔤 ℂ (Set.Ici (0 : ℤ))} is the subspace of \lean{ℤ →₀ 𝔤} of Laurent
polynomials supported in degrees $\ge0$, i.e.\ $\ltwog[t]$. So \lean{nonnegSub} is the set of triples
$(f,a,p)$ with $f\in\ltwog[t]$ and $a,p$ arbitrary: the subspace $\hat{\ltwop}$, defined \emph{directly} as a
subspace (the omitted fields prove that it is closed under addition, scalars and the bracket, the last
from \lean{loopBracket_}\allowbreak\lean{mem_}\allowbreak\lean{nonneg} and \lean{loopDeriv_}\allowbreak\lean{mem_}\allowbreak\lean{nonneg}). This differs from \lean{parabolic}
(Section~\ref{l02:sec:affine}), which is defined as a generated subalgebra; the two are equal, but the
library does not need that. \lean{vacuumChar k} sends $X = (f,a,p)$ to multiplication by $ka$ on the
one-dimensional space \lean{ℂ} (\lean{carrier.2.1} is the $C$-coefficient): this is the character $\chi_k$,
packaged as a representation of \lean{nonnegSub} on \lean{ℂ}. Its omitted field \lean{map_lie'} is the check
that $\chi_k$ kills brackets: the $C$-component of a bracket of two elements of $\hat{\ltwop}$ is
$\omega(f,g)$, which vanishes on $\ltwog[t]$ (\lean{cocycle_nonneg}: the cocycle carries the factor $m$ and needs
$m + n = 0$). \lean{tmul_one_mem_balancing_of_mem_vacuumIdeal} says that $a\in J(k)$ implies
$a\otimes1\in N$, the balancing span of this induction; \lean{vacuumVector_ne_zero_of_induction} is
Theorem~\ref{l02:thm:vacind}, with no hypotheses.

\textbf{Same as field (A).} The conclusion is literally \lean{PBW κ hκ k}, unfolded; this is why the proof
term of \lean{pbw} can use it for \lean{affineVacuum}.
\end{dictionary}

\begin{proof}[Proof on paper]
Every generator of $J(k)$ is of the form $\iota(X) - \chi_k(X)$ with $X\in\hat{\ltwop}$: $X = x\otimes t^n$
($n\ge0$, $\chi_k = 0$), $X = C$ ($\chi_k = k$), $X = \ell_0$ ($\chi_k = 0$). For such a generator,
$(\iota(X) - \chi_k(X))\otimes1 = 1\cdot\iota(X)\otimes1 - 1\otimes\chi_k(X)1$ is a balancing relation (with
$y = 1$, $u = 1$). The balancing span is stable under left multiplication by $U(\ltwoghat)$ on the first
tensor factor (replace $y$ by $ry$), so $a\otimes1\in N$ for every $a$ in the left ideal $J(k)$. Now if the
vacuum vector were $0$, then $1\in J(k)$, so $1\otimes1\in N$, i.e.\ the bottom map sends $1\in\C$ to $0$, the
image of $0$. By Theorem~\ref{l02:thm:bottominj} the bottom map is injective, so $1 = 0$ in $\C$, which is
absurd.
\end{proof}

\begin{physics}
This is the statement that the vacuum of a current algebra is consistent at every level $k$, critical
or not, and for every Lie algebra and invariant form, degenerate or not: the vacuum sector is the
generalised Verma module induced from the one-dimensional representation ``level $k$, zero charge, zero
energy'' of the non-negative modes, and its lowest state survives.
\end{physics}

\begin{status}
\lean{VOA.}\allowbreak\lean{AffineLie.}\allowbreak\lean{vacuumVector_}\allowbreak\lean{ne_}\allowbreak\lean{zero_}\allowbreak\lean{of_}\allowbreak\lean{induction}: \lean{[propext, Classical.choice, Quot.sound]} (run for
this chapter). Unconditional.
\end{status}

\subsubsection{Status of \texorpdfstring{\texttt{VOA.pbw}}{VOA.pbw}}

\begin{status}
\texttt{\#print axioms VOA.pbw}, run for this chapter at an earlier commit (not re-run for this edition) with
\texttt{lake env lean -DmaxSynthPendingDepth=3} on a scratch file containing
\texttt{import VOA.Foundations.AffineLie} (output in Section~\ref{l02:sec:verification}), reports
\begin{plaincode}
'VOA.pbw' depends on axioms: [propext, Classical.choice, Quot.sound]
\end{plaincode}
\lean{VOA.pbw} is a theorem, proved from Mathlib alone, with no hypotheses: \textbf{unconditional}. A text
search for declarations beginning with \lean{axiom} in \leanpath{VOA/\allowbreak } finds none (only
docstring prose using the word). The project's \texttt{scripts/check\_repository.py}, which since \lean{pbw}
became a theorem rejects any \lean{axiom} declaration, was run for this chapter at that earlier commit and
reported ``7107 library files scanned; \ldots; 0 sanctioned axiom declaration(s)'' (its overall verdict was
``failed'', solely because two draft files, \leanpath{VOA/Huang/LogAssocIterateInputLinearity.lean} and
\leanpath{VOA/Huang/LogGenusOneModeRecurrence.lean}, were not imported by the umbrella; those drafts have since
been moved out, and a re-run on the working tree of 28 October 2026 printed ``Repository check passed: 10353
library files scanned; \ldots; 0 sanctioned axiom declaration(s)''). The axiom audit
\texttt{scripts/audit/run\_axiom\_audits.py} treats any reappearance of \lean{pbw} in an axiom list as a
failure (according to the commit message of the revision that made \lean{pbw} a theorem; not re-run for this chapter). Mathlib v4.30.0 still contains no PBW theorem
(Section~\ref{l02:sec:lean}); the library proves the \emph{relative} consequence above, and the PBW basis of
the affine vacuum module (Section~\ref{l02:sec:pbwbasis}), but not the PBW basis theorem for $U(\ltwog)$
itself (Theorem~\ref{l02:thm:pbw}).

\emph{History.} In an earlier run (for the first edition of this chapter) the same
command reported \lean{[propext, Classical.choice, Quot.sound, pbw]}: \lean{pbw} was then an axiom, and an
axiom's own report lists itself. A revision of 19 October 2026 replaced it by the theorem.
That revision's own commit message records that the aggregate \texttt{lake build} was not yet certified green at
the time (only the per-target builds of \leanpath{AffineLie.lean} and \leanpath{VacuumBasis.lean} were);
according to the project's build record, the aggregate \texttt{lake build VOA} was later run green
(13868 jobs, 26 October 2026, axiom audit PASS). No \texttt{lake build} was run for this chapter.
\end{status}

%======================================================================
\section{Consequences of the PBW principle, and the history of the axiom}\label{l02:sec:consequences}
%======================================================================

\subsection{Nonvanishing of the affine vacuum and the simple quotient}

\begin{theorem}\label{l02:thm:vacne}
For every complex Lie algebra $\ltwog$ with invariant form $\kappa$ and every $k\in\C$:
$\ltwovac\ne0$ in $V^k(\ltwog)$; hence $V^k(\ltwog)$ is nonzero; it has a maximal proper
$U(\ltwoghat)$-submodule $N$; $\ltwovac\notin N$; and $L(0) = V^k(\ltwog)/N$ is a simple
$U(\ltwoghat)$-module.
\end{theorem}

File \leanpath{VOA/Foundations/AffineLie.lean}.
\begin{leancode}
theorem vacuumVector_ne_zero (k : ℂ) : vacuumVector κ hκ k ≠ 0 :=
  VOA.pbw.affineVacuum κ hκ k

theorem vacuumSpace_nontrivial (k : ℂ) : Nontrivial (vacuumSpace κ hκ k) :=
  ⟨⟨vacuumVector κ hκ k, 0, vacuumVector_ne_zero κ hκ k⟩⟩

theorem exists_vacuumMaxSubmodule (k : ℂ) :
    ∃ P : Submodule (UEnv 𝔤 κ hκ) (vacuumSpace κ hκ k), IsCoatom P := by

noncomputable def vacuumMaxSubmodule (k : ℂ) :
    Submodule (UEnv 𝔤 κ hκ) (vacuumSpace κ hκ k) :=
  (exists_vacuumMaxSubmodule κ hκ k).choose

theorem vacuumVector_not_mem_maxSubmodule (k : ℂ) :
    vacuumVector κ hκ k ∉ vacuumMaxSubmodule κ hκ k := by

noncomputable abbrev vacuumSimple (k : ℂ) : Type _ :=
  vacuumSpace κ hκ k ⧸ vacuumMaxSubmodule κ hκ k

theorem vacuumSimple_isSimpleModule (k : ℂ) :
    IsSimpleModule (UEnv 𝔤 κ hκ) (vacuumSimple κ hκ k) := by
\end{leancode}

\begin{dictionary}
\textbf{Decoding.} \lean{vacuumVector_}\allowbreak\lean{ne_zero} is a one-line term: it is \emph{literally} the field
\lean{affineVacuum} of the theorem \lean{VOA.pbw} applied to \lean{κ hκ k}; the shown proof is the whole
proof.
\lean{Nontrivial X} means $X$ has two distinct elements. \lean{IsCoatom P} means $P\ne\top$ and every
submodule strictly above $P$ is $\top$: $P$ is a maximal proper submodule. \lean{.choose} picks one
(non-constructively). \lean{IsSimpleModule R M} is Mathlib's notion of a simple module (nonzero, with no
submodules other than $0$ and $M$).

\textbf{Same as the textbook?} The textbook $L(0)$ (often $L_k(\ltwog)$) is the quotient by \emph{the}
maximal proper submodule. The library quotients by \emph{a chosen} maximal proper submodule. For
$V^k(\ltwog)$ these agree, because the maximal proper submodule is unique: $\ell_0$ acts diagonalisably
(\lean{ell0_isInternal}), so every submodule is the sum of its intersections with the eigenspaces; a
submodule is proper iff its eigenvalue-$0$ part, a subspace of $\C\ltwovac$, is zero; so the sum of all
proper submodules is proper. The library does not prove uniqueness here, so a referee should read
\lean{vacuumSimple} as ``a simple quotient of $V^k(\ltwog)$'', which by the argument just given is the
simple quotient. The simple affine VOAs of the library (Chapter~\ref{ch:L08}) are built on a different
route, as quotients of vertex algebras by their maximal graded ideal.

\textbf{Dependence.} Every theorem in the block except the definitions depends on \lean{VOA.pbw},
because the existence of a maximal proper submodule needs $V^k(\ltwog)\ne0$ (a zero module has none).
Since \lean{VOA.pbw} is now a theorem, this dependence is an ordinary use of a proved result: the block is
unconditional. (The proof term of \lean{vacuumVector_}\allowbreak\lean{ne_zero} still reads
\lean{VOA.pbw.affineVacuum κ hκ k} rather than calling
\lean{vacuumVector_}\allowbreak\lean{ne_}\allowbreak\lean{zero_}\allowbreak\lean{of_}\allowbreak\lean{induction} directly; the two are the same proof.)
\end{dictionary}

\begin{physics}
$L(0)$ is the vacuum module of the chiral WZW-type theory in which all null states have been set to
zero: the irreducible representation of the current algebra generated by the vacuum. Nonvanishing of
$\ltwovac$ says the theory has a vacuum at all.
\end{physics}

\begin{proof}[Proof on paper]
$\ltwovac\ne0$ is field (A) (Theorem~\ref{l02:thm:A}; in Lean, Theorem~\ref{l02:thm:vacind}). A cyclic module $R\cdot v$ is finitely generated, so
its top submodule is a compact element of the submodule lattice, which is therefore coatomic: every
proper submodule lies in a maximal one (Zorn's lemma, through \lean{CompleteLattice.}\allowbreak\lean{coatomic_}\allowbreak\lean{of_}\allowbreak\lean{top_}\allowbreak\lean{compact}).
The zero submodule is proper because $\ltwovac\ne0$, so a maximal proper $N$ exists. If $\ltwovac\in N$ then
$N\supseteq R\ltwovac = V^k(\ltwog)$, contradiction. A quotient by a maximal proper submodule is simple.
\end{proof}

\begin{status}
\lean{VOA.}\allowbreak\lean{AffineLie.}\allowbreak\lean{vacuumVector_}\allowbreak\lean{ne_}\allowbreak\lean{zero}, \lean{VOA.}\allowbreak\lean{AffineLie.}\allowbreak\lean{vacuumSpace_}\allowbreak\lean{nontrivial},
\lean{VOA.}\allowbreak\lean{AffineLie.}\allowbreak\lean{exists_}\allowbreak\lean{vacuumMax}\allowbreak\lean{Submodule}, \lean{VOA.}\allowbreak\lean{AffineLie.}\allowbreak\lean{vacuumSimple_}\allowbreak\lean{isSimpleModule}:
\lean{[propext, Classical.choice, Quot.sound]} (run for this chapter). Unconditional. They use the
theorem \lean{VOA.pbw} (field \lean{affineVacuum}). \emph{History:} in an earlier run, when \lean{pbw}
was an axiom, all four reported \lean{[propext, Classical.choice, Quot.sound, pbw]}.
\end{status}

\subsection{The two direct consumers}

A text search of the library source for uses of \lean{VOA.pbw} in proof terms (as opposed to
documentation) finds exactly two (re-checked at an earlier commit, after \lean{pbw} became a theorem; not re-run for this edition):
\begin{enumerate}
\item \lean{AffineLie.}\allowbreak\lean{vacuumVector_}\allowbreak\lean{ne_}\allowbreak\lean{zero} above, which uses \lean{pbw.affineVacuum};
\item \lean{VHat.}\allowbreak\lean{DLM3.}\allowbreak\lean{voaPBWBottom}\allowbreak\lean{Injective} in
  \leanpath{VOA/\allowbreak Modules/\allowbreak ZhuVermaAnArchDLM3PBWAxiom.lean}, which uses
  \lean{pbw.induced}\allowbreak\lean{Bottom}:
\end{enumerate}
File \leanpath{VOA/Modules/ZhuVermaAnArchDLM3PBWAxiom.lean}.
\begin{leancode}
theorem voaPBWBottomInjective (P : ParabolicRep W n U) :
    Function.Injective (indBottom P) :=
  VOA.PBWPrinciple.inducedBottom
    (VOA.pbw : VOA.PBWPrinciple.{u}) (Pn W n) P.act
\end{leancode}

\begin{dictionary}
Here $W$ is a vertex operator algebra and the names come from Zhu's theory of modules of $W$
(Chapters~\ref{ch:L12}--\ref{ch:L13}). The modes $a_{(n)}$ of the states of $W$ span a Lie algebra $\hat W$
(\lean{VHat}, Chapter~\ref{ch:L13}, Definition~\ref{l13:def:vhat}); \lean{Pn W n} is its ``parabolic'' Lie subalgebra of
modes of degree $\ge -n$ (Definition~\ref{l13:def:parabolicrep}); a \lean{ParabolicRep W n U} is a representation
\lean{P.act} of \lean{Pn W n} on a vector space $U$ (the would-be lowest-weight space); \lean{IndMod P} is the
induced module $U(\hat W)\otimes_{U(P_n)}U$, the ``generalised Verma module'', and \lean{indBottom P} is
$u\mapsto[1\otimes u]$ (Definition~\ref{l13:def:indmod}). A \lean{BottomProjection} is a projection of the induced
module back onto this bottom copy of $U$, compatible with the structure, and \lean{PBWFreeness} is the
statement that the levels of the induced module are free in the PBW sense
(Definition~\ref{l13:def:bottomprojection}); both follow from injectivity of \lean{indBottom}, and the datum
\lean{Pr} of Zhu's theorems is a \lean{BottomProjection} for the regular module. The hypotheses
\lean{hC2}, \lean{hss}, \lean{hcard}, \lean{hcw}, \lean{hsep}, \lean{hSR}, \lean{hFin}, \lean{hZhu} below are
explained, with their status, in Section~\ref{l02:sec:assumes}, item~5; the same file proves by \lean{rfl} that
\lean{PBWInduction.bottom (Pn W n) P.act} is \emph{definitionally} the map \lean{indBottom P} of the
Zhu chain (\lean{pbwInduction_}\allowbreak\lean{bottom_eq}). So the generic field (B) is applied, with no translation
step that could hide a mismatch. The meaning of \lean{ParabolicRep}, \lean{indBottom} and of the
consequences \lean{voaPBWFreeness} and \lean{voaPBWBottom}\allowbreak\lean{Projection} derived in the same file is the
subject of Chapter~\ref{ch:L13}.
\end{dictionary}

\begin{status}
\lean{VOA.}\allowbreak\lean{VHat.}\allowbreak\lean{DLM3.}\allowbreak\lean{voaPBWBottom}\allowbreak\lean{Injective}: \lean{[propext, Classical.choice, Quot.sound]}
(run for this chapter at an earlier commit; not re-run for this edition). Unconditional; its proof term uses the theorem \lean{VOA.pbw} (field \lean{inducedBottom}).
\emph{History:} in an earlier run, while \lean{pbw} was an axiom, it reported \lean{[propext, Classical.choice, Quot.sound, pbw]}. The file
name \leanpath{ZhuVermaAnArchDLM3PBWAxiom.lean} and its closing comment (``visibly depend on the sole
sanctioned project axiom'') date from that time and are now out of date.
\end{status}

\subsection{What used to inherit the axiom, and the twin policy}

Everything that depends on \lean{VOA.pbw} depends on one of the two declarations above. While
\lean{pbw} was an axiom, this delimited a family of results whose \texttt{\#print axioms} report contained
\lean{pbw}; since \lean{pbw} became a theorem every one of them reports only the three standard axioms. The
probes run for this chapter (Section~\ref{l02:sec:verification}) confirm this for the representatives
below; Chapter~\ref{ch:15} records the library-wide audit.

\paragraph{Through \lean{affineVacuum}.} The legacy affine $\mathfrak{sl}_2$ route:
\lean{AffineVOA.}\allowbreak\lean{vacuumA_}\allowbreak\lean{ne_}\allowbreak\lean{zero} (\leanpath{VOA/\allowbreak Examples/\allowbreak AffineSl2VOA.lean}),
\lean{AffineVOA.}\allowbreak\lean{Inst.}\allowbreak\lean{closed_}\allowbreak\lean{vacuum_}\allowbreak\lean{ne_}\allowbreak\lean{zero} and \lean{voaSl2Adm_}\allowbreak\lean{vacuum_}\allowbreak\lean{ne_}\allowbreak\lean{zero}
(\leanpath{AffineSl2VOAInstance.lean}), \lean{QuasiLisseSl2.}\allowbreak\lean{vacuum_}\allowbreak\lean{simpleSl2QL_}\allowbreak\lean{ne_}\allowbreak\lean{zero} and the
Arakawa--Kawasetsu payoff \lean{QuasiLisseSl2.}\allowbreak\lean{character_}\allowbreak\lean{formalODE_}\allowbreak\lean{simpleSl2QL}
(Chapter~\ref{ch:L16}). In an earlier run all reported \lean{[propext, Classical.choice, Quot.sound,
VOA.pbw]}; in a later run, after \lean{pbw} became a theorem, \lean{vacuumA_}\allowbreak\lean{ne_}\allowbreak\lean{zero}, \lean{closed_}\allowbreak\lean{vacuum_}\allowbreak\lean{ne_}\allowbreak\lean{zero},
\lean{vacuum_}\allowbreak\lean{simpleSl2QL_}\allowbreak\lean{ne_}\allowbreak\lean{zero} and \lean{character_}\allowbreak\lean{formalODE_}\allowbreak\lean{simpleSl2QL} report
\lean{[propext, Classical.choice, Quot.sound]}. Their remaining hypotheses (if any) are those printed in
their signatures, discussed in Chapters~\ref{ch:L08} and~\ref{ch:L16}; none of them is a PBW hypothesis.

\paragraph{Through \lean{inducedBottom}.} The canonical Zhu-theory wrappers:
\lean{zhuBijective_}\allowbreak\lean{of_voaPBW}, \lean{zhuModularity_}\allowbreak\lean{pinned_}\allowbreak\lean{of_}\allowbreak\lean{C2Cofinite_}\allowbreak\lean{voaPBW},
\lean{zhuModularity_}\allowbreak\lean{strongly}\allowbreak\lean{Rational_}\allowbreak\lean{voaPBW} and the results of \leanpath{Modules/\allowbreak ZhuInduced.lean}
and \leanpath{Modularity/\allowbreak GNProp8PBW.lean} built on them (Chapters~\ref{ch:L13}
and~\ref{ch:L15}). All three named wrappers now report \lean{[propext, Classical.choice, Quot.sound]} (run
for this chapter). This does \emph{not} make all of them unconditional: \lean{zhuBijective_}\allowbreak\lean{of_voaPBW}
takes only the vertex operator algebra \lean{W} and is unconditional, but
\lean{zhuModularity_}\allowbreak\lean{pinned_}\allowbreak\lean{of_}\allowbreak\lean{C2Cofinite_}\allowbreak\lean{voaPBW} still carries
the instance hypothesis \lean{[Finite (Φ W)]} and \lean{hC2}, \lean{hss}, \lean{hcard}, \lean{hcw},
\lean{hsep} (besides the data \lean{Mfam}, \lean{h} it speaks about), and
\lean{zhuModularity_}\allowbreak\lean{strongly}\allowbreak\lean{Rational_}\allowbreak\lean{voaPBW} still carries \lean{hSR}, \lean{hFin}, \lean{hZhu}
(the citation hypotheses analysed in Chapter~\ref{ch:L15}). They are \texttt{\#print axioms}-clean,
conditional on those hypotheses. The PBW input is no longer one of their conditions, and the other inputs
are now discharged by separate theorems: \lean{[Finite (Φ W)]} by \lean{finite_Φ_of_C2Cofinite_pbw}, \lean{hss}
from \lean{AdmissiblyRational} by \lean{isSemisimpleRing_}\allowbreak\lean{A_of_}\allowbreak\lean{admissiblyRational},
\lean{hFin} by \lean{dlmFinitelyManyIrredAdmissible_of_c2Cofinite} and \lean{hZhu} by
\lean{zhuAdmissibleSemisimple_holds}. The headline form, \lean{zhuModularity_}\allowbreak\lean{canonical_of_}\allowbreak\lean{C2Cofinite_}\allowbreak\lean{admissiblyRational}, assumes only
\lean{hC2} and \lean{AdmissiblyRational}.

The axiom-free ``twins'' of these wrappers (\lean{zhuBijective}, \lean{zhuModularity_}\allowbreak\lean{pinned_}\allowbreak\lean{of_}\allowbreak\lean{C2Cofinite},
which take \lean{Pr} as an explicit hypothesis) date from the time when \lean{pbw} was an axiom; they
are explained in Section~\ref{l02:sec:history}. Their \lean{Pr} hypothesis is now always supplied by
\lean{voaPBWBottom}\allowbreak\lean{Projection}.

\subsection{Results proved without \texorpdfstring{\texttt{VOA.pbw}}{VOA.pbw}}

Every result of the library now reports only the three standard axioms, so the distinction drawn in
this subsection is no longer one of trust. It is still informative about \emph{method}: the following
headline results were proved, while PBW was an axiom, by routes that do not pass through \lean{VOA.pbw} at
all (they reported exactly \lean{[propext, Classical.choice, Quot.sound]} in an earlier run while PBW was an axiom and still did
in a later run, after \lean{pbw} became a theorem, run for this chapter; not re-run for this edition):
\begin{itemize}
\item the Moonshine construction, \lean{VOA.}\allowbreak\lean{Moonshine.}\allowbreak\lean{moonshine} (Chapters~\ref{ch:02}--\ref{ch:14});
\item graded first-order rigidity of the Virasoro minimal models, \lean{VOA.}\allowbreak\lean{VirModule.}\allowbreak\lean{minimal_}\allowbreak\lean{graded_}\allowbreak\lean{rigidity}
  (Chapter~\ref{ch:L07}), and the Virasoro vacuum nonvanishing \lean{VOA.vacVec_}\allowbreak\lean{ne_zero};
\item uniform graded first-order rigidity of affine $\mathfrak{sl}_2$ at negative admissible level,
  \lean{VOA.}\allowbreak\lean{AffineSl2}\allowbreak\lean{GeneralRigidity.}\allowbreak\lean{negative_}\allowbreak\lean{admissible_}\allowbreak\lean{graded_}\allowbreak\lean{rigidity} (Chapter~\ref{ch:L09});
\item non-simplicity of $V^{-8/5}(\mathfrak{sl}_2)$, \lean{VOA.}\allowbreak\lean{Sl2NotSimple85.}\allowbreak\lean{not_}\allowbreak\lean{isSimple_}\allowbreak\lean{85}, and the
  affine $\mathfrak{sl}_2$ VOA itself, \lean{VOA.}\allowbreak\lean{AffineVOA.}\allowbreak\lean{affineSl2VOA_}\allowbreak\lean{closed} (Chapter~\ref{ch:L08});
\item the affine $\mathfrak{sl}_3$ simple VOA \lean{VOA.}\allowbreak\lean{AffineSl3.}\allowbreak\lean{simpleVOA}, its vacuum nonvanishing
  \lean{VOA.}\allowbreak\lean{AffineSl3.}\allowbreak\lean{vacuumA_}\allowbreak\lean{ne_}\allowbreak\lean{zero} (proved by a free-field evaluation), and
  \lean{VOA.}\allowbreak\lean{AffineSl3.}\allowbreak\lean{simple_}\allowbreak\lean{regular_}\allowbreak\lean{isTrivial} (declarations of the
  bespoke affine $\mathfrak{sl}_3$ files under \leanpath{VOA/\allowbreak Examples/}; the chapter documenting them is in
  preparation, so they are not refereed in this edition);
\item ordinary PBW for affine $\mathfrak{sl}_2$ at noncritical level,
  \lean{VOA.}\allowbreak\lean{AffineSl2}\allowbreak\lean{Wakimoto}\allowbreak\lean{LiLeading.}\allowbreak\lean{affineSl2}\allowbreak\lean{PBWBasis}, and the resulting \lean{pbw}-independent vacuum
  nonvanishing \lean{VOA.}\allowbreak\lean{AffineSl2}\allowbreak\lean{Boundary}\allowbreak\lean{MaximalIdeal.}\allowbreak\lean{affineSl2_}\allowbreak\lean{vacuum_}\allowbreak\lean{ne_}\allowbreak\lean{zero} and
  \lean{VOA.}\allowbreak\lean{Wak.}\allowbreak\lean{vacuumA_}\allowbreak\lean{ne_}\allowbreak\lean{zero_}\allowbreak\lean{free} (Section~\ref{l02:sec:pbwbasis} and Chapter~\ref{ch:L08}).
\end{itemize}

File \leanpath{VOA/Examples/AffineSl2BoundaryMaximalIdeal.lean}.
\begin{leancode}
theorem affineSl2_vacuum_ne_zero : (W k hk).vacuum ≠ 0 := by
  change AffineLie.vacuumVector kappaForm kappaForm_isInvariant k ≠ 0
  exact (AffineSl2WakimotoLiLeading.affineSl2PBWBasis k hk).vacuumVector_ne_zero
\end{leancode}
File \leanpath{VOA/Examples/WakimotoFree.lean}.
\begin{leancode}
theorem vacuumA_ne_zero_free (k : ℂ) (hk : k + 2 ≠ 0) : AffineVOA.vacuumA k ≠ 0 := by
\end{leancode}

\begin{dictionary}
These are the \emph{same} statement as an instance of \lean{AffineLie.}\allowbreak\lean{vacuumVector_}\allowbreak\lean{ne_}\allowbreak\lean{zero}
($\ltwog = \mathfrak{sl}_2$, $\kappa$ the normalised form \lean{kappaForm} with $\kappa(e,f) = 1$,
$\kappa(h,h) = 2$, \leanpath{Examples/\allowbreak AffineSl2.lean}), restricted to $k\ne-2$, but proved without
\lean{VOA.pbw}: the first from the Wakimoto-based PBW basis of Section~\ref{l02:sec:pbwbasis}, the second by
mapping $V^k(\mathfrak{sl}_2)$ to the Wakimoto free-field module (the same strategy as for Virasoro in
Section~\ref{l02:sec:virmod}). The critical level $k=-2$ is not covered by these routes; there the
library proves $\ltwovac\ne0$ through the theorem \lean{VOA.pbw} (Theorem~\ref{l02:thm:vacind}), and the
general \lean{pbwBasis_holds} of Section~\ref{l02:sec:pbwbasis} covers it as well. When \lean{pbw} was an
axiom, the critical level was the one place where $\mathfrak{sl}_2$ vacuum nonvanishing rested on it.
\end{dictionary}

\begin{warning}
\textbf{Documentation lags about PBW (Lean docstrings, not statements).}
\begin{itemize}
\item The docstring of \leanpath{Examples/\allowbreak AffineSl2VOAInstance.lean} says that ``a \emph{zero} module
  satisfies every VOA axiom, so no \lean{pbw}-free theorem can assert that this VOA is nonzero''. The first
  half is right and important (the VOA axioms are satisfied by the zero space, so nonvanishing is extra
  information; the release audit of 2026-10-27 confirms, finding m-4, that the zero space satisfies every field of
  \lean{VertexOperatorAlgebra}, including \lean{vacuum_grade}, since $\C\cdot 0 = 0$). The second half was already superseded in the first edition of this chapter
  (\lean{affineSl2_}\allowbreak\lean{vacuum_}\allowbreak\lean{ne_}\allowbreak\lean{zero} and \lean{vacuumA_}\allowbreak\lean{ne_zero_}\allowbreak\lean{free} prove
  that nonvanishing without \lean{pbw} for $k\ne-2$), and since \lean{pbw} became a theorem the whole distinction it draws
  is moot, \lean{pbw} being a theorem.
\item The module docstring of \leanpath{Foundations/\allowbreak PBWBasis.lean} still calls \lean{VOA.pbw} ``the
  project's sole sanctioned axiom'', says ``\lean{VOA.pbw} is left exactly as it is'' and, under Status,
  ``This file proves nothing about non-degenerate $\ltwog$; \lean{PBWBasis} is a hypothesis here''; the
  docstring of \lean{PBWBasis.}\allowbreak\lean{vacuumVector_}\allowbreak\lean{ne_zero} calls \lean{PBW} ``the first field of the
  sanctioned axiom''. The last Status sentence is literally true of that one file, but \lean{PBWBasis} is
  now proved for every $\ltwog$ by \lean{pbwBasis_holds} in \leanpath{PBWTheorem/\allowbreak VacuumBasis.lean}
  (Section~\ref{l02:sec:pbwbasis}), and \lean{pbw} is a theorem.
\item The closing comment of \leanpath{Modules/\allowbreak ZhuVermaAnArchDLM3PBWAxiom.lean} (``The theorem-level
  wrappers visibly depend on the sole sanctioned project axiom'') is out of date for the same reason.
\end{itemize}
None of these affects a Lean statement. This list is a sample, not an exhaustive inventory: comments in
the library that describe \lean{VOA.pbw} as an axiom are out of date wherever they occur (for instance in
\leanpath{VOA/Modules.lean} and in several \leanpath{VOA/Classification/COne/} files, whose comments still
report ``Lean's three plus \lean{VOA.pbw}''); the book-side occurrences are listed in
\texttt{Notes/Textbook/referee/pbw\_sweep.md}.
\end{warning}

\begin{physics}
The pattern is general and physically natural: a universal vacuum module is nonzero as soon as one
exhibits a concrete theory in which the vacuum conditions hold --- a free-field realisation (free
bosons for Virasoro, Wakimoto's $\beta\gamma$ system plus a boson for $\widehat{\mathfrak{sl}}_2$, analogous
free fields for $\widehat{\mathfrak{sl}}_3$). The relative PBW theorem of Section~\ref{l02:sec:pbwproof} is the
same idea made universal: the normal-ordered Fock space $S(\mathfrak c)\otimes U$ is a ``free-field
realisation'' that exists for every Lie algebra, every level (including the critical one) and every
induced module of Zhu's theory.
\end{physics}

%======================================================================
\section{PBW bases: \texorpdfstring{\texttt{PBWBasis}}{PBWBasis}, from named hypothesis to theorem}\label{l02:sec:pbwbasis}
%======================================================================

All declarations of this section are in \leanpath{VOA/\allowbreak Foundations/\allowbreak PBWBasis.lean}
unless stated otherwise; the general theorem \lean{pbwBasis_holds} is in
\leanpath{VOA/\allowbreak Foundations/\allowbreak PBWTheorem/\allowbreak VacuumBasis.lean}.

\subsection{Why a second PBW statement}

Neither field of \lean{PBWPrinciple} gives \emph{linear independence} of creation monomials. Statements
such as ``the $C_2$-algebra of $V^k(\ltwog)$ is the symmetric algebra $S(\ltwog)$'' need independence: one
must be able to read off the coefficient of a particular monomial in an arbitrary vector. The library
therefore has a second PBW statement, \lean{PBWBasis}: the ordered creation monomials form a basis of the
vacuum module. It was introduced (decision of 23 August 2026, recorded in the file's docstring) under the
policy that a PBW input may be \emph{assumed}, but only as a \emph{named hypothesis} in the signature of
every theorem that uses it, never as a further axiom: \texttt{\#print axioms} then remains a reliable gate
distinguishing unconditional from conditional results; a hypothesis is visible in the signature, where
every previous vacuity failure of the project was found; and a hypothesis claims less (``if PBW, then
\ldots''). Since \lean{pbw} became a theorem that hypothesis is a \emph{theorem} for every $\ltwog$, every
invariant form, every ordered basis and every level (\lean{pbwBasis_holds},
Section~\ref{l02:sec:pbwbasisholds}). The definition and the consequences below are unchanged; what has
changed is that every hypothesis \lean{h : PBWBasis κ hκ b k} in the library can now be discharged.

\begin{definition}[Ordered creation word; PBW basis hypothesis]\label{l02:def:pbwbasis}
Let $(b_i)_{i\in I}$ be a basis of $\ltwog$ indexed by a totally ordered set $I$, and order $I\times\N$
lexicographically. For a finitely supported exponent function $\mu : I\times\N\to\N$, the
\emph{ordered creation word} lists the pairs $(i,n)$ in the support in increasing order, each repeated
$\mu(i,n)$ times; the corresponding \emph{ordered creation monomial} is
\begin{equation*}
 v_\mu = \prod_{(i,n)\ \text{increasing}}(b_i\otimes t^{-n-1})^{\mu(i,n)}\ltwovac\ \in V^k(\ltwog),
\end{equation*}
with the smallest pair leftmost. \emph{$\mathrm{PBWBasis}(\kappa,b,k)$} is the statement that
$(v_\mu)_\mu$ is a basis of $V^k(\ltwog)$.
\end{definition}

File \leanpath{VOA/Foundations/PBWBasis.lean}.
\begin{leancode}
noncomputable def pbwWord {𝔤 : Type u} [LieRing 𝔤] [LieAlgebra ℂ 𝔤] {ι : Type*} [LinearOrder ι]
    (b : Module.Basis ι ℂ 𝔤) (μ : (ι ×ₗ ℕ) →₀ ℕ) : List (𝔤 × ℕ) :=
  (μ.support.sort (· ≤ ·)).flatMap fun p => List.replicate (μ p) (b (ofLex p).1, (ofLex p).2)

def PBWBasis {𝔤 : Type u} [LieRing 𝔤] [LieAlgebra ℂ 𝔤]
    (κ : LinearMap.BilinForm ℂ 𝔤) (hκ : κ.IsInvariant)
    {ι : Type*} [LinearOrder ι] (b : Module.Basis ι ℂ 𝔤) (k : ℂ) : Prop :=
  ∃ B : Module.Basis ((ι ×ₗ ℕ) →₀ ℕ) ℂ (AffineLie.vacuumSpace κ hκ k),
    ∀ μ, B μ = AffineLie.negWord κ hκ k (pbwWord b μ)
\end{leancode}
and its first consequences:
\begin{leancode}
theorem linearIndependent (h : PBWBasis κ hκ b k) :
    LinearIndependent ℂ fun μ : (ι ×ₗ ℕ) →₀ ℕ => AffineLie.negWord κ hκ k (pbwWord b μ) := by

theorem vacuumVector_ne_zero (h : PBWBasis κ hκ b k) :
    AffineLie.vacuumVector κ hκ k ≠ 0 := by

theorem pbwBasis_iff_of_subsingleton :
    PBWBasis κ hκ (Module.Basis.empty 𝔤 : Module.Basis (Fin 0) ℂ 𝔤) k
      ↔ AffineLie.vacuumVector κ hκ k ≠ 0 :=
\end{leancode}

\begin{dictionary}
\textbf{Decoding.} \lean{Module.Basis ι ℂ 𝔤} is a basis of $\ltwog$ indexed by $\iota$ (a linear
equivalence with \lean{ι →₀ ℂ}); \lean{[LinearOrder ι]} orders the index set. \lean{ι ×ₗ ℕ} is the
product with the \emph{lexicographic} order (\lean{×ₗ} is Mathlib's \lean{Lex} synonym, and \lean{ofLex}
forgets the order). \lean{(ι ×ₗ ℕ) →₀ ℕ} are the exponent functions $\mu$. In \lean{pbwWord},
\lean{μ.support.sort (· ≤ ·)} lists the support in increasing order, and each $p = (i,n)$ contributes
\lean{μ p} copies of the pair $(b_i,n)$; recall from Section~\ref{l02:sec:grading} that $(x,n)$ stands for
$x\otimes t^{-n-1}$ and that the head of the list is the leftmost operator. \lean{PBWBasis κ hκ b k}
asserts that there is a basis \lean{B} of $V^k(\ltwog)$ indexed by exponent functions whose $\mu$-th
vector is the ordered creation monomial $v_\mu$. This is Definition~\ref{l02:def:pbwbasis}.

\textbf{Hypotheses of the consequences.} The theorems take \lean{h : PBWBasis κ hκ b k} as an explicit
hypothesis; nothing in this file is an axiom, and nothing in it uses \lean{VOA.pbw}. \lean{linearIndependent} extracts independence;
\lean{span_eq_top} (not displayed) extracts spanning by \emph{ordered} words, a strengthening of
\lean{negSpan_eq_top}, which spans by \emph{all} words; \lean{vacuumVector_}\allowbreak\lean{ne_zero} shows
$\ltwovac = v_0\ne0$, i.e.\ \lean{PBWBasis} implies field (A) of \lean{PBWPrinciple} at the given
$(\ltwog,\kappa,k)$.

\textbf{Is \lean{PBWBasis} true?} Yes, for every $\ltwog$, $\kappa$, ordered basis $b$ and $k$. By
Theorem~\ref{l02:thm:A}'s proof, $V^k(\ltwog)\cong U(\ltwon_-)$ via $u\mapsto u\ltwovac$ (the map
$\Phi$ there is inverse to it), and $\{b_i\otimes t^{-n-1}\}_{(i,n)}$ is a basis of $\ltwon_-$, ordered
lexicographically; PBW for $\ltwon_-$ (Theorem~\ref{l02:thm:pbw}) says the ordered monomials in this
basis form a basis of $U(\ltwon_-)$. Any total order would do; the lexicographic one is the library's
choice. So \lean{PBWBasis} is a true hypothesis, never a vacuous one --- and since \lean{pbw} became a theorem the
library proves it (Section~\ref{l02:sec:pbwbasisholds}). It is \emph{logically stronger} than field (A)
(it implies (A) for its instance, and asserts a whole basis), which is the sense of the docstring's
``strictly stronger''; both are now theorems.

\textbf{Positive control.} A hypothesis of the form ``there is a basis such that\ldots'' could be
unsatisfiable through a badly chosen index type (two indices forced to the same vector). The theorem
\lean{pbwBasis_}\allowbreak\lean{iff_}\allowbreak\lean{of_}\allowbreak\lean{subsingleton} shows that for the zero Lie algebra (with the empty basis, so a
single exponent function) \lean{PBWBasis} holds \emph{exactly} when $\ltwovac\ne0$; so the shape is
satisfiable whenever the module is nonzero. It is a sanity check of the shape, not a proof of PBW.
\end{dictionary}

\begin{physics}
\lean{PBWBasis} is the precise form of the physicist's counting of states in the vacuum sector of a
current algebra: a basis is given by the normal-ordered products of the creation modes
$J^{i}_{-n-1}$, ordered first by colour $i$ and then by depth. Its graded dimension is
$\prod_{n\ge1}(1-q^n)^{-\dim\ltwog}$.
\end{physics}

\begin{proof}[Proof on paper of \lean{pbwBasis_}\allowbreak\lean{iff_}\allowbreak\lean{of_}\allowbreak\lean{subsingleton}]
If $\ltwog = 0$, every nonempty creation word is $0$ (each mode is $0\otimes t^{-n-1} = 0$), so
$V^k(\ltwog) = \C\ltwovac$ by spanning; there is a single exponent function $\mu = 0$, with $v_0 = \ltwovac$.
A single vector is a basis of the space it spans iff it is nonzero.
\end{proof}

\begin{status}
\lean{VOA.}\allowbreak\lean{PBWBasis.}\allowbreak\lean{linearIndependent}, \lean{VOA.}\allowbreak\lean{PBWBasis.}\allowbreak\lean{vacuumVector_}\allowbreak\lean{ne_}\allowbreak\lean{zero},
\lean{VOA.}\allowbreak\lean{pbwBasis_}\allowbreak\lean{iff_}\allowbreak\lean{of_}\allowbreak\lean{subsingleton}: \lean{[propext, Classical.choice, Quot.sound]} (run for this
chapter). Unconditional as implications. \lean{PBWBasis} itself is a definition (a proposition), proved
for all data by \lean{pbwBasis_holds}; a library theorem that takes \lean{h : PBWBasis κ hκ b k} as a
hypothesis is therefore unconditional in substance, though its Lean statement still displays the (now
redundant) hypothesis.
\end{status}

\subsection{The hypothesis discharged for \texorpdfstring{$\widehat{\mathfrak{sl}}_2$}{affine sl2}}

Before the general theorem of Section~\ref{l02:sec:pbwbasisholds}, the hypothesis was first discharged
for $\widehat{\mathfrak{sl}}_2$ at noncritical level, by a free-field argument that is still in the library
and still consumed:

File \leanpath{VOA/Examples/AffineSl2WakimotoLiLeading.lean}.
\begin{leancode}
theorem affineSl2PBWBasis (k : ℂ) (hk : k + 2 ≠ 0) :
    PBWBasis AffineVOA.kappaForm AffineVOA.kappaForm_isInvariant
      (Pi.basisFun ℂ (Fin 3)) k :=
  pbwBasis_of_wakimotoOrderedCreationState_linearIndependent k hk
    (wakimotoOrderedCreationState_linearIndependent k)
\end{leancode}
and it is consumed, e.g., by \lean{universal}\allowbreak\lean{SymbolAlg_}\allowbreak\lean{injective} in the same file (injectivity of the
$C_2$ symbol map of $V^k(\mathfrak{sl}_2)$) and by \leanpath{Examples/\allowbreak AffineSl2JointPBW.lean}.

\begin{dictionary}
$\mathfrak{sl}_2$ is modelled in the library as \lean{Fin 3 → ℂ} with the bracket given by structure
constants in the basis $(e,f,h)$ (indices $0,1,2$; \leanpath{Examples/\allowbreak AffineSl2.lean}), so
\lean{Pi.basisFun ℂ (Fin 3)} is the basis $(e,f,h)$ with the order $e<f<h$ of \lean{Fin 3}.
\lean{kappaForm} is the normalised invariant form. The hypothesis \lean{hk : k + 2 ≠ 0} excludes the
critical level $k = -h^\vee = -2$. The theorem is the instance of Definition~\ref{l02:def:pbwbasis} for
$\mathfrak{sl}_2$ at every noncritical level, proved without any hypothesis. Its proof (detailed in
Chapter~\ref{ch:L08}) maps $V^k(\mathfrak{sl}_2)$ to the Wakimoto free-field module and shows that the images
of the ordered creation monomials are linearly independent by a leading-term (triangularity) argument;
the noncriticality is used by the Wakimoto realisation, not by PBW, which holds at $k=-2$ too.
\end{dictionary}

\begin{physics}
The Wakimoto realisation writes the $\widehat{\mathfrak{sl}}_2$ currents in terms of a $\beta\gamma$ system and
a free boson with background charge depending on $k+2$. Since a free-field Fock space has a manifest
basis, independence of the images of the PBW monomials can be checked there; that independence then
holds in $V^k(\mathfrak{sl}_2)$ itself.
\end{physics}

\begin{status}
\lean{VOA.}\allowbreak\lean{AffineSl2}\allowbreak\lean{Wakimoto}\allowbreak\lean{LiLeading.}\allowbreak\lean{affineSl2}\allowbreak\lean{PBWBasis}: \lean{[propext, Classical.choice,
Quot.sound]} (run for this chapter). Unconditional for $k\ne-2$; does not use \lean{VOA.pbw}. It is now a
special case of \lean{pbwBasis_holds} below (which also covers $k = -2$ and every ordered basis), but it
is a separate proof, and the files that consume it were not rewired.
\end{status}

\subsection{The hypothesis proved for every \texorpdfstring{$\mathfrak g$}{g}}\label{l02:sec:pbwbasisholds}

\begin{theorem}[PBW basis of the affine vacuum module]\label{l02:thm:pbwbasisholds}
For every complex Lie algebra $\ltwog$ (of any dimension), every symmetric invariant bilinear form $\kappa$,
every basis $(b_i)_{i\in I}$ of $\ltwog$ indexed by a totally ordered set, and every $k\in\C$ (including the
critical level), $\mathrm{PBWBasis}(\kappa,b,k)$ holds: the ordered creation monomials $v_\mu$ of
Definition~\ref{l02:def:pbwbasis} form a basis of $V^k(\ltwog)$. More precisely, the linear map
$\Phi : V^k(\ltwog)\to S(\ltwon_-)$, $[a]\mapsto a\cdot1$, given by the relative PBW representation of
$\ltwoghat$ on $S(\ltwon_-)\otimes\C_{\chi_k}\cong S(\ltwon_-)$, is a linear isomorphism.
\end{theorem}

File \leanpath{VOA/Foundations/PBWTheorem/VacuumBasis.lean}.
\begin{leancode}
noncomputable def yMode (a : ι ×ₗ ℕ) : AffineLie 𝔤 κ hκ :=
  ofLoop κ hκ (b (ofLex a).1) (-((ofLex a).2 : ℤ) - 1)

noncomputable def negComp : Submodule ℂ (AffineLie 𝔤 κ hκ) :=
  Submodule.span ℂ (Set.range (yMode κ hκ b))

theorem negComp_isCompl : IsCompl (nonnegSub κ hκ).toSubmodule (negComp κ hκ b) := by

noncomputable def Φ : vacuumSpace κ hκ k →ₗ[ℂ] (Multiset (ι ×ₗ ℕ) →₀ ℂ) :=

theorem Φ_rho (z : AffineLie 𝔤 κ hκ) (v : vacuumSpace κ hκ k) :
    Φ κ hκ b k (vacuumRho κ hκ k z v) = ACT z (Φ κ hκ b k v) := by

noncomputable def ΦEquiv : vacuumSpace κ hκ k ≃ₗ[ℂ] (Multiset (ι ×ₗ ℕ) →₀ ℂ) :=
  LinearEquiv.ofLinear (Φ κ hκ b k) (Ψ κ hκ b k)
    (Finsupp.lhom_ext fun m u => Φ_Ψ_single κ hκ b k m u)
    (LinearMap.ext_on_range (negSpan_eq_top κ hκ k) fun L => Ψ_Φ_negWord κ hκ b k L)

theorem pbwBasis_holds {𝔤 : Type u} [LieRing 𝔤] [LieAlgebra ℂ 𝔤]
    (κ : LinearMap.BilinForm ℂ 𝔤) (hκ : κ.IsInvariant)
    {ι : Type*} [LinearOrder ι] (b : Module.Basis ι ℂ 𝔤) (k : ℂ) : PBWBasis κ hκ b k := by
\end{leancode}

\begin{dictionary}
\textbf{Decoding.} \lean{yMode a}, for \lean{a = (i, n)} in \lean{ι ×ₗ ℕ}, is the creation mode
$b_i\otimes t^{-n-1}$ (depth $n$ means mode $-n-1$, as in Section~\ref{l02:sec:grading}). \lean{negComp} is
their span, i.e.\ $\ltwon_- = \ltwog\otimes t^{-1}\C[t^{-1}]$, and \lean{negComp_isCompl} says it is a vector-space
complement of \lean{nonnegSub} $=\hat{\ltwop}$ in $\ltwoghat$; the omitted \lean{negBasis} is the basis
$(y_{(i,n)})$ of this complement, indexed by the lexicographically ordered \lean{ι ×ₗ ℕ}. With these choices
and the character \lean{vacuumChar k} on the one-dimensional $U = \C$, the relative PBW representation of
Theorem~\ref{l02:thm:relrep} acts on \lean{Multiset (ι ×ₗ ℕ) →₀ ℂ}, which is $S(\ltwon_-)$ in its monomial
basis; \lean{ACT} is local notation for that action. \lean{Φ} sends the class of $a\in U(\ltwoghat)$ to $a$
applied to the vector $\varnothing\otimes1$ (the omitted body checks that $J(k)$ acts by zero on it);
\lean{Φ_rho} says $\Phi$ is $\ltwoghat$-equivariant. \lean{ΦEquiv} upgrades $\Phi$ to a linear equivalence
(\lean{≃ₗ[ℂ]}) with inverse \lean{Ψ}, which sends a monomial $m$ to the sorted creation monomial of $m$
applied to $\ltwovac$. \lean{pbwBasis_holds} is the theorem, with \emph{no} hypotheses beyond the data; its
statement is literally \lean{PBWBasis κ hκ b k} of Definition~\ref{l02:def:pbwbasis}.

\textbf{Same as the textbook?} Yes, and in the strongest form: any $\ltwog$ (no finiteness), any invariant
$\kappa$ (possibly degenerate), any ordered basis, any $k$. The order on \lean{ι ×ₗ ℕ} used to sort is the
same lexicographic order in which \lean{pbwWord} lists letters, so the basis vectors are exactly the
$v_\mu$ of the definition.

\textbf{Could it be vacuous?} No: the conclusion is the existence of a basis \emph{equal} to the given
family of creation monomials, which pins every basis vector; and $\Phi(\ltwovac) = \varnothing\otimes1\ne0$
(\lean{Φ_vacuum}).
\end{dictionary}

\begin{proof}[Proof on paper]
Write $\rho_\Phi$ for the relative PBW representation of $\ltwoghat$ on $S(\ltwon_-) = \bigoplus_m\C\,m$ built from
$\hat{\ltwop}$, its complement $\ltwon_-$ with the ordered basis $y_{(i,n)}$, and the character $\chi_k$ on $\C$
(Theorem~\ref{l02:thm:relrep}). Extend it to $U(\ltwoghat)$ and set $\Phi([a]) = a\cdot\varnothing$. This is
well defined: each generator $X - \chi_k(X)$ of $J(k)$ ($X\in\hat{\ltwop}$) kills $\varnothing$ by (E3), hence so
does the left ideal $J(k)$. $\Phi$ is $\ltwoghat$-equivariant by construction. By (E1), applying an ordered
creation monomial $y_{a_1}\cdots y_{a_r}$ with $a_1\le\dots\le a_r$ to $\varnothing$ gives the monomial
$a_1\cdots a_r$; so $\Phi(v_\mu) = $ the basis monomial of $S(\ltwon_-)$ with exponents $\mu$.

Conversely define $\Psi(m) = $ (the sorted creation monomial of $m$)$\cdot\ltwovac$, i.e.\ for
$m = a_1\cdots a_r$ with $a_1\le\dots\le a_r$, $\Psi(m) = y_{a_1}\cdots y_{a_r}\ltwovac$. We show
$\Psi(\sigma(z)(m\otimes1)) = z\cdot\Psi(m)$ for every generator $z$ of $\ltwoghat$ (a $y_a$ or an element of
$\hat{\ltwop}$) by induction on $|m|$, following the four equations. (E1): if $a\le m$, then
$\Psi((a\,m)\otimes1) = y_a\,y_{a_1}\cdots y_{a_r}\ltwovac = y_a\Psi(m)$ by definition of the sorted monomial. (E3): for
$X\in\hat{\ltwop}$, $X\ltwovac = \chi_k(X)\ltwovac$ because $X - \chi_k(X)\in J(k)$, which matches
$\sigma(X)(\varnothing\otimes1) = \chi_k(X)\varnothing\otimes1$. (E2): for $b<a$, $b\le m$, the left side is
$\Psi$ of the right side of (E2), which by induction (the vectors involved have degree $\le|m|$ and
the $y_b$ applied last is low) equals $y_b\,y_a\Psi(m) + [y_a,y_b]\Psi(m)$; in $U(\ltwoghat)$ we have
$y_ay_b = y_by_a + [y_a,y_b]$, so this is $y_a\,y_b\Psi(m) = y_a\Psi((b\,m)\otimes1)$. (E4) is the same computation with
$x\in\hat{\ltwop}$ in place of $y_a$, using $xy_b = y_bx + [x,y_b]$. Since every $\sigma(z)$ is determined by these
equations, $\Psi$ intertwines $\sigma$ with the action of $\ltwoghat$ on $V^k(\ltwog)$. Then
$\Phi\circ\Psi = \mathrm{id}$ on monomials (by the previous paragraph), and $\Psi\circ\Phi = \mathrm{id}$ on every
creation word $w\ltwovac$: writing $w = z_1\cdots z_r$, equivariance of $\Phi$ and of $\Psi$ gives
$\Psi\Phi(z_1\cdots z_r\ltwovac) = z_1\cdots z_r\,\Psi\Phi(\ltwovac)$, and $\Psi\Phi(\ltwovac) = \Psi(\varnothing) = \ltwovac$, hence on all of $V^k(\ltwog)$ by the spanning theorem
\lean{negSpan_eq_top} (Theorem~\ref{l02:thm:span}). So $\Phi$ is an isomorphism, and the pull-back of the
monomial basis of $S(\ltwon_-)$ is the family $(v_\mu)$, which is therefore a basis.
\end{proof}

\begin{physics}
This is the precise statement behind the physicist's character formula
$\tr_{V^k(\ltwog)}q^{L_0} = \prod_{n\ge1}(1-q^n)^{-\dim\ltwog}$ (for finite-dimensional $\ltwog$, up to the
prefactor $q^{-c/24}$): the vacuum sector of a current algebra is a Fock space of $\dim\ltwog$ bosonic
oscillators per mode number, with no relations among normal-ordered states, at every level including the
critical one. The map $\Phi$ is the identification of the vacuum sector with that Fock space.
\end{physics}

\begin{status}
\lean{VOA.pbwBasis_holds}, \lean{VOA.}\allowbreak\lean{VacuumBasis.}\allowbreak\lean{ΦEquiv}: \lean{[propext, Classical.choice,
Quot.sound]} (run for this chapter, Section~\ref{l02:sec:verification}). Unconditional. It makes every
\lean{h : PBWBasis …} hypothesis in the library dischargeable (the project lists removing such now-redundant
hypotheses, e.g.\ in \leanpath{Examples/\allowbreak AffineSl2UniversalC2PBW.lean}, as optional clean-up, not done
when checked at an earlier commit; not re-checked for this edition).
\end{status}

%======================================================================
\section{Root data, affine Verma modules and the integrability target}\label{l02:sec:roots}
%======================================================================

All declarations of this section are in \leanpath{VOA/\allowbreak Foundations/\allowbreak RootDatum.lean}. The file
is scaffolding for the representation theory of affine algebras at general highest weight; little of
it is consumed elsewhere, and its main theorems are either thin wrappers around Mathlib's root theory or
explicitly conditional.

\subsection{The finite root datum}

\begin{definition}\label{l02:def:roots}
Let $\ltwog$ be finite-dimensional with nondegenerate Killing form, and $\mathfrak h\subseteq\ltwog$ a
splitting Cartan subalgebra. The roots of $(\ltwog,\mathfrak h)$ form a reduced crystallographic root
system $\Phi$; it has a base (system of simple roots); for each root $\alpha$ there are $e_\alpha\in\ltwog_\alpha$,
$f_\alpha\in\ltwog_{-\alpha}$ such that $(\alpha^\vee, e_\alpha, f_\alpha)$ is an $\mathfrak{sl}_2$-triple, $\alpha^\vee$ the
coroot; $\beta(\alpha^\vee)\in\Z$ for roots $\alpha,\beta$; a \emph{highest root} $\theta$ is a positive root of
maximal height; the \emph{affine node} is $e_0 = f_\theta\otimes t$, $f_0 = e_\theta\otimes t^{-1}$.
\end{definition}

File \leanpath{VOA/Foundations/RootDatum.lean}.
\begin{leancode}
variable {𝔤 : Type*} [LieRing 𝔤] [LieAlgebra ℂ 𝔤] [FiniteDimensional ℂ 𝔤]
  [LieAlgebra.IsKilling ℂ 𝔤]
  (H : LieSubalgebra ℂ 𝔤) [H.IsCartanSubalgebra] [LieModule.IsTriangularizable ℂ H 𝔤]

theorem nonempty_base : Nonempty (rootSystem H).Base :=
  (rootSystem H).nonempty_base

theorem exists_chevalleyTriple (α : Weight ℂ H 𝔤) (hα : α.IsNonZero) :
    ∃ e f : 𝔤, IsSl2Triple (↑(coroot α) : 𝔤) e f ∧ e ∈ rootSpace H α ∧ f ∈ rootSpace H (-α) := by

theorem cartanInteger (α β : Weight ℂ H 𝔤) :
    ∃ z : ℤ, β (coroot α) = (z : ℂ) :=
  ⟨_, apply_coroot_eq_cast α β⟩

noncomputable def highestRootIndex [Nonempty H.root] (b : (rootSystem H).Base) : H.root :=
  (exists_highestRootIndex H b).choose

theorem height_le_highestRootIndex [Nonempty H.root] (b : (rootSystem H).Base)
    {α : H.root} (hα : b.IsPos α) : b.height α ≤ b.height (highestRootIndex H b) :=

noncomputable def affineEZero [Nonempty H.root] (b : (rootSystem H).Base) :
    VOA.AffineLie 𝔤 κ hκ :=
  VOA.AffineLie.ofLoop κ hκ (highestChevalleyTriple H b).f 1

noncomputable def affineFZero [Nonempty H.root] (b : (rootSystem H).Base) :
    VOA.AffineLie 𝔤 κ hκ :=
  VOA.AffineLie.ofLoop κ hκ (highestChevalleyTriple H b).e (-1)

theorem affineEZero_bracket_affineFZero [Nonempty H.root] (b : (rootSystem H).Base) :
    ⁅affineEZero H κ hκ b, affineFZero H κ hκ b⁆ =
      VOA.AffineLie.ofLoop κ hκ ⁅(highestChevalleyTriple H b).f,
        (highestChevalleyTriple H b).e⁆ 0 +
      (κ (highestChevalleyTriple H b).f (highestChevalleyTriple H b).e) •
        VOA.AffineLie.C κ hκ := by
\end{leancode}

\begin{dictionary}
\textbf{Hypotheses.} \lean{[LieAlgebra.IsKilling ℂ 𝔤]} says the Killing form is nondegenerate (true for
semisimple $\ltwog$); \lean{H.IsCartan}\allowbreak\lean{Subalgebra} that $H$ is a Cartan subalgebra;
\lean{IsTriangularizable} that $\mathrm{ad}\,H$ is triangularisable (automatic over $\C$). \lean{rootSystem H},
\lean{coroot}, \lean{rootSpace}, \lean{Weight ℂ H 𝔤}, \lean{IsSl2Triple}, \lean{H.root}, \lean{.Base},
\lean{b.IsPos}, \lean{b.height} are Mathlib's root-system vocabulary: the root system, the coroot
$\alpha^\vee\in H$, the root space $\ltwog_\alpha$, a weight of $H$ on $\ltwog$, an $\mathfrak{sl}_2$-triple
$[h,e] = 2e$, $[h,f] = -2f$, $[e,f] = h$, the finite type of roots, a base, positivity and height
with respect to the base. \lean{α.IsNonZero} excludes the zero weight (whose ``root space'' is $H$).

\textbf{Content.} \lean{nonempty_base}, \lean{exists_chevalley}\allowbreak\lean{Triple} and \lean{cartanInteger} are
restatements of Mathlib theorems (\lean{RootPairing.}\allowbreak\lean{nonempty_}\allowbreak\lean{base},
\lean{exists_}\allowbreak\lean{isSl2Triple_}\allowbreak\lean{of_}\allowbreak\lean{weight_}\allowbreak\lean{isNonZero}, \lean{apply_coroot_}\allowbreak\lean{eq_cast}); the library adds only the
packaging. The highest root is defined as a positive root of maximal height (a choice, via
\lean{Finset.}\allowbreak\lean{exists_}\allowbreak\lean{max_}\allowbreak\lean{image}); \lean{[Nonempty H.root]} excludes the case with no roots. For simple
$\ltwog$ the root of maximal height is unique, so this is \emph{the} highest root, but uniqueness is not
proved here. \lean{affineEZero_}\allowbreak\lean{bracket_}\allowbreak\lean{affineFZero} is \lean{affine_bracket} at $m = 1$, $n = -1$:
$[f_\theta\otimes t, e_\theta\otimes t^{-1}] = [f_\theta,e_\theta]\otimes1 + \kappa(f_\theta,e_\theta)C$. The textbook
form $[e_0,f_0] = C-\theta^\vee$ would need the normalisation $\kappa(f_\theta,e_\theta) = 1$, which is not
imposed (the docstring says so).
\end{dictionary}

\begin{physics}
The finite root datum classifies the symmetry algebra of the WZW model; the affine node $e_0,f_0$ is
the extra $\mathfrak{sl}_2$ that mixes zero modes with modes $\pm1$, and integrability under it is the
origin of the level-$k$ truncation of the WZW spectrum (only representations with $\theta(\lambda)\le k$
appear).
\end{physics}

\begin{status}
\lean{VOA.}\allowbreak\lean{RootDatum.}\allowbreak\lean{exists_}\allowbreak\lean{chevalleyTriple}, \lean{VOA.}\allowbreak\lean{RootDatum.}\allowbreak\lean{cartanInteger},
\lean{VOA.}\allowbreak\lean{RootDatum.}\allowbreak\lean{highestRoot_}\allowbreak\lean{isNonZero}, \lean{VOA.}\allowbreak\lean{RootDatum.}\allowbreak\lean{affineEZero_}\allowbreak\lean{bracket_}\allowbreak\lean{affineFZero}:
\lean{[propext, Classical.choice, Quot.sound]} (run for this chapter).
\end{status}

\subsection{Affine Verma modules at general highest weight}

\begin{definition}[Affine Verma module]\label{l02:def:verma}
An \emph{affine Cartan weight} $\lambda$ consists of a linear form $\lambda_{\mathrm{fin}}$ on $\mathfrak h$, a level
$\lambda(C)$ and an energy $\lambda(\ell_0)$. With a base fixed, the \emph{Verma module} is
$M(\lambda) = U(\ltwoghat)/I(\lambda)$, where $I(\lambda)$ is the left ideal generated by $x\otimes t^n$ ($n>0$),
the positive root vectors $x\otimes t^0$ ($x\in\ltwog_\alpha$, $\alpha>0$), $h\otimes t^0-\lambda_{\mathrm{fin}}(h)$
($h\in\mathfrak h$), $C-\lambda(C)$ and $\ell_0-\lambda(\ell_0)$; its highest-weight vector is $v_\lambda = [1]$.
\end{definition}

File \leanpath{VOA/Foundations/RootDatum.lean}.
\begin{leancode}
structure AffineCartanWeight where
  finite : H →ₗ[ℂ] ℂ
  level : ℂ
  energy : ℂ

def vermaGens : Set (UEnv 𝔤 κ hκ) :=
  {a | (∃ (x : 𝔤) (n : ℤ), 0 < n ∧ a = ι ℂ (ofLoop κ hκ x n))
      ∨ (∃ (α : H.root) (x : 𝔤), b.IsPos α ∧ x ∈ rootSpace H (α : Weight ℂ H 𝔤) ∧
          a = ι ℂ (ofLoop κ hκ x 0))
      ∨ (∃ h : H,
          a = ι ℂ (ofLoop κ hκ (h : 𝔤) 0) - algebraMap ℂ (UEnv 𝔤 κ hκ) (lam.finite h))
      ∨ a = ι ℂ (C κ hκ) - algebraMap ℂ (UEnv 𝔤 κ hκ) lam.level
      ∨ a = ι ℂ (ell0 κ hκ) - algebraMap ℂ (UEnv 𝔤 κ hκ) lam.energy}

abbrev verma : Type _ := UEnv 𝔤 κ hκ ⧸ vermaIdeal H κ hκ b lam

noncomputable def highestVector : verma H κ hκ b lam := Submodule.Quotient.mk 1

theorem verma_nontrivial_of_highestVector_ne_zero
    (h : highestVector H κ hκ b lam ≠ 0) : Nontrivial (verma H κ hκ b lam) :=

noncomputable abbrev simpleHWOf (h : highestVector H κ hκ b lam ≠ 0) : Type _ :=
  verma H κ hκ b lam ⧸ simpleHWMaxSubmoduleOf H κ hκ b lam h

theorem simpleHWOf_isSimpleModule (h : highestVector H κ hκ b lam ≠ 0) :
    IsSimpleModule (UEnv 𝔤 κ hκ) (simpleHWOf H κ hκ b lam h) := by
\end{leancode}

\begin{dictionary}
\textbf{Decoding.} The five disjuncts of \lean{vermaGens} are the five kinds of generator of
Definition~\ref{l02:def:verma}, in order; \lean{vermaIdeal} is their left span. Note \lean{0 < n} for the
loop modes (the zero modes enter only through the positive root vectors and the Cartan relations),
which is the Borel-type condition. \lean{simpleHWOf h} is defined only \emph{given} a proof
\lean{h : highestVector … ≠ 0}; every theorem about it carries this hypothesis.

\textbf{The hypothesis is open in the library, though true.} $v_\lambda\ne0$ holds for every $\lambda$: it is
field (B) of \lean{PBWPrinciple} (now the theorem \lean{PBWTheorem.}\allowbreak\lean{bottom_}\allowbreak\lean{injective}) for the Borel-type subalgebra $\hat{\mathfrak b} = \ltwon_+\oplus\mathfrak h\oplus\C C\oplus\C\ell_0$
(with $\ltwon_+$ now including the positive root spaces at mode $0$) and its one-dimensional character
$\lambda$, by the argument of Remark~\ref{l02:rem:AfromB}, which the library formalises for $\lambda = 0$
(Theorem~\ref{l02:thm:vacind}) but not for general $\lambda$; the
docstring of \lean{exists_}\allowbreak\lean{simpleHW}\allowbreak\lean{MaxSubmodule_}\allowbreak\lean{of_}\allowbreak\lean{highestVector_}\allowbreak\lean{ne_}\allowbreak\lean{zero} calls it ``the arbitrary-$\lambda$
PBW nonvanishing result not yet derived here from \lean{pbw.induced}\allowbreak\lean{Bottom}''. So the results of this
subsection are \texttt{\#print axioms}-clean \emph{conditional on} $v_\lambda\ne0$.
\end{dictionary}

\begin{physics}
$M(\lambda)$ is the space of all states of a current-algebra primary of weight $\lambda$ and its
descendants $J^{a_1}_{-n_1}\cdots v_\lambda$ before null states are removed; $L(\lambda)$, here
\lean{simpleHWOf}, is the irreducible representation, i.e.\ the conformal family of the primary in a
theory where all nulls decouple.
\end{physics}

\begin{status}
\lean{VOA.}\allowbreak\lean{AffineLie.}\allowbreak\lean{simpleHWOf_}\allowbreak\lean{isSimpleModule}, \lean{VOA.}\allowbreak\lean{AffineLie.}\allowbreak\lean{simpleHighest}\allowbreak\lean{VectorOf_}\allowbreak\lean{ne_}\allowbreak\lean{zero}:
\lean{[propext, Classical.choice, Quot.sound]} (run for this chapter), \textbf{conditional on}
\lean{highestVector H κ hκ b lam ≠ 0}, which is not discharged in the library.
\end{status}

\subsection{The integrability target}

File \leanpath{VOA/Foundations/RootDatum.lean}.
\begin{leancode}
def LocallyNilpotentEnd {V : Type*} [AddCommMonoid V] [Module ℂ V]
    (T : Module.End ℂ V) : Prop :=
  ∀ v : V, ∃ n : ℕ, (T ^ n) v = 0

def SimpleHWIntegrable [Nonempty H.root] (hnz : highestVector H κ hκ b lam ≠ 0) : Prop :=
  Integrable κ hκ (simpleHWOf H κ hκ b lam hnz) (simpleHWRhoOf H κ hκ b lam hnz)
    (chevalleyGeneratorSet H κ hκ b)

theorem simpleHWOf_integrable_of_chevalley_locallyNilpotent [Nonempty H.root]
    (hnz : highestVector H κ hκ b lam ≠ 0)
    (hKG : ∀ X ∈ chevalleyGeneratorSet H κ hκ b,
      LocallyNilpotentEnd (simpleHWRhoOf H κ hκ b lam hnz X)) :
    SimpleHWIntegrable H κ hκ b lam hnz := by
  exact hKG
\end{leancode}

\begin{warning}
\textbf{A tautological wiring lemma.} \lean{Integrable … S} is defined as ``every $X\in S$ acts locally
nilpotently'', and \lean{SimpleHW}\allowbreak\lean{Integrable} is \lean{Integrable} for the set of Chevalley generators
(all $e_\alpha,f_\alpha$ at mode $0$ and $e_0,f_0$). The theorem
\lean{simpleHWOf_}\allowbreak\lean{integrable_}\allowbreak\lean{of_}\allowbreak\lean{chevalley_}\allowbreak\lean{locallyNilpotent} takes exactly that statement as its
hypothesis \lean{hKG} and returns it (\lean{exact hKG}). It has no mathematical content; it records the
shape of the open target (the Kac--Garland theorem: $L(\lambda)$ is integrable for dominant integral
$\lambda$ at level $\lambda(C)\in\N$). Integrability of $L(\lambda)$ is \textbf{not proved} in the library. A
referee should not count this lemma as a result.
\end{warning}

\begin{physics}
Integrability is the algebraic shadow of unitarity of the WZW model: the representations of
$\ltwoghat$ that appear in the level-$k$ WZW model are exactly the integrable ones, in which every
$\mathfrak{sl}_2$-subalgebra attached to a simple root acts through finite-dimensional representations.
\end{physics}

\begin{status}
\lean{VOA.}\allowbreak\lean{AffineLie.}\allowbreak\lean{simpleHWOf_}\allowbreak\lean{integrable_}\allowbreak\lean{of_}\allowbreak\lean{chevalley_}\allowbreak\lean{locallyNilpotent}:
\lean{[propext, Classical.choice, Quot.sound]} (run for this chapter); a tautology.
\end{status}

%======================================================================
\section{Lattices, cocycles and parity}\label{l02:sec:lattices}
%======================================================================

The five remaining files supply the data from which lattice vertex algebras (and through them the
Moonshine module) are built, and the sign bookkeeping of super vertex algebras. They are presented in
full, with their use, in Chapter~\ref{ch:02} (Section~\ref{ch02:sec:voalattice}, integral lattices),
Chapter~\ref{ch:03} (Sections~\ref{ch03:sec:cocycle} and~\ref{ch03:sec:frame}, the cocycle and the
orthonormal frame) and Chapter~\ref{ch:01} (Section~\ref{ch01:sec:super}, parity). Here we give the
statements, their meaning, and their machine status, so that this chapter covers its files; all of them
report only the three standard axioms and none of them involves PBW.

\subsection{Integral lattices}

\begin{definition}[Integral lattice]\label{l02:def:lattice}
An \emph{integral lattice} is a free $\Z$-module $L$ of finite rank with a symmetric $\Z$-valued
bilinear form $\langle\cdot,\cdot\rangle$. It is \emph{even} if $\langle v,v\rangle\in2\Z$ for all $v$,
\emph{positive definite} if $\langle v,v\rangle>0$ for $v\ne0$. Its \emph{dual} is
$L^* = \{x\in L\otimes\mathbb{R} : \langle x,v\rangle\in\Z\ \forall v\in L\}$, and $L$ is
\emph{self-dual} (unimodular) if $L = L^*$.
\end{definition}

File \leanpath{VOA/Foundations/Lattices.lean}.
\begin{leancode}
structure IntegralLattice where
  carrier : Type u
  [addCommGroup : AddCommGroup carrier]
  [free : Module.Free ℤ carrier]
  [finite : Module.Finite ℤ carrier]
  form : LinearMap.BilinForm ℤ carrier
  isSymm : form.IsSymm

def Even : Prop := ∀ v : L.carrier, 2 ∣ L.form v v

def PosDef : Prop := ∀ v : L.carrier, v ≠ 0 → 0 < L.form v v

def SelfDual : Prop := ∀ x : ℝ ⊗[ℤ] L.carrier, x ∈ L.dual ↔ ∃ v : L.carrier, L.toReal v = x

theorem stdLattice_selfDual (n : ℕ) : (stdLattice n).SelfDual := by

noncomputable abbrev latticeAlgebra (L : IntegralLattice) : Type _ :=
  AddMonoidAlgebra ℂ L.carrier
\end{leancode}

\begin{dictionary}
\lean{IntegralLattice} bundles a type with its $\Z$-module structure, freeness, finite generation, the
form and its symmetry; the bracketed fields are made instances. \lean{L.dual} is defined (not shown) as
the additive subgroup of \lean{ℝ ⊗[ℤ] L.carrier} of elements whose real pairing with every
\lean{toReal v} $= 1\otimes v$ is an integer, and \lean{SelfDual} says it is exactly the image of $L$.
\textbf{One difference from the textbook}: \lean{PosDef} is positivity on \emph{integer} points; for an
integral form this is equivalent to positive definiteness over $\mathbb{R}$ (positivity on $\Q^r\setminus0$
gives semidefiniteness over $\mathbb{R}$ by density, and a real null direction of an integral form would be
a rational one, since the kernel of an integer matrix is defined over $\Q$). The library proves the
needed consequence (Theorem~\ref{l02:thm:bf} below), so nothing rests on this remark. Nondegeneracy is
not part of \lean{IntegralLattice}; it follows from \lean{PosDef} when that is assumed.
\lean{latticeAlgebra L} is the group algebra $\C[L]$ with $e^ue^v = e^{u+v}$ (untwisted).
\end{dictionary}

\begin{physics}
$L$ is the momentum (charge) lattice of $r$ chiral bosons compactified on the torus
$L\otimes\mathbb{R}/L$; evenness makes all vertex operators $e^{i\alpha\cdot X}$ have integer conformal weight
$\langle\alpha,\alpha\rangle/2$ (bosonic statistics), self-duality makes the chiral theory modular invariant on its
own (holomorphic).
\end{physics}

\begin{status}
\lean{VOA.}\allowbreak\lean{IntegralLattice.}\allowbreak\lean{le_}\allowbreak\lean{dual}, \lean{VOA.}\allowbreak\lean{stdLattice_}\allowbreak\lean{selfDual}: \lean{[propext, Classical.choice,
Quot.sound]} (run for this chapter).
\end{status}

\subsection{The FLM cocycle and the twisted group algebra}

\begin{theorem}[FLM cocycle]\label{l02:thm:cocycle}
Let $L$ have $\Z$-basis $B_1,\dots,B_n$ and put
$\mathrm{coc}(\alpha,\beta) = \sum_{j<i}\langle B_i,B_j\rangle\alpha_i\beta_j$ (coordinates in the basis) and
$\varepsilon(\alpha,\beta) = (-1)^{\mathrm{coc}(\alpha,\beta)}$. Then $\varepsilon$ is bimultiplicative, normalised
($\varepsilon(0,\cdot) = \varepsilon(\cdot,0) = 1$), a 2-cocycle,
$\varepsilon(\alpha,\beta)\varepsilon(\alpha+\beta,\gamma) = \varepsilon(\beta,\gamma)\varepsilon(\alpha,\beta+\gamma)$, and, if $L$ is
even, $\varepsilon(\alpha,\beta)\varepsilon(\beta,\alpha) = (-1)^{\langle\alpha,\beta\rangle}$. Consequently the product
$e^\alpha\star e^\beta = \varepsilon(\alpha,\beta)e^{\alpha+\beta}$ on $\C[L]$ is associative with unit $e^0$.
\end{theorem}

File \leanpath{VOA/Foundations/LatticeCocycle.lean}. The section variables are
\begin{leancode}
variable (L : IntegralLattice) {n : ℕ} (B : Basis (Fin n) ℤ L.carrier)
\end{leancode}
and, with them, the general cocycle and its properties:
\begin{leancode}
noncomputable def cocExp (α β : L.carrier) : ℤ :=
  ∑ i : Fin n, ∑ j : Fin n,
    (if j < i then L.form (B i) (B j) else 0) * B.repr α i * B.repr β j

noncomputable def cocycle (α β : L.carrier) : ℂ := (-1 : ℂ) ^ (cocExp L B α β)

theorem cocycle_add_left (α α' β : L.carrier) :
    cocycle L B (α + α') β = cocycle L B α β * cocycle L B α' β := by

theorem cocycle_add_right (α β β' : L.carrier) :
    cocycle L B α (β + β') = cocycle L B α β * cocycle L B α β' := by

theorem cocycle_zero_left (β : L.carrier) : cocycle L B 0 β = 1 := by

theorem cocycle_zero_right (α : L.carrier) : cocycle L B α 0 = 1 := by

theorem cocycle_sq (α β : L.carrier) : cocycle L B α β * cocycle L B α β = 1 := by

theorem cocycle_cocycleLaw (α β γ : L.carrier) :
    cocycle L B α β * cocycle L B (α + β) γ = cocycle L B β γ * cocycle L B α (β + γ) := by

theorem cocycle_comm (hEven : L.Even) (α β : L.carrier) :
    cocycle L B α β * cocycle L B β α = (-1 : ℂ) ^ (L.form α β) := by
\end{leancode}
and the canonical cocycle of a lattice, obtained by fixing one basis, with the two specialisations cited
in the status box below (the library also states \lean{latticeCocycle_zero_left}, \lean{_zero_right},
\lean{_add_left}, \lean{_add_right}, \lean{_sq} in the same way, each a one-line instance of the general
theorem):
\begin{leancode}
noncomputable def latticeBasis :
    Basis (Fin (Fintype.card (Module.Free.ChooseBasisIndex ℤ L.carrier))) ℤ L.carrier :=
  (Module.Free.chooseBasis ℤ L.carrier).reindex (Fintype.equivFin _)

noncomputable def latticeCocycle (α β : L.carrier) : ℂ := cocycle L (latticeBasis L) α β

theorem latticeCocycle_cocycleLaw (α β γ : L.carrier) :
    latticeCocycle L α β * latticeCocycle L (α + β) γ
      = latticeCocycle L β γ * latticeCocycle L α (β + γ) :=
  cocycle_cocycleLaw L (latticeBasis L) α β γ

theorem latticeCocycle_comm (hEven : L.Even) (α β : L.carrier) :
    latticeCocycle L α β * latticeCocycle L β α = (-1 : ℂ) ^ (L.form α β) :=
  cocycle_comm L (latticeBasis L) hEven α β
\end{leancode}
File \leanpath{VOA/Foundations/LatticeAlgebraTwisted.lean}.
\begin{leancode}
noncomputable def twMul (f g : L.carrier →₀ ℂ) : L.carrier →₀ ℂ := twMulL L f g

theorem twMul_e (α β : L.carrier) :
    twMul L (latticeAlgebra.e α) (latticeAlgebra.e β)
      = Finsupp.single (α + β) (latticeCocycle L α β) := by

theorem twMul_assoc (f g h : L.carrier →₀ ℂ) :
    twMul L (twMul L f g) h = twMul L f (twMul L g h) := by

theorem one_twMul (g : L.carrier →₀ ℂ) : twMul L (latticeAlgebra.e 0) g = g := by
\end{leancode}

\begin{dictionary}
\textbf{Context.} \lean{L : IntegralLattice} is an integral lattice (Definition~\ref{l02:def:lattice});
\lean{L.carrier} is its underlying free $\Z$-module and \lean{L.form} its $\Z$-valued symmetric form. \lean{n} is
the rank, and \lean{B : Basis (Fin n) ℤ L.carrier} is a $\Z$-basis $B_0,\dots,B_{n-1}$ of $L$ indexed by
\lean{Fin n} $=\{0,\dots,n-1\}$ with its usual order (the order enters through \lean{j < i}). The general
theorems hold for \emph{every} such basis. \textbf{Matching the theorem.} \lean{cocycle_add_left},
\lean{cocycle_add_right} are bimultiplicativity, \lean{cocycle_zero_left}, \lean{cocycle_zero_right}
normalisation, \lean{cocycle_cocycleLaw} the 2-cocycle law, \lean{cocycle_comm} the commutator property.
\lean{latticeBasis L} is a basis chosen once and for all (from Mathlib's \lean{Module.Free.chooseBasis},
reindexed by \lean{Fin}); \lean{latticeCocycle} is \lean{cocycle} at that basis, and its theorems are, by their
one-line proof terms, literally the general theorems instantiated at \lean{B := latticeBasis L}.
\lean{B.repr α i} is the $i$-th coordinate of $\alpha$ in the basis \lean{B}; \lean{cocExp} is the
``upper-triangular Gram'' exponent; \lean{(-1 : ℂ) ^ z} with $z\in\Z$ is an integer power, so $\varepsilon$ is
$\pm1$ (\lean{cocycle_sq}). \lean{latticeCocycle} fixes the basis \lean{latticeBasis L} (a chosen basis);
all properties hold for every basis. \lean{cocycle_comm} needs \lean{hEven}: the diagonal terms
$\langle B_i,B_i\rangle\alpha_i\beta_i$ are dropped from $\mathrm{coc}$, which is harmless exactly when they are even.
The FLM requirement is $\varepsilon(\alpha,\beta)/\varepsilon(\beta,\alpha) = (-1)^{\langle\alpha,\beta\rangle}$; since $\varepsilon=\pm1$,
quotient and product agree. \lean{twMul} is the bilinear extension of
$e^\alpha\star e^\beta = \varepsilon(\alpha,\beta)e^{\alpha+\beta}$ (\lean{twMul_e}), kept as an operation rather than a
ring instance so as not to clash with the untwisted product on the same carrier.
\end{dictionary}

\begin{physics}
$\varepsilon(\alpha,\beta)$ are the Klein factors (cocycle factors) of the compactified boson: without them,
vertex operators $e^{i\alpha X}$, $e^{i\beta X}$ with odd $\langle\alpha,\beta\rangle$ would anticommute instead of commute
(their operator product picks up $(-1)^{\langle\alpha,\beta\rangle}$ under exchange), and the lattice theory would not
be a consistent bosonic chiral algebra. The twisted group algebra $\C_\varepsilon[L]$ is the algebra of the
charge-shift (zero-mode) operators.
\end{physics}

\begin{proof}[Proof on paper]
$\mathrm{coc}$ is $\Z$-bilinear, so $\varepsilon$ is bimultiplicative, and both sides of the cocycle law equal
$(-1)^{\mathrm{coc}(\alpha,\beta)+\mathrm{coc}(\alpha,\gamma)+\mathrm{coc}(\beta,\gamma)}$. For the commutator,
$\mathrm{coc}(\alpha,\beta)+\mathrm{coc}(\beta,\alpha) = \sum_{i\ne j}\langle B_i,B_j\rangle\alpha_i\beta_j =
\langle\alpha,\beta\rangle-\sum_i\langle B_i,B_i\rangle\alpha_i\beta_i$ (using symmetry), and the last sum is even when $L$ is
even. Associativity of $\star$ on basis vectors is the cocycle law; normalisation gives the unit.
\end{proof}

\begin{status}
\lean{VOA.}\allowbreak\lean{IntegralLattice.}\allowbreak\lean{latticeCocycle_}\allowbreak\lean{cocycleLaw}, \lean{VOA.}\allowbreak\lean{IntegralLattice.}\allowbreak\lean{latticeCocycle_}\allowbreak\lean{comm},
\lean{VOA.}\allowbreak\lean{IntegralLattice.}\allowbreak\lean{twMul_}\allowbreak\lean{assoc}, \lean{VOA.}\allowbreak\lean{IntegralLattice.}\allowbreak\lean{one_}\allowbreak\lean{twMul}:
\lean{[propext, Classical.choice, Quot.sound]} (run for this chapter). Likewise \lean{VOA.}\allowbreak\lean{IntegralLattice.}\allowbreak\lean{cocycle_}\allowbreak\lean{add_}\allowbreak\lean{left},
\lean{cocycle_}\allowbreak\lean{add_}\allowbreak\lean{right}, \lean{cocycle_}\allowbreak\lean{zero_}\allowbreak\lean{left}, \lean{cocycle_}\allowbreak\lean{zero_}\allowbreak\lean{right},
\lean{cocycle_sq} (separate probe at an earlier commit, not re-run for this edition, importing \lean{VOA.Foundations.LatticeCocycle}). Unconditional.
\end{status}

\subsection{Finiteness of bounded vectors and the orthonormal frame}

\begin{theorem}\label{l02:thm:bf}
If $L$ is positive definite, then for every $N\in\Z$ the set $\{\alpha\in L : \langle\alpha,\alpha\rangle\le N\}$ is
finite. Moreover, writing $G$ for the real Gram matrix in the chosen basis and $C = \sqrt G$ (its positive
square root), the real coordinates $\tilde\alpha = C\cdot\mathrm{coord}(\alpha)$ satisfy
$\sum_a\tilde\alpha_a\tilde\beta_a = \langle\alpha,\beta\rangle$.
\end{theorem}

File \leanpath{VOA/Foundations/LatticeCocycle.lean}.
\begin{leancode}
def BoundedFinite (L : IntegralLattice) : Prop :=
  ∀ N : ℤ, {α : L.carrier | L.form α α ≤ N}.Finite
\end{leancode}
File \leanpath{VOA/Foundations/LatticeFinite.lean}.
\begin{leancode}
theorem boundedFinite (hPos : L.PosDef) : BoundedFinite L :=
  boundedFinite_of_coercive L hPos (coercive L hPos)
\end{leancode}
File \leanpath{VOA/Foundations/LatticeOrthoFrame.lean}.
\begin{leancode}
noncomputable def frame : Matrix (Fin (latRank L)) (Fin (latRank L)) ℝ := CFC.sqrt (gramR L)

noncomputable def orthCoord (α : L.carrier) : Fin (latRank L) → ℝ := frame L *ᵥ coordR L α

theorem orthCoord_inner (hPos : L.PosDef) (α β : L.carrier) :
    ∑ a, orthCoord L α a * orthCoord L β a = ((L.form α β : ℤ) : ℝ) := by
\end{leancode}

\begin{dictionary}
\lean{BoundedFinite} is the finiteness statement; \lean{boundedFinite} derives it from integer
positivity through real coercivity ($\exists m>0$, $m\|c\|^2\le Q(c)$), proving along the way exactly the
$\Z\Rightarrow\mathbb{R}$ positivity bridge discussed after Definition~\ref{l02:def:lattice}. \lean{CFC.sqrt} is
Mathlib's continuous-functional-calculus square root; \lean{*ᵥ} is matrix--vector multiplication;
\lean{coordR α} the real-cast coordinates. \lean{orthCoord_inner} is the second sentence of the
theorem, with the hypothesis \lean{hPos} needed for $C^2 = G$.
\end{dictionary}

\begin{physics}
Finiteness of bounded vectors makes each $L_0$-eigenspace of the lattice VOA finite-dimensional
(finitely many momenta below any energy), so that the partition function
$\theta_L(q)/\eta(q)^r$ is well defined. The orthonormal frame expresses the lattice bosons
$\alpha(z) = \sum_a\tilde\alpha_a\,\partial X^a(z)$ in terms of $r$ canonically normalised free bosons.
\end{physics}

\begin{status}
\lean{VOA.}\allowbreak\lean{IntegralLattice.}\allowbreak\lean{boundedFinite}, \lean{VOA.}\allowbreak\lean{IntegralLattice.}\allowbreak\lean{orthCoord_}\allowbreak\lean{inner}:
\lean{[propext, Classical.choice, Quot.sound]} (run for this chapter).
\end{status}

\subsection{Parity and the Koszul sign}

\begin{definition}\label{l02:def:parity}
For $p,q\in\Z/2$ let $(-1)^p$ be the sign character and $\epsilon(p,q) = (-1)^{pq}$ the Koszul sign. A
\emph{parity grading} of a vector space is a decomposition $V = V_{\bar0}\oplus V_{\bar1}$; its
\emph{fermion-number involution} $(-1)^F$ is $+1$ on $V_{\bar0}$ and $-1$ on $V_{\bar1}$.
\end{definition}

File \leanpath{VOA/Foundations/Parity.lean}.
\begin{leancode}
def parSign (p : ZMod 2) : ℂ := if p = 0 then 1 else -1

def koszul (p q : ZMod 2) : ℂ := parSign (p * q)

theorem koszul_add_left (p p' q : ZMod 2) :
    koszul (p + p') q = koszul p q * koszul p' q := by

structure ParityGrading (𝒱 : Type*) [AddCommGroup 𝒱] [Module ℂ 𝒱] where
  parity : ZMod 2 → Submodule ℂ 𝒱
  isInternal : DirectSum.IsInternal parity

noncomputable def parityInvol (P : ParityGrading 𝒱) : Module.End ℂ 𝒱 :=
  (2 : ℂ) • ((P.parity 0).projection (P.parity 1) P.isCompl01) - LinearMap.id

theorem koszulOp_homog (P : ParityGrading 𝒱) (p : ZMod 2) {q : ZMod 2} {u : 𝒱}
    (hu : u ∈ P.parity q) : P.koszulOp p u = koszul p q • u := by
\end{leancode}

\begin{dictionary}
\lean{ZMod 2} is $\Z/2$. \lean{parSign p} is $(-1)^p$ and \lean{koszul p q} is $(-1)^{pq}$.
\lean{ParityGrading} is an internal direct-sum decomposition indexed by $\Z/2$; it is deliberately
independent of the conformal-weight grading (no spin--statistics relation is imposed, so that the
$\beta\gamma$ and $bc$ systems fit). \lean{parityInvol} is $2\pi_0 - \mathrm{id}$ with $\pi_0$ the projection onto
$V_{\bar0}$ along $V_{\bar1}$, i.e.\ $(-1)^F$; \lean{koszulOp p} is $\mathrm{id}$ for $p = 0$ and $(-1)^F$ for
$p=1$, and \lean{koszulOp_homog} says it multiplies a vector of parity $q$ by $(-1)^{pq}$.
\end{dictionary}

\begin{physics}
The Koszul sign is Fermi statistics: exchanging two odd (fermionic) fields costs a sign. $(-1)^F$ is the
fermion-number operator, which also implements the Neveu--Schwarz/Ramond twist in super vertex
algebras (Chapter~\ref{ch:L17}).
\end{physics}

\begin{status}
\lean{VOA.koszul_}\allowbreak\lean{add_left}, \lean{VOA.}\allowbreak\lean{ParityGrading.}\allowbreak\lean{koszulOp_}\allowbreak\lean{homog}: \lean{[propext,
Classical.choice, Quot.sound]} (run for this chapter).
\end{status}

%======================================================================
\section{The verification record}\label{l02:sec:verification}
%======================================================================

The following was run by the author of this chapter at an earlier commit, after \lean{pbw} became a theorem
(not re-run for this edition), on three scratch files, with the command
\begin{plaincode}
lake env lean -DmaxSynthPendingDepth=3 PbwProbe.lean
\end{plaincode}
(never a bare \lean{lake env lean}, which can silently change elaboration; see Chapter~\ref{ch:15}). The
compiled library was already present and its \texttt{.olean} files were checked to be newer than the
corresponding sources; no \lean{lake build} was run for this chapter. The first file contained exactly
\begin{plaincode}
import VOA.Foundations.AffineLie
#print axioms VOA.pbw
#print axioms VOA.PBWTheorem.bottom_injective
#print axioms VOA.AffineLie.vacuumVector_ne_zero_of_induction
#print axioms VOA.AffineLie.vacuumVector_ne_zero
#print axioms VOA.AffineLie.vacuumSpace_nontrivial
#print axioms VOA.AffineLie.exists_vacuumMaxSubmodule
#print axioms VOA.AffineLie.vacuumSimple_isSimpleModule
\end{plaincode}
and its complete output was
\begin{plaincode}
'VOA.pbw' depends on axioms: [propext, Classical.choice, Quot.sound]
'VOA.PBWTheorem.bottom_injective' depends on axioms: [propext, Classical.choice, Quot.sound]
'VOA.AffineLie.vacuumVector_ne_zero_of_induction' depends on axioms: [propext, Classical.choice, Quot.sound]
'VOA.AffineLie.vacuumVector_ne_zero' depends on axioms: [propext, Classical.choice, Quot.sound]
'VOA.AffineLie.vacuumSpace_nontrivial' depends on axioms: [propext, Classical.choice, Quot.sound]
'VOA.AffineLie.exists_vacuumMaxSubmodule' depends on axioms: [propext, Classical.choice, Quot.sound]
'VOA.AffineLie.vacuumSimple_isSimpleModule' depends on axioms: [propext, Classical.choice, Quot.sound]
\end{plaincode}
The second file imported \lean{VOA.Foundations.PBWTheorem.VacuumBasis} and the files of the former
consumers (\lean{VOA.Modules.ZhuVermaAnArchDLM3PBWAxiom}, \lean{VOA.Modules.ZhuVermaAnArchLUBijectionPBW},
\lean{VOA.Modularity.ZhuStronglyRational}, \lean{VOA.Modularity.GNProp8PBW}, \lean{VOA.Modularity.AKSl2Payoff},
\lean{VOA.Examples.AffineSl2VOAInstance}); output:
\begin{plaincode}
'VOA.pbwBasis_holds' depends on axioms: [propext, Classical.choice, Quot.sound]
'VOA.VacuumBasis.ΦEquiv' depends on axioms: [propext, Classical.choice, Quot.sound]
'VOA.PBWBasis.vacuumVector_ne_zero' depends on axioms: [propext, Classical.choice, Quot.sound]
'VOA.pbwBasis_iff_of_subsingleton' depends on axioms: [propext, Classical.choice, Quot.sound]
'VOA.VHat.DLM3.voaPBWBottomInjective' depends on axioms: [propext, Classical.choice, Quot.sound]
'VOA.zhuBijective_of_voaPBW' depends on axioms: [propext, Classical.choice, Quot.sound]
'VOA.zhuBijective' depends on axioms: [propext, Classical.choice, Quot.sound]
'VOA.zhuModularity_stronglyRational_voaPBW' depends on axioms: [propext, Classical.choice, Quot.sound]
'VOA.zhuModularity_pinned_of_C2Cofinite_voaPBW' depends on axioms: [propext, Classical.choice, Quot.sound]
'VOA.zhuModularity_pinned_of_C2Cofinite' depends on axioms: [propext, Classical.choice, Quot.sound]
'VOA.AffineVOA.vacuumA_ne_zero' depends on axioms: [propext, Classical.choice, Quot.sound]
'VOA.AffineVOA.Inst.closed_vacuum_ne_zero' depends on axioms: [propext, Classical.choice, Quot.sound]
'VOA.QuasiLisseSl2.vacuum_simpleSl2QL_ne_zero' depends on axioms: [propext, Classical.choice, Quot.sound]
'VOA.QuasiLisseSl2.character_formalODE_simpleSl2QL' depends on axioms: [propext, Classical.choice, Quot.sound]
\end{plaincode}
The third file imported \lean{VOA.Foundations} and the files of the remaining headline results
(Virasoro nonvanishing and rigidity, the affine $\mathfrak{sl}_2$ and $\mathfrak{sl}_3$ files, and
\lean{VOA.Moonshine.TopPairs}) and re-probed every other declaration named in a status box of this chapter;
output (52 lines; line breaks inserted by the printer inside one list have been joined):
\begin{plaincode}
'VOA.Witt.gen_bracket' depends on axioms: [propext, Classical.choice, Quot.sound]
'VOA.Witt.cocycle_cond' depends on axioms: [propext, Classical.choice, Quot.sound]
'VOA.virasoro_bracket' depends on axioms: [propext, Classical.choice, Quot.sound]
'VOA.VirModule.vac_cyclic' depends on axioms: [propext, Classical.choice, Quot.sound]
'VOA.vacVec_ne_zero' depends on axioms: [propext, Classical.choice, Quot.sound]
'VOA.AffineLie.affine_bracket' depends on axioms: [propext, Classical.choice, Quot.sound]
'VOA.AffineLie.vacuum_acts_central' depends on axioms: [propext, Classical.choice, Quot.sound]
'VOA.AffineLie.vacuum_cyclic' depends on axioms: [propext, Classical.choice, Quot.sound]
'VOA.AffineLie.negSpan_eq_top' depends on axioms: [propext, Classical.choice, Quot.sound]
'VOA.AffineLie.ell0_isInternal' depends on axioms: [propext, Classical.choice, Quot.sound]
'VOA.AffineLie.eigenspace_zero_eq' depends on axioms: [propext, Classical.choice, Quot.sound]
'VOA.AffineLie.eigenspace_finiteDimensional' depends on axioms: [propext, Classical.choice, Quot.sound]
'VOA.RootDatum.affineEZero_bracket_affineFZero' depends on axioms: [propext, Classical.choice, Quot.sound]
'VOA.AffineLie.simpleHWOf_isSimpleModule' depends on axioms: [propext, Classical.choice, Quot.sound]
'VOA.IntegralLattice.latticeCocycle_comm' depends on axioms: [propext, Classical.choice, Quot.sound]
'VOA.IntegralLattice.boundedFinite' depends on axioms: [propext, Classical.choice, Quot.sound]
'VOA.IntegralLattice.twMul_assoc' depends on axioms: [propext, Classical.choice, Quot.sound]
'VOA.ParityGrading.koszulOp_homog' depends on axioms: [propext, Classical.choice, Quot.sound]
'VOA.AffineVOA.affineSl2VOA_closed' depends on axioms: [propext, Classical.choice, Quot.sound]
'VOA.AffineSl2WakimotoLiLeading.affineSl2PBWBasis' depends on axioms: [propext, Classical.choice, Quot.sound]
'VOA.AffineSl2BoundaryMaximalIdeal.affineSl2_vacuum_ne_zero' depends on axioms: [propext, Classical.choice, Quot.sound]
'VOA.Wak.vacuumA_ne_zero_free' depends on axioms: [propext, Classical.choice, Quot.sound]
'VOA.AffineSl2GeneralRigidity.negative_admissible_graded_rigidity' depends on axioms: [propext, Classical.choice, Quot.sound]
'VOA.VirModule.minimal_graded_rigidity' depends on axioms: [propext, Classical.choice, Quot.sound]
'VOA.Sl2NotSimple85.not_isSimple_85' depends on axioms: [propext, Classical.choice, Quot.sound]
'VOA.AffineSl3.simpleVOA' depends on axioms: [propext, Classical.choice, Quot.sound]
'VOA.AffineSl3.simple_regular_isTrivial' depends on axioms: [propext, Classical.choice, Quot.sound]
'VOA.Moonshine.moonshine' depends on axioms: [propext, Classical.choice, Quot.sound]
'VOA.AffineLie.C_lie' depends on axioms: [propext, Classical.choice, Quot.sound]
'VOA.AffineLie.ell0_lie_ofLoop' depends on axioms: [propext, Classical.choice, Quot.sound]
'VOA.AffineLie.vacuum_pos_annihilate' depends on axioms: [propext, Classical.choice, Quot.sound]
'VOA.AffineLie.eigenspace_eq_bot_of_not_energy' depends on axioms: [propext, Classical.choice, Quot.sound]
'VOA.AffineLie.simpleHighestVectorOf_ne_zero' depends on axioms: [propext, Classical.choice, Quot.sound]
'VOA.AffineLie.simpleHWOf_integrable_of_chevalley_locallyNilpotent' depends on axioms: [propext, Classical.choice, Quot.sound]
'VOA.RootDatum.cartanInteger' depends on axioms: [propext, Classical.choice, Quot.sound]
'VOA.RootDatum.exists_chevalleyTriple' depends on axioms: [propext, Classical.choice, Quot.sound]
'VOA.RootDatum.highestRoot_isNonZero' depends on axioms: [propext, Classical.choice, Quot.sound]
'VOA.C_central' depends on axioms: [propext, Classical.choice, Quot.sound]
'VOA.VirModule.vac_acts_central' depends on axioms: [propext, Classical.choice, Quot.sound]
'VOA.VirModule.vac_pos_annihilate' depends on axioms: [propext, Classical.choice, Quot.sound]
'VOA.vacuumSpace_nontrivial' depends on axioms: [propext, Classical.choice, Quot.sound]
'VOA.Witt.rep_injective' depends on axioms: [propext, Classical.choice, Quot.sound]
'VOA.PBWBasis.linearIndependent' depends on axioms: [propext, Classical.choice, Quot.sound]
'VOA.PBWTheorem.rep' depends on axioms: [propext, Classical.choice, Quot.sound]
'VOA.PBWTheorem.rep_p_zero' depends on axioms: [propext, Classical.choice, Quot.sound]
'VOA.IntegralLattice.latticeCocycle_cocycleLaw' depends on axioms: [propext, Classical.choice, Quot.sound]
'VOA.IntegralLattice.le_dual' depends on axioms: [propext, Classical.choice, Quot.sound]
'VOA.IntegralLattice.one_twMul' depends on axioms: [propext, Classical.choice, Quot.sound]
'VOA.IntegralLattice.orthCoord_inner' depends on axioms: [propext, Classical.choice, Quot.sound]
'VOA.stdLattice_selfDual' depends on axioms: [propext, Classical.choice, Quot.sound]
'VOA.koszul_add_left' depends on axioms: [propext, Classical.choice, Quot.sound]
'VOA.AffineSl3.vacuumA_ne_zero' depends on axioms: [propext, Classical.choice, Quot.sound]
\end{plaincode}
Together the three files cover every declaration named in a status box of this chapter except the five
cocycle lemmas \lean{cocycle_add_left}, \lean{cocycle_add_right}, \lean{cocycle_zero_left},
\lean{cocycle_zero_right}, \lean{cocycle_sq}, which were checked by a separate probe (Section~\ref{l02:sec:lattices});
the 73 displayed probes all report exactly \lean{[propext, Classical.choice, Quot.sound]}. A text search for uses of
\lean{VOA.pbw} in proof terms at that commit finds exactly the two direct consumers of
Section~\ref{l02:sec:consequences}. According to the project's build record, the aggregate
\texttt{lake build VOA} was last run green on 26 October 2026 (13868 jobs) with the axiom audit passing; this
was not re-run here.

\paragraph{Release audit 2026-10-27.} \texttt{Notes/\allowbreak Textbook/\allowbreak referee/\allowbreak RELEASE\_AUDIT\_FOUNDATIONS.md}
independently re-ran \texttt{\#print axioms VOA.pbw} (the three standard axioms) and confirmed that
\lean{VOA.pbw} is a theorem at \leanpath{VOA/\allowbreak Foundations/\allowbreak AffineLie.lean} and that the repository declares no
axiom. Its findings bearing on this chapter: \lean{AffineVOA.WakModule} is the universal vacuum module $V^k(\mathfrak{sl}_2)$
of this chapter, not a Wakimoto module; the derivation $\ell_0$ is killed by the vacuum ideal, which gives the same vacuum
module; the Virasoro VOA's carrier is \lean{VirUniv c}, the universal $V(c,0)$; and the zero space counts as a VOA, so
nonvanishing results such as \lean{vacuumVector_ne_zero} carry real content.

\paragraph{Historical record (first edition).} The first edition's probe, on a
scratch file importing \lean{VOA}, gave the following (abridged to one line per declaration). It is kept
because it documents what depended on the axiom while there was one.
\begin{plaincode}
'VOA.pbw' depends on axioms: [propext, Classical.choice, Quot.sound, pbw]
'VOA.AffineLie.affine_bracket' depends on axioms: [propext, Classical.choice, Quot.sound]
'VOA.AffineLie.vacuum_acts_central' depends on axioms: [propext, Classical.choice, Quot.sound]
'VOA.AffineLie.vacuum_cyclic' depends on axioms: [propext, Classical.choice, Quot.sound]
'VOA.AffineLie.vacuumVector_ne_zero' depends on axioms: [propext, Classical.choice, Quot.sound, pbw]
'VOA.AffineLie.vacuumSpace_nontrivial' depends on axioms: [propext, Classical.choice, Quot.sound, pbw]
'VOA.AffineLie.exists_vacuumMaxSubmodule' depends on axioms: [propext, Classical.choice, Quot.sound, pbw]
'VOA.AffineLie.vacuumSimple_isSimpleModule' depends on axioms: [propext, Classical.choice, Quot.sound, pbw]
'VOA.AffineLie.negSpan_eq_top' depends on axioms: [propext, Classical.choice, Quot.sound]
'VOA.AffineLie.ell0_isInternal' depends on axioms: [propext, Classical.choice, Quot.sound]
'VOA.AffineLie.eigenspace_zero_eq' depends on axioms: [propext, Classical.choice, Quot.sound]
'VOA.AffineLie.eigenspace_finiteDimensional' depends on axioms: [propext, Classical.choice, Quot.sound]
'VOA.PBWBasis.vacuumVector_ne_zero' depends on axioms: [propext, Classical.choice, Quot.sound]
'VOA.pbwBasis_iff_of_subsingleton' depends on axioms: [propext, Classical.choice, Quot.sound]
'VOA.Witt.gen_bracket' depends on axioms: [propext, Classical.choice, Quot.sound]
'VOA.Witt.cocycle_cond' depends on axioms: [propext, Classical.choice, Quot.sound]
'VOA.virasoro_bracket' depends on axioms: [propext, Classical.choice, Quot.sound]
'VOA.VirModule.vac_cyclic' depends on axioms: [propext, Classical.choice, Quot.sound]
'VOA.vacVec_ne_zero' depends on axioms: [propext, Classical.choice, Quot.sound]
'VOA.RootDatum.affineEZero_bracket_affineFZero' depends on axioms: [propext, Classical.choice, Quot.sound]
'VOA.AffineLie.simpleHWOf_isSimpleModule' depends on axioms: [propext, Classical.choice, Quot.sound]
'VOA.IntegralLattice.latticeCocycle_comm' depends on axioms: [propext, Classical.choice, Quot.sound]
'VOA.IntegralLattice.boundedFinite' depends on axioms: [propext, Classical.choice, Quot.sound]
'VOA.IntegralLattice.twMul_assoc' depends on axioms: [propext, Classical.choice, Quot.sound]
'VOA.ParityGrading.koszulOp_homog' depends on axioms: [propext, Classical.choice, Quot.sound]
'VOA.VHat.DLM3.voaPBWBottomInjective' depends on axioms: [propext, Classical.choice, Quot.sound, pbw]
'VOA.zhuBijective_of_voaPBW' depends on axioms: [propext, Classical.choice, Quot.sound, pbw]
'VOA.zhuBijective' depends on axioms: [propext, Classical.choice, Quot.sound]
'VOA.zhuModularity_stronglyRational_voaPBW' depends on axioms: [propext, Classical.choice, Quot.sound, pbw]
'VOA.AffineVOA.Inst.closed_vacuum_ne_zero' depends on axioms: [propext, Classical.choice, Quot.sound, pbw]
'VOA.AffineVOA.affineSl2VOA_closed' depends on axioms: [propext, Classical.choice, Quot.sound]
'VOA.AffineSl2WakimotoLiLeading.affineSl2PBWBasis' depends on axioms: [propext, Classical.choice, Quot.sound]
'VOA.AffineSl2BoundaryMaximalIdeal.affineSl2_vacuum_ne_zero' depends on axioms: [propext, Classical.choice, Quot.sound]
'VOA.AffineSl2GeneralRigidity.negative_admissible_graded_rigidity' depends on axioms: [propext, Classical.choice, Quot.sound]
'VOA.VirModule.minimal_graded_rigidity' depends on axioms: [propext, Classical.choice, Quot.sound]
'VOA.Sl2NotSimple85.not_isSimple_85' depends on axioms: [propext, Classical.choice, Quot.sound]
'VOA.Moonshine.moonshine' depends on axioms: [propext, Classical.choice, Quot.sound]
\end{plaincode}
A second probe then gave \lean{[propext, Classical.choice, Quot.sound]} for
\lean{VOA.}\allowbreak\lean{Wak.}\allowbreak\lean{vacuumA_}\allowbreak\lean{ne_}\allowbreak\lean{zero_}\allowbreak\lean{free}, \lean{VOA.}\allowbreak\lean{zhuModularity_}\allowbreak\lean{pinned_}\allowbreak\lean{of_}\allowbreak\lean{C2Cofinite},
\lean{VOA.}\allowbreak\lean{AffineSl3.}\allowbreak\lean{simpleVOA}, \lean{VOA.}\allowbreak\lean{AffineSl3.}\allowbreak\lean{simple_}\allowbreak\lean{regular_}\allowbreak\lean{isTrivial},
\lean{VOA.}\allowbreak\lean{AffineSl3.}\allowbreak\lean{vacuumA_}\allowbreak\lean{ne_}\allowbreak\lean{zero}, and \lean{[propext, Classical.choice, Quot.sound, VOA.pbw]} for
\lean{VOA.}\allowbreak\lean{AffineVOA.}\allowbreak\lean{vacuumA_}\allowbreak\lean{ne_}\allowbreak\lean{zero}, \lean{VOA.}\allowbreak\lean{QuasiLisseSl2.}\allowbreak\lean{vacuum_}\allowbreak\lean{simpleSl2QL_}\allowbreak\lean{ne_}\allowbreak\lean{zero},
\lean{VOA.}\allowbreak\lean{QuasiLisseSl2.}\allowbreak\lean{character_}\allowbreak\lean{formalODE_}\allowbreak\lean{simpleSl2QL} and
\lean{VOA.}\allowbreak\lean{zhuModularity_}\allowbreak\lean{pinned_}\allowbreak\lean{of_}\allowbreak\lean{C2Cofinite_}\allowbreak\lean{voaPBW} (58 declarations were probed in all).

%======================================================================
\section{Referee's checklist}\label{l02:sec:checklist}
%======================================================================

The table lists the chapter's headline declarations with their intended meaning, their hypotheses
and machine status, and the place where the fidelity check is carried out. Every declaration listed
reports exactly \lean{[propext, Classical.choice, Quot.sound]} (Section~\ref{l02:sec:verification}); the
library has no project axiom. Entries in bold mark a condition not discharged in the library, or a
tautology.

\begingroup\small
\begin{longtable}{p{0.25\textwidth}p{0.29\textwidth}p{0.18\textwidth}p{0.14\textwidth}}
\toprule
Lean declaration & Intended meaning & Hypotheses / status & Fidelity check\\
\midrule\endhead
\lean{VOA.pbw} (\lean{PBWPrinciple}) & (A) affine vacuum $\ne0$ for all $\ltwog,\kappa,k$; (B) $U\to U(\ltwog)\otimes_{U(\ltwop)}U$ injective for all $\ltwog,\ltwop,\rho$ & theorem, unconditional (formerly a project axiom) & \S\ref{l02:sec:axiom}\\
\lean{PBWTheorem.}\allowbreak\lean{rep}, \lean{rep_p_zero} & relative PBW representation on $S(\mathfrak c)\otimes U$ & unconditional, any field & Thm.~\ref{l02:thm:relrep}\\
\lean{PBWTheorem.}\allowbreak\lean{bottom_}\allowbreak\lean{injective} & field (B), any field & unconditional & Thm.~\ref{l02:thm:bottominj}\\
\lean{AffineLie.}\allowbreak\lean{vacuumVector_}\allowbreak\lean{ne_}\allowbreak\lean{zero_}\allowbreak\lean{of_}\allowbreak\lean{induction} & field (A) from (B) & unconditional & Thm.~\ref{l02:thm:vacind}\\
\lean{PBWInduction.}\allowbreak\lean{Induced}, \lean{bottom} & induced module, $u\mapsto[1\otimes u]$ & definitions & Prop.~\ref{l02:prop:induced}\\
\lean{AffineLie.}\allowbreak\lean{vacuumVector_}\allowbreak\lean{ne_}\allowbreak\lean{zero} & $\ltwovac\ne0$ in $V^k(\ltwog)$ & none; via theorem \lean{VOA.pbw} & \S\ref{l02:sec:consequences}\\
\lean{AffineLie.}\allowbreak\lean{vacuumSimple_}\allowbreak\lean{isSimpleModule} & $L(0)$ simple & none; chosen coatom & \S\ref{l02:sec:consequences}\\
\lean{VHat.}\allowbreak\lean{DLM3.}\allowbreak\lean{voaPBWBottom}\allowbreak\lean{Injective} & bottom of DLM induced module embeds & none; via theorem \lean{VOA.pbw} & \S\ref{l02:sec:consequences}\\
\lean{Witt.gen_bracket} & $[\ell_m,\ell_n]=(m-n)\ell_{m+n}$ & unconditional & \S\ref{l02:sec:witt}\\
\lean{virasoro_bracket} & Virasoro relation, $c/12\,(m^3-m)$ & unconditional & \S\ref{l02:sec:witt}\\
\lean{vacVec_ne_zero} & Virasoro vacuum $\ne0$ & unconditional, pbw-free & \S\ref{l02:sec:virmod}\\
\lean{AffineLie.}\allowbreak\lean{affine_}\allowbreak\lean{bracket} & affine relation with $m\kappa\delta C$ & \lean{hκ} (invariance) & \S\ref{l02:sec:affine}\\
\lean{AffineLie.}\allowbreak\lean{vacuum_}\allowbreak\lean{cyclic} etc. & vacuum conditions, cyclicity & unconditional & \S\ref{l02:sec:vacuum}\\
\lean{AffineLie.}\allowbreak\lean{negSpan_}\allowbreak\lean{eq_}\allowbreak\lean{top} & creation monomials span & unconditional & \S\ref{l02:sec:grading}\\
\lean{AffineLie.}\allowbreak\lean{ell0_}\allowbreak\lean{isInternal}, \lean{eigenspace_}\allowbreak\lean{zero_eq}, \lean{eigenspace_}\allowbreak\lean{finiteDimensional} & energy grading, $V_0=\C\ltwovac$ (spanned), finite pieces & last needs $\dim\ltwog<\infty$ & \S\ref{l02:sec:grading}\\
\lean{PBWBasis} & ordered monomials form a basis & definition (a proposition); proved by \lean{pbwBasis_holds} & \S\ref{l02:sec:pbwbasis}\\
\lean{pbwBasis_holds} & PBW basis of $V^k(\ltwog)$, all $\ltwog,\kappa,b,k$ & unconditional & Thm.~\ref{l02:thm:pbwbasisholds}\\
\lean{PBWBasis.}\allowbreak\lean{vacuumVector_}\allowbreak\lean{ne_}\allowbreak\lean{zero} & PBWBasis $\Rightarrow\ltwovac\ne0$ & \lean{h : PBWBasis …} & \S\ref{l02:sec:pbwbasis}\\
\lean{pbwBasis_}\allowbreak\lean{iff_}\allowbreak\lean{of_}\allowbreak\lean{subsingleton} & shape is satisfiable & $\ltwog=0$ & \S\ref{l02:sec:pbwbasis}\\
\lean{affineSl2}\allowbreak\lean{PBWBasis} & PBW basis for $\widehat{\mathfrak{sl}}_2$ & $k\ne-2$; Wakimoto proof, no \lean{VOA.pbw} & \S\ref{l02:sec:pbwbasis}\\
\lean{RootDatum.}\allowbreak\lean{exists_}\allowbreak\lean{chevalleyTriple}, \lean{cartanInteger} & root datum & \lean{IsKilling}, Cartan & \S\ref{l02:sec:roots}\\
\lean{AffineLie.}\allowbreak\lean{simpleHWOf_}\allowbreak\lean{isSimpleModule} & $L(\lambda)$ simple & \textbf{conditional} on $v_\lambda\ne0$ (open in library) & \S\ref{l02:sec:roots}\\
\lean{simpleHWOf_}\allowbreak\lean{integrable_}\allowbreak\lean{of_}\allowbreak\lean{chevalley_}\allowbreak\lean{locallyNilpotent} & integrability & \textbf{tautology}; Kac--Garland open & \S\ref{l02:sec:roots}\\
\lean{stdLattice_}\allowbreak\lean{selfDual} & $\Z^n$ unimodular & unconditional & \S\ref{l02:sec:lattices}\\
\lean{latticeCocycle_}\allowbreak\lean{comm} & $\varepsilon(\alpha,\beta)\varepsilon(\beta,\alpha)=(-1)^{\langle\alpha,\beta\rangle}$ & \lean{L.Even} & \S\ref{l02:sec:lattices}\\
\lean{twMul_assoc} & $\C_\varepsilon[L]$ associative & unconditional & \S\ref{l02:sec:lattices}\\
\lean{boundedFinite} & finitely many short vectors & \lean{L.PosDef} & \S\ref{l02:sec:lattices}\\
\lean{orthCoord_inner} & orthonormal frame realises form & \lean{L.PosDef} & \S\ref{l02:sec:lattices}\\
\lean{koszulOp_homog} & $(-1)^{pq}$ on homogeneous vectors & unconditional & \S\ref{l02:sec:lattices}\\
\bottomrule
\end{longtable}
\endgroup

Points the referee should take away:
\begin{enumerate}
\item The library has \emph{no} project axiom. \lean{VOA.pbw}, formerly the project's sole axiom,
  is now a theorem proved from Mathlib alone by relative PBW
  (Section~\ref{l02:sec:pbwproof}); \texttt{\#print axioms VOA.pbw} reports the three standard axioms.
\item \lean{PBWPrinciple} is the intended statement: Section~\ref{l02:sec:axiom} unfolds each definition it
  mentions, identifies its two fields with classical consequences of PBW, and lists the slips that would
  have made it false or empty.
\item \lean{VOA.pbw} is used directly in exactly two places. The results that used to inherit the axiom
  (the \lean{_voaPBW} wrappers and the legacy affine $\mathfrak{sl}_2$ nonvanishing route) are now
  \texttt{\#print axioms}-clean; those that carry other hypotheses (e.g.\ \lean{hC2}, \lean{hSR},
  \lean{hFin}, \lean{hZhu}) remain conditional on them as stated. Of these, \lean{hFin}, \lean{hZhu},
  \lean{hss} and \lean{[Finite (Φ W)]} are now discharged by library theorems, and Zhu's theorem holds for every
  $C_2$-cofinite \lean{AdmissiblyRational} VOA (\lean{zhuModularity_}\allowbreak\lean{canonical_of_}\allowbreak\lean{C2Cofinite_}\allowbreak\lean{admissiblyRational}); only the relation
  of \lean{AdmissiblyRational} to Dong--Li--Mason rationality is open. The axiom-free twins with \lean{Pr : BottomProjection …},
  \lean{PBWFreeness} or \lean{h : PBWBasis …} hypotheses are historical; those hypotheses are now always
  dischargeable.
\item \lean{PBWBasis}, introduced as a named hypothesis, is a theorem for every $\ltwog$, $\kappa$, ordered basis
  and level (\lean{pbwBasis_holds}).
\item Documentation lags found (a sample, not an exhaustive list): the \leanpath{PBWBasis.lean} module
  docstring, which still calls \lean{pbw} an axiom and \lean{PBWBasis} an unproved hypothesis; the
  \leanpath{AffineSl2VOAInstance.lean} remark that no pbw-free proof of nonvanishing is possible; and the
  closing comment of \leanpath{ZhuVermaAnArchDLM3PBWAxiom.lean}. Comments calling \lean{VOA.pbw} an axiom are
  out of date throughout the library (Notes/Textbook/referee/pbw\_sweep.md lists the book-side
  occurrences). None affects a Lean statement.
\item Two results of \leanpath{RootDatum.lean} are conditional on an undischarged (but true) hypothesis,
  and one is a tautology; none of them is used by a headline theorem.
\end{enumerate}

%======================================================================
\section{Notes and references}\label{l02:sec:notes}
%======================================================================

\textbf{PBW.} The theorem goes back to H.~Poincar\'e (1900), G.~Birkhoff (1937) and E.~Witt (1937). Our
proof sketch follows J.~E.~Humphreys, \emph{Introduction to Lie Algebras and Representation Theory},
Graduate Texts in Mathematics 9, Springer, 1972, \S17.3--17.4; the general (any dimension, any field)
statement and the freeness of $U(\ltwog)$ over $U(\ltwop)$ are in J.~Dixmier, \emph{Enveloping Algebras},
North-Holland, 1977 (reprinted by the AMS, Graduate Studies in Mathematics 11, 1996), \S2.1
(Theorem 2.1.11, the reference cited in the library's docstrings) and \S2.2. The diamond-lemma proof is in G.~M.~Bergman, ``The diamond lemma for ring theory'',
\emph{Adv.\ Math.} 29 (1978) 178--218.

\textbf{Affine and Virasoro algebras.} V.~G.~Kac, \emph{Infinite-Dimensional Lie Algebras}, 3rd ed.,
Cambridge University Press, 1990: Chapter 7 for the realisation of affine algebras as central extensions
of loop algebras, Chapter 9 for highest-weight and Verma modules, Chapter 10 for integrable
modules, Chapters 1--4 for Cartan matrices and root systems. V.~G.~Kac and A.~K.~Raina, \emph{Bombay Lectures on Highest
Weight Representations of Infinite Dimensional Lie Algebras}, World Scientific, 1987, for the Virasoro
algebra and its cocycle. T.~Gannon, \emph{Moonshine beyond the Monster}, Cambridge University Press, 2006,
\S\S3.1--3.2, is the source the library's docstrings follow. The vacuum modules as vertex algebras are
treated in E.~Frenkel and D.~Ben-Zvi, \emph{Vertex Algebras and Algebraic Curves}, 2nd ed., AMS
Mathematical Surveys and Monographs 88, 2004, Chapter 2, and in V.~G.~Kac, \emph{Vertex Algebras for
Beginners}, 2nd ed., AMS University Lecture Series 10, 1998. The Wakimoto realisation
is due to M.~Wakimoto, ``Fock representations of the affine Lie algebra $A_1^{(1)}$'', \emph{Comm.\ Math.\
Phys.} 104 (1986) 605--609. The integrability theorem is H.~Garland, ``The arithmetic theory of loop
algebras'', \emph{J.\ Algebra} 53 (1978) 480--551, and Kac, \emph{op.\ cit.}, Chapter 10.

\textbf{Physics.} The WZW model and its current algebra: E.~Witten, ``Nonabelian bosonization in two
dimensions'', \emph{Comm.\ Math.\ Phys.} 92 (1984) 455--472; V.~G.~Knizhnik and A.~B.~Zamolodchikov,
\emph{Nucl.\ Phys.} B247 (1984) 83--103; and the textbook P.~Di~Francesco, P.~Mathieu and
D.~S\'en\'echal, \emph{Conformal Field Theory}, Springer, 1997, Part A (the Virasoro algebra and its
representations) and Part C (affine Lie algebras and WZW models).

\textbf{Lattices and cocycles.} I.~Frenkel, J.~Lepowsky and A.~Meurman, \emph{Vertex Operator Algebras
and the Monster}, Academic Press, 1988, Chapters 5--8; J.~H.~Conway and N.~J.~A.~Sloane, \emph{Sphere
Packings, Lattices and Groups}, 3rd ed., Springer, 1999. Super signs: Kac, \emph{Vertex Algebras for
Beginners}, the sections on the super case.

\textbf{Lean.} The consistency of Lean's type theory relative to set theory with inaccessible cardinals is
M.~Carneiro, \emph{The Type Theory of Lean}, Master's thesis, Carnegie Mellon University, 2019. Mathlib's
universal enveloping algebra is in \leanpath{Mathlib/\allowbreak Algebra/\allowbreak Lie/\allowbreak UniversalEnveloping.lean}
(pinned version v4.30.0).

\endgroup
